% !TeX program = pdflatex
% Zbior zadan uzupelniajacych rozdzialy 1--35.
% Zadania celowo nie powtarzaja dowodow i konstrukcji wykonanych w skrypcie.
\documentclass[11pt,a4paper,oneside,openany]{book}
\usepackage[T1]{fontenc}
\usepackage[utf8]{inputenc}
\usepackage[polish]{babel}
\usepackage{lmodern}
\usepackage{amsmath,amssymb,amsthm,mathtools}
\usepackage{microtype}
\usepackage{geometry}
\usepackage{enumitem}
\usepackage[unicode,colorlinks=true,linkcolor=blue!40!black]{hyperref}
\usepackage{xcolor}
\geometry{margin=2.55cm}
\pagestyle{plain}
\setlength{\parindent}{0pt}
\setlength{\emergencystretch}{2em}
\setlist[enumerate]{leftmargin=*,itemsep=.65em,topsep=.7em}
\newcommand{\R}{\mathbb R}
\newcommand{\C}{\mathbb C}
\newcommand{\Z}{\mathbb Z}
\newcommand{\Q}{\mathbb Q}
\newcommand{\N}{\mathbb N}
\newcommand{\T}{\mathbb T}
\newcommand{\dd}{\mathrm d}
\newcommand{\Hom}{\operatorname{Hom}}
\newcommand{\rank}{\operatorname{rank}}
\newcommand{\im}{\operatorname{im}}
\newcommand{\kerop}{\operatorname{ker}}
\newcommand{\Cl}{\operatorname{Cl}}
\newcommand{\Spin}{\operatorname{Spin}}
\newcommand{\Gr}{\operatorname{Gr}}
\newcommand{\Ric}{\operatorname{Ric}}
\newcommand{\Scal}{\operatorname{Scal}}
\newcommand{\Wh}{\operatorname{Wh}}
\newcommand{\id}{\operatorname{id}}
\newcommand{\wskazowka}[1]{%
  \par\smallskip\noindent{\footnotesize\color{black!75}\emph{Wskazówka.}\enspace #1}\par
}
\begin{document}
\frontmatter
\begin{center}
{\LARGE\bfseries Zadania do rozdziałów 1--35}\par\vspace{.8em}
{\large Geometria i topologia różniczkowa}\par\vspace{.8em}
Zadania rachunkowe, konstrukcyjne i dowodowe. Wskazówki do trudniejszych zadań, bez rozwiązań.
\end{center}
\bigskip
\noindent\textit{Uwagi.} Zestaw uzupełnia skrypt: nie poleca ponownie
udowadniać wyników ani odtwarzać przykładów, dla których
pełny dowód lub rachunek podano już w rozdziałach.
Liczba zadań w rozdziale wynika z zakresu materiału, a nie ze sztywnego limitu.
Rozdział 21 stanowi pogłębienie zagadnień torsji Whiteheada
z rozdziału 20. Rozdziały 31--35 dotyczą szczególnych zjawisk
w topologii czterowymiarowej; twierdzenia o istnieniu przyjmujemy
tam zgodnie z zakresem skryptu.
Wolno stosować bez dowodu twierdzenia ze skryptu, chyba że treść
konkretnego zadania mówi inaczej. Każde zadanie może mieć kilka podpunktów. Krótkie wskazówki zamieszczono przy wybranych trudniejszych zadaniach; wskazują metodę, nie podają rozwiązań.
\tableofcontents
\mainmatter

\chapter{Rozmaitości różniczkowe}
\begin{enumerate}
\item \textbf{Rozmaitość macierzy specjalnych.} Niech $SL(2,\R)=\{A\in M_2(\R):\det A=1\}$. Traktując wyznacznik jako funkcję na $\R^4$, sprawdź regularność wartości $1$ i oblicz $T_A SL(2,\R)$. Wyraź odpowiedź bez współrzędnych za pomocą $\operatorname{tr}(A^{-1}H)$.
\item \textbf{Macierze ustalonego rzędu.} Pokaż, że zbiór macierzy $2\times3$ o rzędzie dokładnie $1$ jest gładką podrozmaitością $M_{2\times3}(\R)$. Oblicz jej wymiar i przestrzeń styczną w punkcie $\begin{psmallmatrix}1&0&0\\0&0&0\end{psmallmatrix}$.
\wskazowka{Wybierz niezerowy minor $1\times1$ i rozwiąż równania zerowania minorów $2\times2$ w jego otoczeniu. Przestrzeń styczną otrzymasz przez zróżniczkowanie tych równań.}
\item \textbf{Funkcja odwrotności.} Dla $F:GL(n,\R)\to GL(n,\R)$, $F(A)=A^{-1}$, oblicz $\dd F_A(H)$. Następnie znajdź różniczkę $A\mapsto\log|\det A|$ i porównaj oba rachunki.
\item \textbf{Rozmaitość ortonormalnych par.} Rozważ $V_2(\R^3)=\{(u,v)\in\R^3\times\R^3:|u|=|v|=1,\ u\cdot v=0\}$. Zbadaj regularność wartości $(1,1,0)$ odwzorowania $F(u,v)=(|u|^2,|v|^2,u\cdot v)$, wyznacz przestrzeń styczną do $V_2(\R^3)$ i sprawdź, że $(u,v)\mapsto(u,v,u\times v)$ definiuje dyfeomorfizm $V_2(\R^3)\to SO(3)$.
\wskazowka{Najpierw oblicz $dF_{(u,v)}(h,k)$, a potem sprawdź niezależność trzech funkcjonałów. Trzecią kolumnę macierzy ortogonalnej wyznacza orientacja.}
\item \textbf{Stożek i wartość osobliwa.} Rozważ $X=\{(x,y,z)\in\R^3:z^2=x^2+y^2\}$. Wykaż, że $X\setminus\{0\}$ jest gładką powierzchnią i opisz jej przestrzenie styczne. Wyjaśnij, dlaczego przy wierzchołku nie może być ono $C^1$-zanurzoną powierzchnią w $\R^3$.
\wskazowka{Poza początkiem zastosuj twierdzenie o wartości regularnej do $z^2-x^2-y^2$. W wierzchołku porównaj graniczne płaszczyzny styczne wzdłuż różnych tworzących stożka.}
\item \textbf{Dwa równania i transwersalność.} Niech $F:\R^4\to\R^2$, $F(x_1,x_2,x_3,x_4)=(x_1^2+x_2^2,x_3+x_4)$. Opisz włókna $F^{-1}(a,b)$ dla $a>0$, ich wymiar i przestrzenie styczne. Co zmienia się przy $a=0$?
\item \textbf{Różniczka odwzorowania zadanego na ilorazie.} Na torusie $\T^2=\R^2/\Z^2$ określ $f([x,y])=[2x+y,x+3y]$. Sprawdź poprawność definicji, oblicz $\dd f$ i rozstrzygnij, czy $f$ jest lokalnym dyfeomorfizmem, dyfeomorfizmem czy nakryciem wielokrotnym.
\wskazowka{Zbadaj homomorfizm $\mathbb Z^2\to\mathbb Z^2$ dany przez macierz odwzorowania. Liczbę punktów włókna powiąż z indeksem podgrupy obrazu.}
\item \textbf{Wiązka nad okręgiem.} Sklej $[0,1]\times\R^2$ relacją $(0,v)\sim(1,Av)$, gdzie $A=\operatorname{diag}(-1,2)$. Opisz funkcje przejścia, wyznacz wiązkę wyznacznikową i rozstrzygnij, czy otrzymana wiązka rzędu $2$ jest trywialna.
\wskazowka{Wiązka wyznacznikowa ma funkcję sklejającą $\det A$. Przy klasyfikacji rzeczywistych wiązek nad $S^1$ istotna jest składowa $GL(2,\R)$ zawierająca $A$.}
\item \textbf{Wiązka indukowana ze wstęgi Möbiusa.} Dla $p_k:S^1\to S^1$, $p_k(z)=z^k$, zbadaj wiązkę $p_k^*L$, gdzie $L$ jest tautologiczną wiązką rzeczywistą nad $\mathbb{RP}^1\cong S^1$. Dla jakich $k$ ma ona sekcję nigdzie niezerową?
\wskazowka{Rozważ monodromię sklejania wstęgi Möbiusa oraz jej zachowanie po $k$-krotnym obejściu okręgu. Alternatywnie oblicz cofnięcie klasy $w_1$.}
\item \textbf{Wektor styczny a derywacja.} Dla krzywej $\gamma(t)=(\cos(2t),\sin(2t),t^2)$ w $\R^3$ i funkcji $f(x,y,z)=x^2y+e^z$, oblicz derywację odpowiadającą $\dot\gamma(0)$. Porównaj wynik z różniczką restrykcji $f$ do obrazu krzywej.
\item \textbf{Pochodna na przestrzeni rzutowej.} Dla $\mathbb{RP}^2$ rozważ na wspólnym obszarze mapy $[1:u:v]$ i $[s:1:t]$. Oblicz jakobian przejścia i przelicz składowe wektora $(a,b)$ z pierwszej mapy na drugą w wybranym punkcie $u\ne0$.
\item \textbf{Przestrzenie styczne produktu poziomic.} Niech $M=\{x\in\R^3:\|x\|=2\}$, $N=\{(u,v)\in\R^2:u^2-v^2=1\}$. Opisz jawnie $T_{(p,q)}(M\times N)$ jako podprzestrzeń $\R^5$ i oblicz jej wymiar.
\item \textbf{Krzywa parametryczna a osadzenie.} Zbadaj $\gamma:\R\to\R^3$, $\gamma(t)=(t^2,t^3,t^4)$. Znajdź punkty, w których ma pełny rząd, i rozstrzygnij, czy jej obraz z naturalną topologią podprzestrzeni jest gładką podrozmaitością jednowymiarową $\R^3$.
\wskazowka{Sprawdź pochodną w $t=0$, ale nie poprzestawaj na nieimmersyjności parametryzacji: porównaj stożek styczny obrazu z prostą styczną hipotetycznej podrozmaitości.}
\item \textbf{Morfizm wiązek.} Dla gładkiego $F:M\to N$ rozważ $\dd F:TM\to TN$. W przykładzie $F(x,y)=(x^2-y^2,2xy)$ wyznacz miejsca, w których morfizm wiązek jest izomorfizmem na włóknach, i scharakteryzuj jego jądro na pozostałych włóknach.
\item \textbf{Zespolona krzywa algebraiczna.} Dla $C=\{(z,w)\in\C^2:w^2=z^3-z\}$ sprawdź, że $C$ jest jednowymiarową rozmaitością zespoloną. Wyznacz jej przestrzeń styczną w dowolnym punkcie i wskaż, na których punktach projekcja $(z,w)\mapsto z$ przestaje być lokalnie biholomorficzna.
\wskazowka{Użyj holomorficznego twierdzenia o funkcji uwikłanej dla $F(z,w)=w^2-z^3+z$. Dla projekcji badaj, kiedy $dz$ znika na $\ker dF$.}
\item \textbf{Działanie skończonej grupy i iloraz.} Na $\R^2\setminus\{0\}$ rozważ działanie $(x,y)\mapsto(-x,-y)$. Wykaż, że iloraz jest gładką rozmaitością, oblicz lokalny model i porównaj sytuację z działaniem na całej $\R^2$. Wskaż dokładnie, które założenie twierdzenia o gładkim ilorazie przestaje zachodzić w punkcie stałym.
\wskazowka{Po usunięciu zera działanie jest wolne i właściwe. W początku porównaj grupę stabilizującą oraz lokalny model ilorazowy z twierdzeniem o wolnym działaniu.}
\end{enumerate}


\chapter{Pola wektorowe, przepływy i pochodna Liego}
\begin{enumerate}
\item \textbf{Nawias pól trójkątnych.} Dla $X=\partial_x+2x\partial_y$, $Y=\partial_y+y\partial_z$ na $\R^3$ oblicz $[X,Y]$, $[X,[X,Y]]$ i $[Y,[X,Y]]$. Rozstrzygnij, czy rozpięta przez $X,Y$ dystrybucja jest involutywna.
\item \textbf{Jednorodność i pole Eulera.} Niech $E=\sum_{j=1}^n x_j\partial_{x_j}$. Dla pola $Y=x^2\partial_y+xy\partial_z$ na $\R^3$ oblicz $[E,Y]$ i sprawdź wynik za pomocą jawnego przepływu $E$.
\item \textbf{Przepływy dwóch ścinań.} Wyznacz przepływy pól $X=y\partial_x$ i $Y=x\partial_y$ na $\R^2$. Porównaj $\Phi_t^X\circ\Phi_s^Y$ i $\Phi_s^Y\circ\Phi_t^X$; wyznacz wyraz rzędu $st$ ich różnicy i powiąż go z nawiasem Liego.
\wskazowka{Zapisz oba przepływy jako macierze ścinające i porównaj ich iloczyny do wyrazów rzędu $st$. Znak zależy od kolejności składania.}
\item \textbf{Pole z globalnym hamowaniem.} Rozwiąż $\dot x=-x^3$ dla wszystkich danych początkowych. Ustal maksymalne przedziały istnienia, zupełność pola oraz granice orbit dla $t\to\pm\infty$.
\wskazowka{Dla niezerowego $x$ rozdziel zmienne i zwróć uwagę na ujemne czasy: zachowanie przy $t\to+\infty$ nie wystarcza do zupełności.}
\item \textbf{Pola na okręgu.} Dla $X(\theta)=\sin\theta\,\partial_\theta$ na $S^1$ sklasyfikuj punkty stacjonarne, kierunek pozostałych orbit i ich zbiory graniczne. Rozstrzygnij, czy pole jest zupełne.
\item \textbf{Prędkość a zupełność.} Na $\R$ porównaj pola $X=\partial_x$ i $Y=(1+x^2)\partial_x$. Mają te same zorientowane orbity. Oblicz czas osiągnięcia $+\infty$ i ustal, które z pól jest zupełne.
\wskazowka{Czas przebiegania otrzymasz z całki $\int_{x_0}^{R}du/(1+u^2)$. O zupełności decyduje możliwość przedłużenia rozwiązania w obie strony czasu.}
\item \textbf{Całka pierwsza.} Niech $H(x,y)=2x^2+y^2$ oraz $X=H_y\partial_x-H_x\partial_y$. Wyznacz $X$, sprawdź $X(H)=0$, rozwiąż równania ruchu i opisz wszystkie niezerowe orbity.
\item \textbf{Przeniesienie pola przez ścinanie.} Dla dyfeomorfizmu $F(x,y)=(x,y+x^3)$ oblicz $F_*(\partial_x+y\partial_y)$ w nowych współrzędnych i sprawdź wynik na funkcjach współrzędnych.
\item \textbf{Niezmienniczy podzbiór.} Dla $X=x(1-y)\partial_x+y(x-1)\partial_y$ określ, czy proste $x=0$, $y=0$ i hiperbola $xy=1$ są niezmiennicze. Uzasadnij odpowiedzi, badając styczność pola.
\wskazowka{Podzbiór dany równaniem $G=0$ jest niezmienniczy, jeśli $X(G)$ znika na tym zbiorze; osobno sprawdź każdy kandydat.}
\item \textbf{Równanie na transporcie funkcji.} Rozwiąż $\partial_tu+x\partial_yu=0$ na $\R^2$ z warunkiem $u(0,x,y)=e^{-x^2-y^2}$. Użyj przepływu pola i zweryfikuj otrzymaną funkcję przez podstawienie.
\item \textbf{Pochodna Liego kowektora.} Dla $X=x\partial_x-y\partial_y$ i $\alpha=y\,dx+x\,dy$ oblicz $\mathcal L_X\alpha$ bezpośrednio oraz za pomocą wzoru Cartana $\mathcal L_X=i_X\dd+\dd i_X$.
\item \textbf{Pole zależne od czasu.} Na $\R^2$ rozważ $X_t=-t y\partial_x+t x\partial_y$. Znajdź ewolucję $\Phi_{t,0}$ i rozstrzygnij, dlaczego na ogół nie zachodzi $\Phi_{t+s,0}=\Phi_{t,0}\circ\Phi_{s,0}$.
\wskazowka{Macierze pól $X_t$ są wielokrotnościami tej samej macierzy obrotu. Całkuj zależny od czasu kąt, zamiast podstawiać $t$ do przepływu pola autonomicznego.}
\item \textbf{Izotopia dylatacji nieliniowej.} Dla $F_t(x)=x+t\tanh x$, $0\le t\le1$, sprawdź, że $F_t$ jest dyfeomorfizmem $\R$. Wyznacz generator $X_t=(\partial_tF_t)\circ F_t^{-1}$ i podaj jego ograniczenia przy $x\to\pm\infty$.
\wskazowka{Wykorzystaj $F_t'(x)>0$ i granice przy $x\to\pm\infty$. Generator można podać niejawnie, przez punkt $u=F_t^{-1}(x)$.}
\item \textbf{Warunek zachowania objętości.} Niech $X=x^2\partial_x-2xy\partial_y$ na $\R^2$. Oblicz $\mathcal L_X(dx\wedge dy)$ i sprawdź, czy lokalny przepływ zachowuje standardowe pole powierzchni.
\item \textbf{Całkowalność dystrybucji.} Na $\R^3$ porównaj dystrybucje $D_1=\ker(dz-x\,dy)$ oraz $D_2=\ker(dz-x\,dx)$. Wyznacz lokalne bazy i ich nawiasy. Dla której dystrybucji istnieją lokalne powierzchnie całkowe przez każdy punkt?
\wskazowka{Zastosuj kryterium Frobeniusa $\alpha\wedge d\alpha=0$ dla jednowymiarowego anihilatora; sprawdź też jedną próbę całkowania równań Pfaffa.}
\item \textbf{Dwie całki pierwsze.} Dla pola $X=(y-z)\partial_x+(z-x)\partial_y+(x-y)\partial_z$ znajdź dwie funkcyjnie niezależne całki pierwsze na odpowiednim podzbiorze $\R^3$ i opisz geometrycznie krzywe całkowe. Zbadaj zupełność pola.
\wskazowka{Poszukaj wielkości liniowej $x+y+z$ oraz kwadratowej $x^2+y^2+z^2$. Wektorowe równanie ma stałe współczynniki, co pomaga przy zupełności.}
\end{enumerate}


\chapter{Algebra tensorowa i zewnętrzna}
\begin{enumerate}
\item \textbf{Rząd tensorowy.} Tensor $T\in\R^3\otimes\R^2$ ma macierz współrzędnych $\begin{psmallmatrix}1&2\\2&4\\3&6\end{psmallmatrix}$. Zapisz go jako tensor prosty. Powtórz dla macierzy $\begin{psmallmatrix}1&2\\2&4\\3&5\end{psmallmatrix}$ i wyznacz minimalną liczbę tensorów prostych w rozkładzie.
\item \textbf{Dwa sposoby kontrakcji.} Dla $T^i{}_{jk}=a^i b_jc_k+\delta^i_kd_j$ w wymiarze $n$ oblicz osobno kontrakcje indeksów $(i,j)$ oraz $(i,k)$. Podaj warunki, przy których oba otrzymane kowektory są równe.
\wskazowka{Wypisz sumowanie po powtarzającym się indeksie oddzielnie w obu przypadkach. Zauważ, że symbol Kroneckera kontrahuje inne czynniki w każdej z dwóch operacji.}
\item \textbf{Iloczyn tensorowy endomorfizmów.} Dla $A=\begin{psmallmatrix}1&1\\0&2\end{psmallmatrix}$ i $B=\begin{psmallmatrix}0&1\\-1&0\end{psmallmatrix}$ wyznacz macierz $A\otimes B$ oraz jej wyznacznik i ślad. Porównaj otrzymane liczby z wyznacznikami i śladami $A,B$.
\item \textbf{Tensor symetryczny i forma kwadratowa.} Dla formy $q(x,y,z)=2x^2+4xy+3y^2-z^2$ znajdź tensor symetryczny $B$, dla którego $q(v)=B(v,v)$. Wyznacz sygnaturę $B$ metodą dopełniania kwadratów.
\item \textbf{Antysymetryzacja w wymiarze trzy.} Zapisz pełną antysymetryzację tensora $T=e^1\otimes e^2\otimes e^3+2e^2\otimes e^3\otimes e^1$. Oblicz współczynnik przy $e^1\wedge e^2\wedge e^3$, przyjmując konwencję $\operatorname{Alt}(T)=\frac1{3!}\sum_{\sigma}\operatorname{sgn}(\sigma)\,\sigma T$.
\item \textbf{Warunki na biwektor.} Niech $u=e_1\wedge e_2+a e_1\wedge e_3+b e_2\wedge e_4+c e_3\wedge e_4$ w $\Lambda^2\R^4$. Wyznacz wszystkie $(a,b,c)$, dla których $u$ jest rozkładalny, i skonstruuj wtedy dwa wektory $v,w$ takie, że $u=v\wedge w$.
\wskazowka{W czterech wymiarach dwuwektor jest rozkładalny dokładnie wtedy, gdy $u\wedge u=0$. Do konstrukcji $v,w$ wybierz współczynnik przy $e_1\wedge e_2$ jako niezerowy.}
\item \textbf{Minory odwzorowania.} Dla $A:\R^3\to\R^2$ o macierzy $\begin{psmallmatrix}1&2&-1\\0&3&4\end{psmallmatrix}$ oblicz macierz $\Lambda^2A:\Lambda^2\R^3\to\Lambda^2\R^2$. Określ jej jądro i powiąż je z przestrzenią $\ker A$.
\wskazowka{Macierz $\Lambda^2 A$ tworzą minory $2\times2$ macierzy $A$. Związek z jądrem $A$ zobacz, badając $v\wedge w$ dla $v\in\ker A$.}
\item \textbf{Tożsamość Cauchy'ego--Bineta w liczbach.} Niech $A=\begin{psmallmatrix}1&2&0\\0&1&3\end{psmallmatrix}$ i $B=\begin{psmallmatrix}2&0\\1&1\\0&-1\end{psmallmatrix}$. Oblicz $\det(AB)$ oraz sumę iloczynów odpowiednich minorów rzędu dwa, wyjaśniając ich pochodzenie z $\Lambda^2(AB)$.
\item \textbf{Parowanie na potęgach zewnętrznych.} Dla $\alpha_1=dx+dy$, $\alpha_2=dy+dz$ oraz $v_1=(1,0,1)$, $v_2=(0,2,1)$ oblicz $(\alpha_1\wedge\alpha_2)(v_1,v_2)$ na dwa sposoby: przez wyznacznik i przez rozwinięcie formy.
\item \textbf{Dualna baza po zmianie bazy.} W $V=\R^3$ wprowadź bazę $f_1=e_1+e_2$, $f_2=e_2+e_3$, $f_3=e_1+e_3$. Wyznacz dualną bazę $(f^1,f^2,f^3)$ i wartość $f^1\wedge f^2\wedge f^3(e_1,e_2,e_3)$.
\item \textbf{Tensor i zmiana mapy.} Dla $T=x\,dx\otimes\partial_y+y\,dy\otimes\partial_x$ na dodatniej ćwiartce $\R^2$ znajdź współrzędne tensora w mapie $u=\log x$, $v=\log y$ i oblicz jego ślad w obu mapach.
\wskazowka{Przekształć zarówno wektory bazowe, jak i kowektory: $\partial_x=x^{-1}\partial_u$ oraz $dx=x\,du$. Ślad endomorfizmu nie zależy od mapy.}
\item \textbf{Wykładnicza rodzina odwzorowań.} Niech $F_t(x,y)=(e^t x,e^{-t}y)$. Oblicz $(F_t)^*(dx\wedge dy)$ oraz $(F_t)^*(dx\otimes dy)$; zbadaj, które z tych tensorów są niezmiennicze.
\item \textbf{Równanie dla pochodnej Liego.} Znajdź wszystkie gładkie funkcje $f(x)$ na $\R$, dla których tensor $T=f(x)\,dx\otimes dx$ spełnia $\mathcal L_{x\partial_x}T=0$. Rozstrzygnij, czy istnieje niezerowe rozwiązanie gładkie w $x=0$.
\wskazowka{Użyj $\mathcal L_X(dx)=d(Xx)$ i reguły Leibniza. Po otrzymaniu ODE na $f$ zbadaj gładkość rozwiązania w początku.}
\item \textbf{Ślad iloczynu tensorowego.} Niech $T$ odpowiada operatorowi $A:V\to V$, a $S$ operatorowi $B:W\to W$. Oblicz kontrakcję $T\otimes S$ wzdłuż pary indeksów należących do $V$ i porównaj wynik ze śladem $A$. Zastosuj do $A=\operatorname{diag}(1,2)$ i $B=\begin{psmallmatrix}0&1\\1&0\end{psmallmatrix}$.
\item \textbf{Iloczyn zewnętrzny i niezależność liniowa.} Niech $\alpha=dx+2dy$, $\beta=dy+3dz$, $\gamma=2dx+dz$ na $\R^3$. Oblicz $\alpha\wedge\beta\wedge\gamma$ i na tej podstawie rozstrzygnij, czy formy $\alpha,\beta,\gamma$ tworzą bazę przestrzeni kostycznej.
\item \textbf{Kontrakcja dwóch tensorów typu $(0,3)$.} Załóż, że $V$ ma ustalony iloczyn skalarny i $A,B\in(V^*)^{\otimes3}$. Zbuduj za pomocą metryki tensor typu $(0,4)$ otrzymany przez kontrakcję jednego indeksu. Wyjaśnij, dlaczego bez dodatkowej struktury nie ma kanonicznego sposobu uzyskania w ten sposób tensora typu $(0,3)$.
\wskazowka{Do kontrakcji dwóch dolnych indeksów potrzebujesz odwrotnej metryki $g^{ij}$. Policz typ tensorowy przed i po kontrakcji; sama parzystość sumy rzędów jest istotna.}
\end{enumerate}


\chapter{Grassmanniany, algebra Clifforda i spinory}
\begin{enumerate}
\item \textbf{Płaszczyzny jako wykresy.} Dla $A_t:\R^2\to\R^2$ o macierzy $\begin{psmallmatrix}t&t^2\\0&2t\end{psmallmatrix}$ wyznacz krzywą $L_t=\operatorname{graph}(A_t)$ w $\Gr_2(\R^4)$, jej wektor styczny w $t=0$ oraz zmiany współrzędnych Plückera do pierwszego rzędu.
\item \textbf{Przecięcie map grassmannianu.} Niech $L\subset\R^4$ będzie rozpięta przez $v=(1,0,1,1)$ i $w=(0,1,2,1)$. Oblicz macierze wykresowe $L$ względem dwóch różnych rozkładów $\R^4=\R^2\oplus\R^2$, o ile obie mapy są określone; zapisz przejście między nimi.
\wskazowka{Każda mapa Grassmannianu wymaga odwracalności odpowiedniej macierzy projekcji $L\to\R^2$. Macierz przejścia wyznacz przez rozwiązanie równania blokowego.}
\item \textbf{Wiązka tautologiczna nad krzywą.} Niech $L_t=\operatorname{span}(\cos t,\sin t,1)\in\Gr_1(\R^3)$, $t\in\R/2\pi\Z$. Zbadaj, czy wiązka $L_t$ nad okręgiem ma sekcję nigdzie niezerową; skonstruuj taką sekcję albo wykaż niemożliwość.
\item \textbf{Suma podprzestrzeni.} Dla $L_A=\operatorname{graph}(A)$ i $L_B=\operatorname{graph}(B)$ w $\R^k\oplus\R^k$ wykaż, że $L_A\cap L_B=\{0\}$ wtedy i tylko wtedy, gdy $A-B$ jest odwracalna. Zastosuj do dwóch zadanych przez siebie macierzy $2\times2$.
\item \textbf{Plücker i równania kwadratowe.} Zbadaj, czy $u=e_1\wedge e_2+e_1\wedge e_4+e_2\wedge e_3+\lambda e_3\wedge e_4$ reprezentuje płaszczyznę w $\R^4$. Wyznacz dopuszczalne $\lambda$ i odzyskaj bazę płaszczyzny.
\wskazowka{Wykorzystaj relację Plückera albo kryterium $u\wedge u=0$. Następnie szukaj bazy w postaci wykresu względem dogodnej pary osi.}
\item \textbf{Gauss dla wykresu.} Dla powierzchni $z=x^2+2y^2$ w $\R^3$ oblicz odwzorowanie Gaussa w otoczeniu $(0,0)$, jego różniczkę w początku i określ rząd tej różniczki.
\item \textbf{Równanie Clifforda.} W $\Cl(\R^3,-\langle\cdot,\cdot\rangle)$ uprość $(e_1+e_2)(e_1-e_2)$, $(e_1e_2)(e_2e_3)$ oraz $(e_1+e_2+e_3)^2$. Porównaj ostatni wynik z kwadratem normy wektora.
\item \textbf{Odwracalność wektora.} Dla dowolnego niezerowego $v\in\R^{p,q}\subset\Cl_{p,q}$ o $q(v)\ne0$ znajdź jego odwrotność. Podaj przykład wektora niezerowego i nieodwracalnego w algebrze Clifforda formy o sygnaturze $(1,1)$.
\item \textbf{Elementy parzyste.} W $\Cl_3$ (konwencja $e_i^2=-1$) oblicz kwadraty $e_1e_2$, $e_2e_3$, $e_3e_1$ oraz ich iloczyny w obu kolejnościach. Zapisz element $(1+e_1e_2)/\sqrt2$ i oblicz jego odwrotność.
\wskazowka{Zamiast mnożyć wszystkie wyrazy od zera, używaj relacji $e_ie_j=-e_je_i$ dla $i\ne j$ i $e_i^2=-1$.}
\item \textbf{Obroty wokół wybranej osi.} W $\Cl_3$ rozważ $s=\cos(\theta/2)+\sin(\theta/2)e_1e_2$. Oblicz sprzężenie $v\mapsto svs^{-1}$ na $e_1,e_2,e_3$, a następnie dla $\theta=\pi/3$ wyznacz macierz odpowiadającej rotacji.
\wskazowka{Ogranicz rachunek do płaszczyzny rozpiętej przez $e_1,e_2$. Pamiętaj, że w konwencji $e_i^2=-1$ zwrot obrotu zależy od znaku w wykładniku spinowym.}
\item \textbf{Spinor jako kolumna.} Dla macierzy Pauliego $\sigma_1,\sigma_2,\sigma_3$ oblicz $\sigma_1\sigma_2\sigma_3$ i sprawdź antykomutatory $\sigma_i\sigma_j+\sigma_j\sigma_i$. Zbadaj działanie elementu $\exp(i\theta\sigma_3/2)$ na spinor $(1,1)^T$.
\item \textbf{Chiralność w konkretnym modelu.} Niech $\sigma_j$ oznaczają macierze Pauliego, $c(e_1)=i\sigma_1$, $c(e_2)=i\sigma_2$, a $\Gamma=-i c(e_1)c(e_2)$. Oblicz $\Gamma$, jej wartości własne i rzutowania na podprzestrzenie własne. Sprawdź działanie $c(e_1)$ na każdej z nich.
\wskazowka{Wyznacz najpierw iloczyn $\sigma_1\sigma_2$; do rzutowania na przestrzenie chiralne użyj $(I\pm\Gamma)/2$.}
\item \textbf{Obrót o $2\pi$ spinora.} Rozpatrz pętlę obrotów wokół osi $z$ od kąta $0$ do $2\pi$. Oblicz podniesienie tej pętli do $\Spin(3)$ z wartością początkową $1$ i rozstrzygnij, czy jest pętlą zamkniętą.
\wskazowka{W $\Spin(3)$ obrót o kąt $\theta$ reprezentuje element zależny od połowy kąta. Porównaj wartości jego podniesienia w punktach $0$ i $2\pi$.}
\item \textbf{Spin a suma wiązek.} Załóż, że rzeczywiste zorientowane wiązki $E,F$ spełniają $w_2(E)=0$ i $w_2(F)=0$. Korzystając ze wzoru Whitneya dla klas Stiefela--Whitneya, zbadaj $w_2(E\oplus F)$. Co się zmienia, gdy wiązki nie są orientowalne?
\wskazowka{Rozwiń $w(E\oplus F)=w(E)\smile w(F)$ do stopnia drugiego; nie pomijaj składnika $w_1(E)w_1(F)$ w przypadku nieorientowalnym.}
\item \textbf{Struktura spinowa na cylindrze.} Rozważ płaski cylinder $S^1\times\R$. Opisz dwie możliwe struktury spinowe jako warunki sklejenia spinorów po pełnym obrocie wokół $S^1$. Dla każdej z nich rozstrzygnij, czy istnieje niezerowy spinor równoległy względem płaskiej koneksji.
\wskazowka{Dwa podniesienia holonomii okręgu do $\Spin(2)$ różnią się znakiem. Spinor równoległy musi wrócić do tej samej wartości po sklejeniu.}
\end{enumerate}


\chapter{Podstawy algebry abstrakcyjnej}
\begin{enumerate}
\item \textbf{Prezentacja z redundantnymi generatorami.} Przekształć prezentację $G=\langle a,b,c\mid c=ab,\ a^3=1,\ b^2=1,\ ab=ba\rangle$ do prezentacji na dwóch generatorach. Wyznacz rząd grupy i sprawdź, czy jest cykliczna.
\item \textbf{Abelianizacja grupy z relacją.} Oblicz abelianizację $\langle a,b\mid a^2b^{-3}=1\rangle$ i ustal, czy jest grupą wolną abelową. Wykorzystaj opis abelianizacji przez iloraz $\Z^2$.
\wskazowka{Abelianizacja zamienia relację grupową na relację liniową w $\Z^2$. Wykorzystaj postać Smitha jednoelementowej macierzy relacji.}
\item \textbf{Iloczyn półprosty w nowym przykładzie.} Niech $\Z/2\Z$ działa na $\Z/5\Z$ przez $x\mapsto-x$. Oblicz prawo mnożenia, centrum i rzędy elementów w $G=(\Z/5\Z)\rtimes(\Z/2\Z)$.
\item \textbf{Dwa działania, dwa iloczyny półproste.} Dla $N=\Z/8\Z$ rozważ działania grupy $\Z/2\Z$ dane przez automorfizmy $x\mapsto-x$ i $x\mapsto3x$. Porównaj otrzymane grupy rzędu $16$, licząc liczbę elementów rzędu $2$.
\wskazowka{W grupie półprostej kwadrat elementu $(x,1)$ zależy od wybranego automorfizmu $\Z/8$. Porównaj liczbę rozwiązań równania $(x,1)^2=e$.}
\item \textbf{Homomorfizmy z grupy wolnej.} Zlicz homomorfizmy $F(a,b,c)\to S_3$. Ile z nich jest surjektywnych? Uzasadnij odpowiedź, korzystając z możliwych obrazów generatorów.
\item \textbf{Jądro homomorfizmu do grupy cyklicznej.} Dla $f:\Z^3\to\Z/12\Z$, $f(a,b,c)=[2a+3b+5c]$, wyznacz obraz i opisz jądro jako podgrupę $\Z^3$ z trzema generatorami.
\wskazowka{Najpierw sprawdź surjektywność przez $\gcd(2,3,5,12)$. Potem sprowadź równanie kongruencyjne do bazy jądra przy użyciu operacji całkowitoliczbowych.}
\item \textbf{Pierścienie ilorazowe.} Wyznacz wszystkie elementy idempotentne w $\Z/30\Z$. Porównaj wynik z rozkładem chińskiego twierdzenia o resztach.
\item \textbf{Nilpotenty.} W $A=\R[x,y]/(x^2,xy,y^3)$ wyznacz bazę $\R$-liniową, wszystkie nilpotentne elementy i jedyny ideał maksymalny. Sprawdź, które elementy są odwracalne.
\item \textbf{Homomorfizm modułów.} Niech $A=\begin{psmallmatrix}2&4\\6&8\end{psmallmatrix}$ określa $f:\Z^2\to\Z^2$. Wyznacz $\ker f$, $\operatorname{coker}f$ i rangę obrazu, korzystając z postaci normalnej Smitha.
\item \textbf{Tensorowanie modułów skończonych.} Oblicz $(\Z/8\Z\oplus\Z/9\Z)\otimes_\Z(\Z/12\Z\oplus\Z/15\Z)$ i zapisz wynik w postaci sumy prostych grup cyklicznych.
\wskazowka{Rozdziel iloczyn tensorowy po sumach prostych i użyj $\Z/m\otimes_\Z\Z/n\cong\Z/\gcd(m,n)$.}
\item \textbf{Homomorfizmy między modułami cyklicznymi.} Opisz $\Hom_\Z(\Z/18\Z,\Z/12\Z)$ oraz $\Hom_\Z(\Z/12\Z,\Z)$, z podaniem rzędu każdej grupy.
\item \textbf{Lokalizacja.} Niech $S=\{5^k:k\ge0\}$. Oblicz $S^{-1}(\Z/20\Z)$ i $S^{-1}(\Z/18\Z)$ jako moduły nad $\Z[1/5]$.
\wskazowka{W lokalizacji wielokrotności potęg $5$ stają się odwracalne. Rozbij moduł na składowe pierwotne i sprawdź, które z nich zanikają.}
\item \textbf{Lokalizacja i jądro.} Dla $f:\Z\to\Z/15\Z$, $f(n)=[n]$, porównaj jądro $S^{-1}f$ przy $S=\{3^k:k\ge0\}$ z lokalizacją $S^{-1}\ker f$.
\wskazowka{Zastosuj dokładność funktora lokalizacji, ale pokaż też bezpośrednio, jak wygląda jądro po odwróceniu $3$.}
\item \textbf{Ranga modułów.} Niech $M$ będzie zadany przez $\Z^3/\langle(2,0,0),(0,0,6)\rangle$. Oblicz jego rangę i podmoduł torsyjny oraz porównaj $M\otimes\Q$ z $M\otimes\Z/2\Z$.
\item \textbf{Jawne rozszczepienie.} Dla surjekcji $f:\Z^3\to\Z^2$, $f(a,b,c)=(a+2b,b+3c)$ skonstruuj homomorfizm $s:\Z^2\to\Z^3$ spełniający $f\circ s=\id$. Wyznacz generator $\ker f$ i zapisz jawny izomorfizm $\Z^3\cong\ker f\oplus\Z^2$.
\wskazowka{Szukaj obrazów wektorów bazowych $\Z^2$, które odwzorowanie $f$ wysyła na odpowiednie wektory standardowe. Potem rozłóż dowolny $v$ na część w jądrze i w obrazie sekcji.}
\end{enumerate}


\chapter{Algebra homologiczna: kompleksy i (ko)homologie}
\begin{enumerate}
\item \textbf{Kompleks z trzema stopniami.} Niech $C_2=\Z^2$, $C_1=\Z^3$, $C_0=\Z$, przy czym $\partial_2(a,b)=(2a,0,3b)$, $\partial_1(x,y,z)=y$. Wyznacz wszystkie grupy $H_k(C)$ i wskaż reprezentantów ich generatorów.
\item \textbf{Dobieranie różniczki.} Niech $C_2=\Z$, $C_1=\Z^2$, $C_0=\Z$, $\partial_2(1)=(4,6)$. Wyznacz wszystkie homomorfizmy $\partial_1:\Z^2\to\Z$ spełniające $\partial_1\partial_2=0$. Wybierz taki, który ma obraz $2\Z$, i oblicz homologię.
\wskazowka{Warunek $\partial_1\partial_2=0$ to pojedyncze równanie diofantyczne na współczynniki funkcjonału. Do homologii użyj bazy $\ker\partial_1$.}
\item \textbf{Dwie redukcje współczynników.} Dla kompleksu $C_1=\Z^2\xrightarrow{\operatorname{diag}(4,9)}C_0=\Z^2$ oblicz homologię nad $\Z$, po tensorowaniu z $\Z/2\Z$ oraz po tensorowaniu z $\Z/3\Z$. Wskaż zjawiska, których nie wykrywa samo tensorowanie grup homologii.
\wskazowka{Po tensorowaniu różniczki redukują się modulo $p$. Współczynniki mogą tworzyć nowe cykle; porównaj wynik z członem $\operatorname{Tor}$ w twierdzeniu o współczynnikach uniwersalnych.}
\item \textbf{Kohomologia dualnego kompleksu.} Dla $C_1=\Z^2\xrightarrow{\begin{psmallmatrix}3&6\\0&5\end{psmallmatrix}}C_0=\Z^2$ oblicz $H_*(C)$ i $H^*(\Hom(C,\Z))$. Wyjaśnij, w jakich stopniach pojawia się torsja.
\item \textbf{Łańcuchowa mapa diagonalna.} Niech $C_1=\Z\xrightarrow{\cdot12}C_0=\Z$. Dla liczb całkowitych $a,b$ rozstrzygnij, kiedy $f_1=\cdot a$, $f_0=\cdot b$ wyznaczają endomorfizm kompleksu. Wyznacz wszystkie możliwe endomorfizmy indukowane na $H_0(C)$.
\item \textbf{Stożek macierzy.} Zbuduj stożek odwzorowania $f:\Z^2[0]\to\Z^2[0]$ o macierzy $\begin{psmallmatrix}2&1\\0&3\end{psmallmatrix}$. Oblicz jego homologię i powiąż z jądrem oraz kokernlem $f$.
\wskazowka{Dla kompleksu skoncentrowanego w stopniu zerowym stożek ma tylko dwa niezerowe stopnie. Sprawdź dokładnie przyjętą konwencję znaków różniczki stożka.}
\item \textbf{Acykliczny, lecz nieściągalny kompleks.} Zbuduj dokładny kompleks długości trzy z modułów nad $\Z/4\Z$, który nie jest łańcuchowo homotopijny z zerem. Uzasadnij brak homotopii, korzystając z konsekwencji rozszczepialności.
\wskazowka{Wykorzystaj nierozszczepialny krótki ciąg dokładny $0\to 2R\to R\to R/(2)\to0$ dla $R=\Z/4\Z$. Kompleks ściągalny wymuszałby istnienie odpowiedniego rozszczepienia.}
\item \textbf{Homotopia łańcuchowa w macierzach.} Na kompleksie $0\to\R\xrightarrow{\binom{1}{0}}\R^2\xrightarrow{(0\ \ 1)}\R\to0$ skonstruuj jawny operator stopnia $+1$, $h$, spełniający $\partial h+h\partial=\id$.
\wskazowka{Zacznij od prawego końca kompleksu, dobierając prawą odwrotność surjekcji; następnie zbuduj homotopię na składniku dopełniającym obraz.}
\item \textbf{Pięciowyrazowy ciąg.} Załóż, że $0\to A\to\Z^3\xrightarrow{(2,4,6)}\Z\to B\to0$ jest dokładny. Oblicz $A$ i $B$ i ustal, czy $A$ jest wolny. Wskaż, jakie informacje daje sama znajomość rzędów homomorfizmów.
\item \textbf{Długi ciąg na konkretnym ilorazie.} Niech $B$ będzie kompleksem $B_1=\Z\xrightarrow{\cdot6}B_0=\Z$, a $A\subset B$ podkompleksem zadanym przez $A_1=0$ i $A_0=3\Z$. Oblicz homologie $A$, $B$, $B/A$ oraz homomorfizm łączący.
\wskazowka{Oblicz najpierw $B/A$ w każdym stopniu. Homomorfizm łączący polega na podniesieniu cyklu z ilorazu, zastosowaniu różniczki i odczytaniu klasy w $A$.}
\item \textbf{Filtracja trójstopniowa.} Skonstruuj filtrację kompleksu $C$ z zadania 1 przez podkompleksy, w której w kolejnych krokach dodajesz generatory w stopniach $0,1,2$. Zapisz względne grupy homologii kolejnych par i wyjaśnij, kiedy nowy generator zabija wcześniejszą klasę.
\wskazowka{W filtracji przez szkielety każda nowa komórka daje generator względny, ale jego brzeg może trafić do wcześniejszego filtru.}
\item \textbf{Mayer--Vietoris dla podkompleksów.} Niech $C$ będzie kompleksem symplicjalnym dwóch trójkątów mających jedną wspólną krawędź. Przyjmij za $A,B$ podkompleksy pojedynczych wypełnionych trójkątów, tak że $A\cup B=C$ i $A\cap B$ jest jedną krawędzią. Wypisz łańcuchowy ciąg Mayera--Vietorisa i oblicz jego homomorfizm łączący w stopniu $1$.
\item \textbf{Charakterystyka Eulera bez obliczania homologii.} Kompleks skończonych przestrzeni nad ciałem ma wymiary $(\dim C_3,\dim C_2,\dim C_1,\dim C_0)=(2,5,7,3)$. Załóż, że $H_3=H_0=0$ i $\dim H_2=1$. Wyznacz $\dim H_1$ albo wykaż, że dane są sprzeczne.
\item \textbf{Zastosowanie lematu pięciu.} W przemiennym diagramie długich ciągów dokładnych załóż, że znasz izomorfizmy na czterech spośród pięciu kolejnych grup. Sformułuj precyzyjnie, jakie warunki wystarczą, by uzyskać izomorfizm w środku. Znajdź przykład pokazujący, że osłabienie jednego warunku do samej różnowartościowości nie zawsze wystarcza.
\wskazowka{Rozdziel twierdzenie pięciu na dowody różnowartościowości i surjektywności za pomocą pościgu za diagramem; osłabione założenie sprawdź na prostych ciągach zerowych.}
\item \textbf{Torsja a charakterystyka Eulera.} Porównaj kompleksy $0\to\Z\xrightarrow{\cdot m}\Z\to0$ dla $m=0,1,7$. Oblicz homologię i charakterystykę Eulera każdego z nich. Wytłumacz, dlaczego ostatnia nie odróżnia tych przypadków.
\item \textbf{Uwolnienie torsji przez współczynniki.} Dla kompleksu $0\to\Z\xrightarrow{6}\Z\to0$ oblicz jego homologie po tensorowaniu z $\Z/4$ i $\Z/9$. Porównaj z przewidywaniami twierdzenia o współczynnikach uniwersalnych i odpowiednimi grupami $\operatorname{Tor}$.
\wskazowka{Dla homomorfizmu mnożenia przez $6$ po redukcji modulo $m$ oblicz oddzielnie jądro i obraz. Porównaj z $\operatorname{Tor}_1^{\Z}(\Z/6,\Z/m)$.}
\end{enumerate}


\chapter{Homotopia, nakrycia i CW-kompleksy}
\begin{enumerate}
\item \textbf{Uzupełniony graf.} Niech $X$ powstanie z okręgu przez dołączenie odcinka, którego oba końce przyklejamy do dwóch różnych punktów okręgu. Wyznacz $\pi_1(X)$ oraz liczbę niezależnych cykli grafu.
\item \textbf{Trzy dziury w płaszczyźnie.} Dla $X=\R^2\setminus\{(0,0),(1,0),(0,1)\}$ skonstruuj konkretne pętle generujące $\pi_1(X)$, oblicz obraz komutatora pierwszych dwóch generatorów w abelianizacji i uzasadnij, dlaczego komutator nie jest trywialny w $\pi_1(X)$.
\item \textbf{Przestrzeń z dwoma komórkami.} Niech $X_m=S^1\cup_{z\mapsto z^m}D^2$ dla $m\ge2$. Oblicz $\pi_1(X_m)$. Następnie dołącz drugą $2$-komórkę odwzorowaniem stopnia $n$ i wyznacz grupę podstawową nowej przestrzeni.
\item \textbf{Nakrycie skończonego grafu.} Zbuduj spójne nakrycie stopnia $3$ bukietu dwóch okręgów, opisując działanie generatorów na trzy punkty włókna przez permutacje. Wyznacz grupę automorfizmów dla swojego przykładu.
\wskazowka{Nakrycie grafu zakoduj dwiema permutacjami włókna. Grupa przekształceń pokładowych jest centralizatorem działania monodromii w $S_3$.}
\item \textbf{Nakrycia torusa z macierzy.} Niech $A=\begin{psmallmatrix}2&1\\0&3\end{psmallmatrix}$, $f_A:\T^2\to\T^2$, $f_A([x])=[Ax]$. Wyznacz stopień nakrycia, grupę przekształceń pokładowych i podgrupę $\pi_1(\T^2)$ odpowiadającą nakryciu.
\wskazowka{Podgrupą odpowiadającą nakryciu jest $A\Z^2\subset\Z^2$. Grupę przekształceń pokładowych znajdziesz jako $\Z^2/A\Z^2$.}
\item \textbf{Podnoszenie do nakrycia torusa.} Dla macierzy całkowitych $A,B\in M_2(\Z)$ i $\det A\ne0$ znajdź warunek konieczny i dostateczny na istnienie $g:\T^2\to\T^2$ z $f_A\circ g=f_B$. Sprawdź go dla $A=\operatorname{diag}(2,3)$ oraz $B=\begin{psmallmatrix}4&2\\0&6\end{psmallmatrix}$.
\wskazowka{Po przejściu na grupy podstawowe warunek podniesienia sprowadza się do zawierania $B\Z^2\subseteq A\Z^2$; następnie uwzględnij ewentualne przesunięcie podniesienia.}
\item \textbf{Nakrycie niegaloisowskie.} Skonstruuj spójne trójkrotne nakrycie bukietu dwóch okręgów, którego grupa przekształceń pokładowych jest trywialna. Użyj dwóch odpowiednio dobranych permutacji z $S_3$.
\wskazowka{Wybierz permutacje generujące przechodnie działanie z trywialnym centralizatorem, na przykład transpozycję i cykl długości trzy.}
\item \textbf{Prezentacja z komórek.} Zbuduj $2$-kompleks CW o jednej $0$-komórce, dwóch $1$-komórkach i jednej $2$-komórce, której odwzorowanie dołączające opisuje słowo $a^2b^3$. Zapisz prezentację $\pi_1$ i oblicz jego abelianizację.
\wskazowka{Relację dołączającej komórki odczytaj z wyrazu w grupie wolnej. Abelianizację otrzymasz, zastępując ją wektorem wykładników.}
\item \textbf{Połączenie dwóch torusów punktem.} Wyznacz $\pi_1(\T^2\vee\T^2)$ przez prezentację. Rozstrzygnij, czy grupa jest abelowa, i znajdź jej epimorfizm na wolną grupę rzędu dwa.
\item \textbf{Sklejanie walców.} Dwa cylindry $S^1\times[0,1]$ sklej wzdłuż jednego okręgu brzegowego odwzorowaniem stopnia $\pm1$. Ustal typ homotopii wyniku. Co zmieni się, jeśli użyjesz stopnia $2$ i zamiast homeomorfizmu rozważysz iloraz sklejony wzdłuż takiego odwzorowania?
\item \textbf{Wolne grupy i drzewa.} Niech $G$ będzie grafem powstałym z pięciokąta przez dodanie dwóch przekątnych nieprzecinających się poza wierzchołkami. Znajdź drzewo rozpinające i bazę $\pi_1(G)$; opisz zmianę bazy przy wyborze innego drzewa.
\item \textbf{Pętla a stopień.} Dla $f:S^1\to S^1$, $f(e^{it})=e^{i(5t+\sin t)}$, oblicz $f_*:\pi_1(S^1)\to\pi_1(S^1)$. Rozstrzygnij, przez które z nakryć $z\mapsto z^n$ odwzorowanie $f$ ma podniesienie.
\item \textbf{Grupa podstawowa sumy spójnej powierzchni.} Skonstruuj CW-model orientowalnej powierzchni rodzaju $2$ z jednym wierzchołkiem i jedną $2$-komórką. Wyznacz prezentację grupy podstawowej i jej abelianizację.
\wskazowka{Użyj standardowego wielokąta ośmiobocznego z parami utożsamianych boków; słowo brzegu daje relację komórki dwuwymiarowej.}
\item \textbf{Graf a przyczepienie komórki.} Zacznij od bukietu trzech okręgów z generatorami $a,b,c$. Dołącz $2$-komórki o słowach $ab^{-1}$ i $bc^{-1}$. Oblicz grupę podstawową i porównaj ją z wynikiem po dołączeniu dodatkowej komórki o słowie $a^4$.
\item \textbf{Kontrprzykład do intuicji o $\pi_1$.} Znajdź dwa skończone CW-kompleksy o izomorficznej grupie podstawowej, które nie są homotopijnie równoważne. Wykaż brak równoważności za pomocą odpowiedniego niezmiennika z dalszych rozdziałów.
\item \textbf{Nakrycia butelki Kleina.} Dla prezentacji $\pi_1(K)=\langle a,b\mid aba^{-1}=b^{-1}\rangle$ wyznacz jądro homomorfizmu $\pi_1(K)\to\Z/2$ danego przez $a\mapsto1$, $b\mapsto0$. Zidentyfikuj odpowiadające mu dwukrotne nakrycie i oblicz jego grupę podstawową.
\wskazowka{Rozważ homomorfizm określający orientację; w jądrze sprawdź relację pomiędzy $a^2$ oraz $b$.}
\end{enumerate}


\chapter{Homologia przestrzeni}
\begin{enumerate}
\item \textbf{Pięciokąt bez wypełnienia.} Zbuduj macierz różniczki dla kompleksu symplicjalnego będącego brzegiem pięciokąta. Oblicz $H_0,H_1$ nad $\Z$ i porównaj wyniki po dodaniu wszystkich trójkątów triangulujących wnętrze.
\item \textbf{Dwa okręgi połączone krawędzią.} Dla grafu utworzonego z dwóch rozłącznych trójkątnych cykli przez dołączenie jednej krawędzi łączącej ich wierzchołki oblicz homologie symplicjalne i podaj jawny układ generatorów $H_1$.
\item \textbf{Kompleks z torsją.} Niech $X_m=S^1\cup_m e^2$ dla $m\ge2$. Zapisz wszystkie komórkowe operatory brzegowe i oblicz $H_k(X_m;\Z)$ oraz $H_k(X_m;\Z/p\Z)$ dla liczby pierwszej $p$.
\wskazowka{Użyj kompleksu komórkowego z pojedynczą różniczką mnożenia przez $m$. Po redukcji modulo $p$ operator może stać się zerowy.}
\item \textbf{Dwa dołączenia do okręgu.} Dla $X=S^1\cup_m e^2\cup_n e^2$ oblicz $H_2(X;\Z)$ i $H_1(X;\Z)$ jako funkcje $m,n\in\Z$. Wyjaśnij, skąd bierze się wolny składnik w $H_2$.
\wskazowka{Operator $\partial_2$ jest wierszem $(m,n)$. Jego jądro i kokernel oblicz przez największy wspólny dzielnik.}
\item \textbf{Kompleks o zadanej macierzy.} Skonstruuj $2$-wymiarowy CW-kompleks z jedną $0$-komórką i dwiema $1$- oraz dwiema $2$-komórkami, którego różniczka $\partial_2$ ma macierz $\begin{psmallmatrix}2&4\\0&6\end{psmallmatrix}$. Oblicz homologie całkowite.
\wskazowka{Macierz różniczki sprowadź do postaci Smitha operacjami elementarnymi nad $\Z$; odpowiednie słowa dołączania można wybrać w wolnej grupie dwóch generatorów.}
\item \textbf{Homologia powierzchni z brzegiem.} Rozważ orientowalną powierzchnię rodzaju $g$ z $b\ge1$ okręgami brzegowymi. Wyznacz jej grupy homologii i charakterystykę Eulera za pomocą odpowiedniego modelu CW.
\item \textbf{Produkt z okręgiem.} Oblicz $H_*(S^2\times S^1;\Z)$ na podstawie struktury CW produktu lub wzoru Künnetha. Wskaż konkretne cykle reprezentujące niezerowe klasy.
\item \textbf{Klein i współczynniki.} Oblicz $H_*(K;\Z/2\Z)$ dla butelki Kleina, zaczynając od jej komórkowego kompleksu nad $\Z$. Porównaj z $H_*(K;\Q)$ i wyjaśnij różnicę.
\wskazowka{W standardowym modelu CW butelki Kleina jedyna nietrywialna różniczka jest mnożeniem przez $2$. Zredukuj ją modulo wybrany pierścień współczynników.}
\item \textbf{Homologia po sklejeniu dwóch dysków.} Sklej dwa dyski $D^2$ wzdłuż całych brzegów przez odwzorowanie stopnia $-1$. Zidentyfikuj przestrzeń ilorazową i oblicz jej homologię bez korzystania z gotowej nazwy tej przestrzeni.
\item \textbf{Zawieszenie.} Dla spójnego CW-kompleksu $X$ oblicz $\widetilde H_k(\Sigma X)$ przez Mayera--Vietorisa. Zastosuj do $X=S^1\vee S^2$, podając wymiary wszystkich grup.
\wskazowka{Pokryj zawieszenie dwoma stożkami, których przecięcie deformacyjnie retraktuje się na $X$. Użyj zredukowanego ciągu Mayera--Vietorisa.}
\item \textbf{Para walca.} Dla $M=S^1\times[0,1]$ i $A=S^1\times\{0,1\}$ oblicz $H_*(M,A;\Z)$ i opisz homomorfizm łączący do $H_0(A)$.
\item \textbf{Homologia z identyfikacją brzegów.} Przestrzeń $X$ otrzymuje się z kwadratu przez sklejenie pionowych boków zgodnie z orientacją, a poziomych z odwróceniem orientacji. Zapisz komórkowy kompleks i oblicz $H_*(X;\Z/3\Z)$.
\item \textbf{Homologia lokalna w węźle grafu.} Niech $X$ będzie grafem z wierzchołkiem $v$ o walencji $5$. Oblicz $H_k(X,X\setminus\{v\};\Z)$ i porównaj z lokalną homologią punktu wewnętrznego krawędzi.
\wskazowka{Wytnij małe gwiazdowe otoczenie wierzchołka. Homologia lokalna w stopniu jeden odpowiada zredukowanej homologii jego łącza.}
\item \textbf{Stopień i stożek odwzorowania.} Niech $f:S^2\to S^2$ ma stopień $m\ne0$. Dołącz do $S^2$ jedną $3$-komórkę przez $f$. Oblicz wszystkie grupy homologii powstałej przestrzeni i sprawdź zależność od $m$.
\wskazowka{W komórkowym kompleksie różniczka od nowej $3$-komórki do $2$-komórki jest mnożeniem przez stopień $m$.}
\item \textbf{Przestrzeń z dwoma torusami.} Rozważ $X=(T^2\sqcup T^2)/(p\sim q)$, gdzie identyfikujesz jeden punkt w każdym torusie. Oblicz homologie $X$ i porównaj je z homologiami rozłącznej sumy dwóch torusów.
\wskazowka{Identyfikacja punktów daje sumę klinową torusów; skorzystaj ze zredukowanych homologii sumy klinowej.}
\end{enumerate}


\chapter{Wyższe grupy homotopii i związek z homologią}
\begin{enumerate}
\item \textbf{Wyższe grupy produktu.} Dla $n\ge2$ wykaż, że $\pi_n(X\times Y)\cong\pi_n(X)\oplus\pi_n(Y)$ dla przestrzeni bazowanych. Zastosuj do $S^2\times S^3$, wyznaczając $\pi_2$ oraz $\pi_3$; w drugim przypadku wolno przyjąć klasyczny wynik $\pi_3(S^2)\cong\Z$.
\item \textbf{Nakrycie i $\pi_2$.} Oblicz $\pi_2(\mathbb{RP}^3)$ za pomocą nakrycia $S^3\to\mathbb{RP}^3$. Porównaj wynik z $H_2(\mathbb{RP}^3;\Z)$.
\wskazowka{Nakrycie indukuje izomorfizmy grup homotopii w stopniach co najmniej dwa; oblicz je po przejściu do nakrycia uniwersalnego.}
\item \textbf{Przestrzeń Moore'a.} Dla $m\ge2$ rozważ $M_m=S^2\cup_m e^3$. Oblicz $\pi_2(M_m)$ za pomocą twierdzenia Hurewicza oraz $H_2(M_m;\Z)$ z kompleksu komórkowego.
\wskazowka{Przestrzeń Moore'a buduje się przez dołączenie komórki z odwzorowaniem określonego stopnia. Zapisz odpowiadający kompleks komórkowy i użyj twierdzenia Hurewicza dopiero po sprawdzeniu spójności.}
\item \textbf{Wykrywanie nietrywialnej wyższej grupy.} Niech $X=S^2\vee S^3$. Oblicz $\pi_2(X)$, $H_2(X)$ i wskaż klasę, którą wykrywa odwzorowanie Hurewicza. Uzasadnij, dlaczego samo $H_2$ nie wyznacza wszystkich wyższych grup homotopii.
\item \textbf{Przyczepienie komórki w stopniu wyższym.} Niech $X=S^3\cup_f e^4$ dla $f:S^3\to S^3$ o stopniu $7$. Oblicz $\pi_3(X)$ i $H_3(X;\Z)$, wyjaśniając, jak działa dołączenie komórki na generator sfery.
\wskazowka{Dołączenie komórki $n$-wymiarowej nie zmienia szkieletu niższego wymiaru; rozważ długi ciąg homotopii pary, zamiast tylko homologii.}
\item \textbf{Długi ciąg homotopii produktu z parą.} Dla pary $(S^2\times D^3,S^2\times S^2)$ oblicz $\pi_3$ względną, wykorzystując równoważność homotopijną i długi ciąg pary. Opisz obraz generatora w $\pi_2(S^2\times S^2)$.
\item \textbf{Stopień i $\pi_n$.} Niech $f:S^n\to S^n$, $n\ge2$, ma stopień $-3$. Oblicz homomorfizm $f_*:\pi_n(S^n)\to\pi_n(S^n)$ i rozstrzygnij, czy $f$ może być równoważnością homotopijną.
\item \textbf{Przyczepienie komórki a spójność.} Niech $X$ będzie jednospójny, a $Y=X\cup_f D^{n+1}$, $n\ge3$. Wyznacz zakres stopni $k$, w których $\pi_k(X)\to\pi_k(Y)$ musi być izomorfizmem, oraz stopień, w którym może nie być surjekcją lub injekcją. Uzasadnij w oparciu o względną aproksymację komórkową.
\wskazowka{Względne grupy homotopii pary $(X\cup_f D^n,X)$ są kontrolowane przez wymiar nowej komórki. Wykorzystaj wynik o spójności pary CW.}
\item \textbf{Przestrzeń o zerowych grupach niższego rzędu.} Rozważ CW-kompleks $X=S^4\vee S^4\vee S^5$. Oblicz $\pi_i(X)$ dla $i\le4$ oraz homomorfizm Hurewicza w stopniu $4$.
\item \textbf{Abelianizacja i odwzorowanie do okręgu.} Niech $X$ będzie powierzchnią orientowalną rodzaju $2$. Wyznacz $\pi_1(X)_{\mathrm{ab}}$ i opisz możliwe homomorfizmy $\pi_1(X)\to\Z$ jako obrazy indukowane przez odwzorowania $X\to S^1$.
\item \textbf{Porównanie dwóch sferycznych nakryć.} Oblicz $\pi_3(\mathbb{RP}^2)$ i $\pi_3(\mathbb{RP}^3)$, odwołując się do odpowiednich nakryć i przyjmując klasyczny wynik Hopfa $\pi_3(S^2)\cong\Z$. Wyjaśnij, dlaczego ostatniego wyniku nie można uzyskać przez bezpośrednie zastosowanie twierdzenia Hurewicza do $S^2$ w stopniu $3$.
\item \textbf{Działanie $\pi_1$ w przestrzeni rzutowej.} Dla parzystego $n\ge2$ określ działanie $\pi_1(\mathbb{RP}^n)$ na $\pi_n(\mathbb{RP}^n)\cong\Z$. Porównaj przypadek nieparzystego $n$ i uzasadnij rolę stopnia odwzorowania antypodycznego.
\wskazowka{Działanie grupy podstawowej na $\pi_2$ można odczytać z działania transformacji pokładowych nakrycia uniwersalnego.}
\item \textbf{Zawieszenie i pierwsza homotopia.} Dla $X=S^1\vee S^1$ oblicz grupy $\pi_2(\Sigma X)$ i $H_2(\Sigma X)$. Wskaż, w którym miejscu stosujesz jednospójność zawieszenia.
\item \textbf{Homologia nie wykrywa grupy podstawowej.} Porównaj $X=S^1\vee S^2$ i $Y=S^1\times S^2$ pod kątem $H_1,H_2,H_3$ oraz $\pi_1$. Znajdź niezmiennik wykrywający różnicę między ich typami homotopii.
\item \textbf{Homologia wymierna to za mało.} Niech $f:\T^2\to\T^2$ będzie odwzorowaniem indukowanym przez macierz $\operatorname{diag}(2,3)$. Oblicz $f_*$ na homologii całkowitej i wymiernej we wszystkich stopniach. Wyjaśnij, dlaczego izomorfizmy w homologii wymiernej nie pozwalają zastosować homologicznego kryterium Whiteheada.
\wskazowka{Homologia wymierna pomija torsję. Znajdź przykład komórkowy, w którym różniczka jest mnożeniem przez liczbę niezerową, lecz nieodwracalną nad $\Z$.}
\item \textbf{Włóknienie Hopfa jako narzędzie obliczeń.} Korzystając z długiego ciągu dokładnego homotopii włóknienia $S^1\hookrightarrow S^3\to S^2$, wyznacz $\pi_2(S^2)$ i $\pi_3(S^2)$ z uwzględnieniem grup homotopii sfer występujących w ciągu. Wskaż homomorfizm brzegowy i jego interpretację.
\wskazowka{Zapisz długi ciąg homotopii włóknienia $S^1\to S^3\to S^2$ w okolicy badanej grupy, pamiętając o znikaniu wyższych grup homotopii $S^1$.}
\end{enumerate}


\chapter{Formy różniczkowe, orientacja i kohomologia de Rhama}
\begin{enumerate}
\item \textbf{Zamkniętość formy z parametrem.} Dla $\omega=(ax^2y+y)\,dx+(x^3+bx)\,dy$ na $\R^2$ wyznacz parametry $a,b\in\R$, dla których $\omega$ jest zamknięta. Znajdź funkcję pierwotną w tych przypadkach.
\wskazowka{Wypisz $d\alpha$ i osobno zbierz współczynniki każdego niezależnego dwuwektora bazowego. Otrzymany warunek będzie równaniem na parametr.}
\item \textbf{Cofnięcie formy objętości.} Niech $F(u,v,w)=(u+v,uv,w+u^2)$. Oblicz $F^*(dx\wedge dy\wedge dz)$ i znajdź zbiór, w którym $F$ zachowuje orientację, odwraca ją lub nie jest lokalnym dyfeomorfizmem.
\item \textbf{Tożsamość Cartana w przykładzie.} Dla $X=y\partial_x-x\partial_y$ i $\omega=x\,dy-y\,dx$ na $\R^2$ oblicz $i_X\omega$, $i_X\dd\omega$ i $\mathcal L_X\omega$. Porównaj wynik z zachowaniem $\omega$ przez obroty.
\item \textbf{Forma na pierścieniu.} Dla $A=\{1<x^2+y^2<9\}$ i $\omega=((2x-y)\,dx+(x+2y)\,dy)/(x^2+y^2)$ rozłóż $\omega$ na część dokładną i wielokrotność formy kąta. Oblicz okres po dodatnio zorientowanym okręgu $r=2$.
\wskazowka{Porównaj całkę po okręgu otaczającym dziurę z całką po pętli ściągalnej. Nie wystarcza jedynie sprawdzenie $d\alpha=0$.}
\item \textbf{Dwie dziury, dwa okresy.} Na $U=\R^2\setminus\{(0,0),(2,0)\}$ zdefiniuj formy $\omega_1,\omega_2$ jako formy kąta wokół dwóch usuniętych punktów. Udowodnij liniową niezależność ich klas w $H^1_{\mathrm{dR}}(U)$, obliczając całki po małych okręgach.
\wskazowka{Wybierz bazę $H_1$ reprezentowaną przez małe dodatnio zorientowane okręgi wokół obu dziur. Okresy są współrzędnymi klasy de Rhama.}
\item \textbf{Stokes na pierścieniu.} Dla $\alpha=(x^3-y)\,dy$ oblicz $\int_{\partial A}\alpha$ oraz $\int_A\dd\alpha$, gdzie $A=\{1\le x^2+y^2\le4\}$. Ustal poprawną orientację wewnętrznej składowej brzegu.
\item \textbf{Całkowanie na walcu.} Rozważ $M=S^1\times[0,2]$ z orientacją $d\theta\wedge dt$ i formę $\omega=t\sin\theta\,d\theta+\cos\theta\,dt$. Oblicz $\dd\omega$ i obie strony twierdzenia Stokesa.
\item \textbf{Forma powierzchni na torusie.} Dla płaskiego torusa $\T^2=\R^2/\Z^2$ sprawdź, czy forma $\omega=(2+\sin(2\pi x))\,dx\wedge dy$ jest dokładna. Oblicz jej klasę w $H^2_{\mathrm{dR}}(\T^2)$ względem generatora $[dx\wedge dy]$.
\item \textbf{Okresy form w stopniu jeden.} Na $\T^2$ rozważ $\alpha=a\,dx+b\,dy+\dd f$, gdzie $f$ jest gładka i okresowa. Oblicz całki po dwóch podstawowych pętlach i rozstrzygnij, kiedy $\alpha$ jest dokładna.
\wskazowka{Sprowadź badanie dokładności do okresów po generatorach pierwszej homologii; do zbudowania potencjału użyj całki po drodze od punktu bazowego.}
\item \textbf{Kohomologia cylindra z usuniętą linią.} Oblicz $H^k_{\mathrm{dR}}((S^1\setminus\{p\})\times\R)$ i porównaj z kohomologią $S^1\times\R$. Zbuduj odpowiednie retrakcje deformacyjne.
\wskazowka{Najpierw zidentyfikuj składowe przestrzeni po usunięciu wskazanego podzbioru; kohomologię liczy się osobno na każdej składowej.}
\item \textbf{Forma magnetyczna na $\R^3$.} Dla $\omega=x\,dy\wedge dz+y\,dz\wedge dx+z\,dx\wedge dy$ oblicz $\dd\omega$ i znajdź $2$-formę $\eta$ o zadanej pochodnej $\dd\eta=\dd\omega$ i prostszych współczynnikach. Zinterpretuj niejednoznaczność wyboru $\eta$.
\item \textbf{Orientacja regularnej poziomicy.} Niech $M=\{(x,y,z):x^2+2y^2+3z^2=1\}$. Dla orientacji otaczającej $\R^3$ i normalnej zewnętrznej wyznacz dodatnią orientację $T_pM$ w punkcie $p=(1,0,0)$ i porównaj ją z orientacją parametryzacji lokalnej.
\wskazowka{Orientację regularnej poziomicy można opisać przez warunek, by dodatni wektor normalny i dodatnia baza przestrzeni stycznej dawały orientację otoczenia.}
\item \textbf{Zależność od orientacji.} Na torusie $\T^2$ rozważ $\omega=\cos(2\pi x)\,dx\wedge dy$ i $\eta=(1+\cos(2\pi x))\,dx\wedge dy$. Oblicz całki dla obu orientacji, a następnie rozstrzygnij, które z tych form mogą być dokładne.
\item \textbf{Kohomologia pary spodni.} Niech $P$ będzie powierzchnią powstałą ze sfery przez usunięcie wnętrz trzech rozłącznych domkniętych dysków. Korzystając z ciągu Mayera--Vietorisa, oblicz $H^k_{\mathrm{dR}}(P)$ i znajdź relację między okresami zamkniętej $1$-formy po trzech zorientowanych składowych brzegu.
\wskazowka{Para spodni retraktuje się na graf z dwiema niezależnymi pętlami. Uwzględnij relację między trzema krzywymi brzegowymi przy całkowaniu zamkniętych form.}
\item \textbf{Dwie osobliwości formy powierzchniowej.} Dla $p=(0,0,1)$ i $q=(0,0,-1)$ niech $U=\R^3\setminus\{p,q\}$ oraz $\omega_r=\bigl((x-r_1)\,dy\wedge dz+(y-r_2)\,dz\wedge dx+(z-r_3)\,dx\wedge dy\bigr)/\|x-r\|^3$ dla $r=p,q$. Oblicz okresy formy $2\omega_p-\omega_q$ po małych dodatnio zorientowanych sferach wokół $p$ i $q$. Wyznacz wymiar podprzestrzeni $H^2_{\mathrm{dR}}(U)$ rozpiętej przez jej klasy składowe.
\wskazowka{Oddziel lokalne okresy wokół poszczególnych osobliwości od okresu po dużej pętli obejmującej je obie; zastosuj addytywność klas homologii.}
\item \textbf{Stopień z całki formy objętości.} Na torusie $T^2=\R^2/\Z^2$ rozważ odwzorowanie $f([x,y])=([2x+y],[x+3y])$. Oblicz $f^*(\dd x\wedge\dd y)$ i wyznacz stopień $f$ przez porównanie całek formy objętości. Wyjaśnij, dlaczego wynik zgadza się z liczbą arkuszy nakrycia i orientacją.
\wskazowka{Do obliczenia stopnia użyj $\int_M f^*\omega=\deg(f)\int_N\omega$ dla znormalizowanej formy objętości; wyjaśnij wpływ orientacji.}
\end{enumerate}


\chapter{Tensor metryczny, długość i miara}
\begin{enumerate}
\item \textbf{Metryka ze stałymi współczynnikami.} Na $\R^2$ rozważ $g=5\,dx^2+4\,dx\,dy+2\,dy^2$, gdzie zapis mieszany oznacza $2g_{12}\,dx\,dy$. Oblicz $g^{-1}$, $v^\flat$ dla $v=(1,-2)$ oraz $\alpha^\sharp$ dla $\alpha=3dx-dy$.
\item \textbf{Metryka w nowych współrzędnych.} W dodatniej ćwiartce płaszczyzny zastosuj do metryki $g=dx^2+4dy^2$ zmianę $u=\log x$, $v=xy$. Oblicz współczynniki metryki w $(u,v)$ i element pola.
\wskazowka{Odwróć zmianę współrzędnych, wyrażając $x,y$ przez $u,v$, a następnie policz macierz Jacobiego. Nową metrykę daje $J^T GJ$.}
\item \textbf{Powierzchnia śrubowa.} Dla $F(u,v)=(u\cos v,u\sin v,av)$, $a>0$, wyznacz metrykę indukowaną, pole fragmentu $u\in[1,2]$, $v\in[0,2\pi]$, i sprawdź ortogonalność krzywych współrzędnych.
\item \textbf{Kula w metryce konforemnej.} Na $\R^2$ niech $g=(1+x^2+y^2)^{-2}(dx^2+dy^2)$. Oblicz długość okręgu euklidesowego o promieniu $R$ i pole dysku $r\le R$ względem $g$. Zbadaj granice przy $R\to\infty$.
\wskazowka{Przejdź na współrzędne biegunowe; długość i pole zawierają odpowiednio pierwszą i drugą potęgę konforemnego czynnika skali.}
\item \textbf{Rozbieżność długości.} Na $(0,1)$ przyjmij $g=x^{-2p}dx^2$, $p\in\R$. Wyznacz odległość między $a,b\in(0,1)$ i ustal, dla których $p$ granica $x\downarrow0$ znajduje się w nieskończonej odległości.
\wskazowka{W wymiarze jeden odległość jest całką z pierwiastka współczynnika metrycznego. Zbieżność przy zerze sprawdzaj testem całkowym.}
\item \textbf{Metryka produktowa.} Na $M=S^1_R\times(0,\infty)$ rozważ $g=R^2d\theta^2+dr^2$. Oblicz długość krzywej $\gamma(t)=(kt,at+b)$ dla ustalonego przedziału parametrów i porównaj ją z długościami rzutów na czynniki.
\item \textbf{Metryka wypaczona.} Niech $g=dr^2+(1+r^2)^2d\theta^2$ na $\R\times S^1$. Wyznacz formę objętości, długości równoleżników $r=r_0$ oraz pole obszaru $r\in[a,b]$.
\item \textbf{Energia a zmiana czasu.} Krzywa $\gamma(t)=(t,t^2)$, $t\in[0,1]$, leży w płaszczyźnie euklidesowej. Oblicz długość i energię $E=\frac12\int\|\dot\gamma\|^2dt$. Znajdź parametryzację o stałej prędkości i wyznacz jej energię.
\wskazowka{Parametr długości łuku $s(t)=\int_0^t\|\dot\gamma(u)\|du$ pozwala uzyskać stałą prędkość. Przy porównaniu energii użyj nierówności Cauchy'ego--Schwarza.}
\item \textbf{Gradient i izolinie.} Dla $g=4dx^2+dy^2$ oraz $f(x,y)=x^2+y^2$ oblicz $\operatorname{grad}_gf$, jego normę i punkty, w których gradient jest styczny do poziomic $f$.
\item \textbf{Strumień w metryce wypaczonej.} Na $M=(0,\infty)\times S^1$ z $g=dr^2+r^4d\theta^2$ oblicz dywergencję pola $X=r^{-2}\partial_r$ oraz jego strumień przez okręgi $r=a$ i $r=b$, z normalnymi zewnętrznymi dla pierścienia.
\wskazowka{Skorzystaj ze wzoru $\operatorname{div}X=|g|^{-1/2}\partial_i(|g|^{1/2}X^i)$. Orientacje normalnych na dwóch składowych brzegu są przeciwne.}
\item \textbf{Objętość elipsoidy.} W $\R^3$ z metryką euklidesową oblicz objętość obszaru $x^2/a^2+y^2/b^2+z^2/c^2\le1$, używając cofnięcia metryki przez liniowy dyfeomorfizm kuli. Porównaj z wyznacznikiem jego macierzy.
\item \textbf{Miara Hausdorffa na krzywej.} Dla wykresu $\gamma(t)=(t,t^2)$, $t\in[0,1]$, oblicz $\mathcal H^1(\gamma([0,1]))$ i sprawdź zgodność z długością z metryki indukowanej.
\item \textbf{Przekrój czasoprzestrzeni FLRW.} Dla $g=-dt^2+a(t)^2(dx^2+dy^2+dz^2)$ i $a(t)=\sqrt{1+t^2}$ oblicz objętość sześcianu współrzędnościowego w chwili $t$ oraz tempo jej zmiany względem czasu własnego obserwatora współporuszającego się.
\wskazowka{Na hiperpłaszczyźnie $t=\mathrm{const}$ wyznacz wyznacznik indukowanej metryki trójwymiarowej; dla obserwatora współporuszającego się $dt$ jest czasem własnym.}
\item \textbf{Zegary w metryce Schwarzschilda.} W obszarze $r>2M$ oblicz czas własny obserwatora o stałych $(r,\theta,\varphi)$ między zdarzeniami odległymi o $\Delta t$. Porównaj dwa takie zegary na różnych promieniach bez zakładania, że $t$ jest czasem własnym.
\wskazowka{Dla nieruchomego obserwatora $dr=d\theta=d\varphi=0$. Podstaw to do metryki, a dopiero potem porównaj $d\tau$ i $dt$.}
\item \textbf{Przyczynowość w metryce Minkowskiego.} Dla $g=-dt^2+dx^2+dy^2$ sklasyfikuj wektory $v_1=(2,1,1)$, $v_2=(1,1,0)$, $v_3=(1,2,0)$ i wyznacz warunki na parametry $(a,b)$, by wektor $(1,a,b)$ był czasowy, świetlny lub przestrzenny.
\item \textbf{Miara i kompletność metryczna.} Dla metryki $g=x^{-2}(\dd x^2+\dd y^2)$ na półpłaszczyźnie $x>0$ oblicz element powierzchni oraz długość pionowej krzywej zbliżającej się do $x=0$. Porównaj geometrię granicy z granicą w metryce euklidesowej.
\wskazowka{Porównaj dwie różne granice: $x\downarrow0$ oraz $x\to\infty$. Wzdłuż ustalonego $y$ długość jest całką $\int dx/x$.}
\end{enumerate}


\chapter{Koneksje i transport równoległy}
\begin{enumerate}
\item \textbf{Koneksja o zadanych współczynnikach.} Na $\R^2$ określ koneksję przez $\nabla_{\partial_x}\partial_x=0$, $\nabla_{\partial_x}\partial_y=\partial_x$, $\nabla_{\partial_y}\partial_x=-\partial_x$, $\nabla_{\partial_y}\partial_y=0$. Oblicz torsję i rozstrzygnij, czy jest to koneksja Leviego-Civity jakiejś metryki riemannowskiej.
\wskazowka{Najpierw policz $T(\partial_x,\partial_y)$. Koneksja Leviego-Civity zawsze ma torsję zerową, niezależnie od wyboru metryki.}
\item \textbf{Wzór Christoffela dla innej metryki.} Dla $g=dx^2+(1+x^2)^2dy^2$ na $\R^2$ oblicz wszystkie niezerowe symbole Christoffela i pochodne kowariantne pól $X=\partial_x+y\partial_y$, $Y=x\partial_x+\partial_y$.
\item \textbf{Koneksja płaska w krzywoliniowej mapie.} Dla zmiany współrzędnych $x=u$, $y=v+u^3$ oblicz symbole koneksji euklidesowej w mapie $(u,v)$. Sprawdź wybraną tożsamość transformacji symboli Christoffela.
\wskazowka{Płaska koneksja może mieć niezerowe symbole w nieliniowych współrzędnych. Wyprowadź je z równania $\nabla_{\partial_u}\partial_u$ w bazie $(\partial_u,\partial_v)$.}
\item \textbf{Transport w metryce diagonalnej.} Na $\R\times S^1$ niech $g=dx^2+e^{4x}d\theta^2$. Napisz układ równań transportu równoległego wzdłuż krzywej $\gamma(t)=(t,0)$ i rozwiąż go dla dowolnego wektora początkowego.
\wskazowka{Wzdłuż $\theta=0$ układ $\nabla_{\dot\gamma}V=0$ rozdziela się na równania dla składowych wektora. Użyj symboli dla metryki diagonalnej.}
\item \textbf{Transport radialny.} Na $(0,\infty)\times S^1$ rozważ $g=dr^2+r^4d\theta^2$. Znajdź wszystkie pola równoległe wzdłuż $\gamma(t)=(t,\theta_0)$, $t\in[a,b]$.
\item \textbf{Torsja i symetria symboli.} Niech na $\R^3$ współczynniki koneksji będą stałe i niezerowe wyłącznie dla $\Gamma^1_{23}=2$, $\Gamma^1_{32}=-1$. Oblicz wszystkie niezerowe wartości torsji na bazowych polach i sprawdź tensorowość otrzymanego wyrażenia na konkretnym przykładzie.
\item \textbf{Hesjan funkcji.} Dla metryki $g=dx^2+e^{2x}dy^2$ oblicz $\operatorname{Hess}_g f$ dla $f(x,y)=e^x+y^2$. Zbadaj jego symetrię i oblicz laplasjan $\Delta_g f=\operatorname{tr}_g\operatorname{Hess}_g f$.
\wskazowka{Hesjan ma składowe $\partial_i\partial_j f-\Gamma^k_{ij}\partial_k f$; laplasjan oblicz dopiero po kontrakcji z $g^{ij}$.}
\item \textbf{Dywergencja przez koneksję.} W metryce $g=dx^2+e^{2x}dy^2$ oblicz $\operatorname{div}_g X$ dla $X=e^{-x}\partial_x+xy\partial_y$ dwiema metodami: przez ślad $Y\mapsto\nabla_YX$ oraz przez gęstość objętości.
\item \textbf{Norma transportowanego wektora.} Dla $g=dx^2+e^{2x}dy^2$ oraz krzywej $\gamma(t)=(t,0)$ oblicz transport wektorów $V(0)=\partial_x+\partial_y$ i $W(0)=2\partial_x-\partial_y$. Sprawdź bezpośrednio zachowanie ich długości i iloczynu skalarnego oraz porównaj z koneksją z zadania 1.
\item \textbf{Koneksja na produktowej powierzchni.} Dla $g=du^2+4dv^2$ na $\T^2$ wyznacz symbole Christoffela w globalnych współrzędnych lokalnych, a następnie transport równoległy wzdłuż pętli $t\mapsto([t],[2t])$.
\item \textbf{Różniczka tensorowa.} Dla metryki $g=dx^2+(1+x^2)^2dy^2$ i tensora $T=dy\otimes\partial_x$ oblicz $\nabla_{\partial_x}T$ oraz $\nabla_{\partial_y}T$. Zwróć uwagę na działania koneksji na indeks górny i dolny.
\wskazowka{Dla tensora $dy\otimes\partial_x$ zastosuj Leibniza. Przy pochodnej kowektora pamiętaj o minusie przed symbolem Christoffela.}
\item \textbf{Dwie koneksje.} Niech $\nabla$ będzie koneksją Leviego-Civity metryki euklidesowej na $\R^2$, a $\widetilde\nabla_XY=\nabla_XY+\alpha(X)Y+\alpha(Y)X$ dla $\alpha=dx$. Wykaż, że $\widetilde\nabla$ jest bez torsji; sprawdź, czy zachowuje metrykę euklidesową.
\wskazowka{Różnica koneksji jest tensorem; jego część antysymetryczna odpowiada różnicy torsji. Zachowanie metryki sprawdzisz na polach bazowych.}
\item \textbf{Koneksja na powierzchni wykresowej.} Dla powierzchni $z=\tfrac12(x^2+y^2)$ w $\R^3$ oblicz metrykę indukowaną i symbole Christoffela w punkcie $(x,y)=(1,0)$. Porównaj z rzutem euklidesowej pochodnej na przestrzeń styczną.
\wskazowka{Wyznacz $g_{ij}=\langle\partial_iF,\partial_jF\rangle$ i jego pochodne w zadanym punkcie. Alternatywnie oblicz składową styczną $\partial_i\partial_jF$.}
\item \textbf{Transport po łamanej.} Na płaszczyźnie z koneksją określoną przez $\Gamma^1_{12}=\Gamma^1_{21}=c$ i wszystkie pozostałe symbole równe zero wyznacz transport wzdłuż odcinka poziomego, a potem pionowego. Zbadaj zależność od kolejności odcinków.
\wskazowka{Traktuj transport jako układ macierzowy $V'=-\Gamma(\dot\gamma)V$ na każdym odcinku osobno; sprawdź, czy uzyskane macierze transportu komutują.}
\item \textbf{Symbole dla Lorentzowskiej metryki diagonalnej.} Dla $g=-dt^2+(1+t^2)^2dx^2$ na $\R^2$ wyznacz symbole Christoffela, pochodną kowariantną pola $\partial_x$ w kierunku $\partial_t$ i równania transportu równoległego po linii $x=0$.
\wskazowka{Wzory na Christoffele działają także dla metryki Lorentzowskiej; istotny jest znak $g^{tt}$. Wzdłuż linii $x=0$ pozostaje układ dwóch ODE.}
\item \textbf{Różnica koneksji.} Dla dwóch koneksji $\nabla$ i $\widetilde\nabla$ na $TM$ rozważ $A(X,Y)=\widetilde\nabla_XY-\nabla_XY$. Sprawdź w konkretnych współrzędnych, że $A$ jest tensorem typu $(1,2)$, i określ zależność symetrii $A$ od torsji obu koneksji.
\end{enumerate}


\chapter{Geodezyjne, odwzorowanie wykładnicze i zupełność}
\begin{enumerate}
\item \textbf{Geodezyjne metryki wypaczonej.} Dla $g=dx^2+e^{2x}dy^2$ wyprowadź równania geodezyjnych i znajdź dwie całki pierwsze ruchu. Rozstrzygnij, które krzywe $y=\mathrm{const}$ są geodezyjnymi.
\wskazowka{Współrzędna $y$ jest cykliczna w lagranżjanie $\tfrac12(\dot x^2+e^{2x}\dot y^2)$. Drugą całką pierwszą jest energia geodezyjnej.}
\item \textbf{Równoleżnik jako geodezyjna.} Dla $g=dr^2+f(r)^2d\theta^2$ wyznacz warunek na $r_0$, aby krzywa $r=r_0$ z odpowiednią parametryzacją była geodezyjną. Zastosuj do $f(r)=1+r^2$ i $f(r)=2+\cos r$.
\wskazowka{W równaniu radialnym geodezyjnej pojawia się wyraz $-f(r)f'(r)\dot\theta^2$. Dla nietrywialnego równoleżnika sprawdź zerowanie odpowiedniego współczynnika.}
\item \textbf{Energia i parametr afiniczny.} Na $\R^2$ rozważ krzywą $\gamma(t)=(t^2,0)$, $t>0$. Rozstrzygnij, czy jest geodezyjną dla metryki euklidesowej, czy jest jej reparametryzacją, i znajdź parametr afiniczny.
\item \textbf{Geodezyjne w metryce stałej.} Niech $g=5dx^2+4dx\,dy+2dy^2$ na $\R^2$. Wyznacz $\exp_p(v)$ i odległość $d_g((0,0),(2,1))$. Scharakteryzuj wszystkie geodezyjne jednostkowej prędkości.
\item \textbf{Zupełność półprostej.} Dla $M=(0,\infty)$ z metryką $g=x^{-2p}dx^2$ ustal, dla jakich $p\in\R$ rozmaitość jest geodezyjnie zupełna. Wykaż odpowiedź, obliczając czasy dochodzenia do obu końców.
\wskazowka{Podstaw nową współrzędną długości $s(x)=\int^x t^{-p}\,dt$. Geodezyjna zupełność wymaga nieskończonej długości w obu końcach półprostej.}
\item \textbf{Przekrój hiperboliczny.} Dla półpłaszczyzny z metryką $g=y^{-2}(dx^2+dy^2)$ rozważ rodzinę półokręgów $x^2+y^2=R^2$, $y>0$. Wyznacz parametryzację jednostkowej prędkości wybranego łuku i jego końce idealne.
\wskazowka{Półokręgi prostopadłe do brzegu są geodezyjnymi. Dla $x=R\tanh s$, $y=R/\cosh s$ sprawdź normę pochodnej.}
\item \textbf{Krótkie łuki na cylindrze.} Na płaskim cylindrze $M=\R/(L\Z)\times\R$ oblicz odległość między $([0],0)$ i $([a],h)$. Wyznacz wszystkie geodezyjne minimalizujące, zwłaszcza gdy $a=L/2$.
\wskazowka{Podnieś punkty do płaszczyzny i minimalizuj długości odcinków do wszystkich przesunięć o $kL$. Wielokrotność minimów odpowiada remisowi.}
\item \textbf{Nieprostokątny torus płaski.} Niech $\Lambda$ będzie kratą generowaną przez $(1,0)$ i $(1/2,\sqrt3/2)$. Wyznacz odległość między $[0]$ a $[(1/2,0)]$ w $\R^2/\Lambda$ i liczbę geodezyjnych minimalizujących łączących te punkty.
\wskazowka{Najkrótsza geodezyjna na torusie to rzut odcinka do odpowiedniego punktu kraty. Porównaj długości wektorów $v-\lambda$, $\lambda\in\Lambda$.}
\item \textbf{Pochodne odwzorowania wykładniczego.} Na $\R^n$ z metryką stałą wyznacz $d(\exp_p)_v$ dla dowolnego $v$. Następnie porównaj z równaniem Jacobiego w zerowej krzywiźnie.
\item \textbf{Pierwsza wariacja na łamanej.} Dla dwóch odcinków w $\R^2$ łączących $p$ z $q$ przez punkt $r$ oblicz pochodną sumy długości przy przesunięciu $r$. Wykaż, że warunek stacjonarności przy $r\ne p,q$ wymusza współliniowość i odpowiednią orientację odcinków.
\wskazowka{Zróżnicz $|p-r|+|r-q|$ po współrzędnych $r$. Warunek krytyczny to równowaga dwóch jednostkowych wektorów.}
\item \textbf{Równania radialne w Schwarzschildzie.} Dla metryki $g=-F(r)dt^2+F(r)^{-1}dr^2+r^2d\Omega^2$, $F=1-2M/r$, rozważ geodezyjne radialne czasowe. Wyznacz zachowaną energię $E=F(r)\dot t$ i równanie dla $\dot r^2$ przy normalizacji $g(\dot\gamma,\dot\gamma)=-1$.
\wskazowka{Energia wynika z pola Killinga $\partial_t$. Po użyciu normowania wektora stycznego wyeliminuj $\dot t$ na rzecz $E$.}
\item \textbf{Geodezyjne zerowe w FLRW.} W płaskim modelu $g=-dt^2+a(t)^2dx^2$ wyprowadź równanie radialnej geodezyjnej świetlnej w postaci $\dd x/\dd t=\pm1/a(t)$. Oblicz zasięg współrzędnościowy sygnału od $t=1$ do $t=T$ dla $a(t)=t^p$.
\wskazowka{Wybierz współrzędne konforemne $d\eta=dt/a(t)$; promienie światła są prostymi w metryce konforemnej, ale parametry afiniczne wymagają osobnej uwagi.}
\item \textbf{Ekstremalność a minimalność.} Na okręgu długości $2\pi R$ znajdź geodezyjne od punktu $p$ do punktu antypody i rozstrzygnij, które łuki są minimalizujące. Wskaż geodezyjną łączącą te same punkty, która nie minimalizuje.
\item \textbf{Promień iniektywności ilorazu.} Niech $M=\R^2/\Lambda$ będzie płaskim torusem. Udowodnij, że jego promień iniektywności jest równy połowie długości najkrótszego niezerowego wektora kraty. Zastosuj do kraty z zadania 8.
\wskazowka{W ilorazie geodezyjne mogą być krótkie mimo płaskości lokalnej. Szukaj najkrótszego nietrywialnego przesunięcia kratowego lub elementu grupy pokładowej.}
\item \textbf{Dwie metryki na tej samej przestrzeni.} Porównaj zupełność $(\R,dx^2)$ i $(\R,e^{-2x^2}dx^2)$. Wyznacz długość promienia $[0,\infty)$ w obu metrykach i wskaż znaczenie twierdzenia Hopfa--Rinowa.
\item \textbf{Punkt sprzężony na sferze.} Na sferze jednostkowej rozważ geodezyjną wielkiego koła i normalne pole Jacobiego z $J(0)=0$, $\nabla_{\dot\gamma}J(0)\ne0$. Wyznacz pierwsze dodatnie miejsce zerowe $J$ i powiąż je z utratą lokalnej minimalności odcinka geodezyjnego.
\wskazowka{Pole Jacobiego normalne do geodezyjnej wielkiego koła spełnia równanie o stałej dodatniej krzywiźnie. Punkty sprzężone odpowiadają jego kolejnym zerom.}
\end{enumerate}


\chapter{Krzywizna: tensor Riemanna, porównania i Gauss--Bonnet}
\begin{enumerate}
\item \textbf{Nowa metryka wypaczona.} Dla $g=dr^2+(1+r^2)^2d\theta^2$ na $\R\times S^1$ oblicz krzywiznę Gaussa, określ jej znak i granice dla $|r|\to\infty$.
\wskazowka{Dla $g=dr^2+f(r)^2d\theta^2$ skorzystaj z $K=-f''/f$. Dopiero potem wyznacz $\Ric$ i skalar krzywizny w wymiarze dwa.}
\item \textbf{Krzywizna formy konforemnej.} Na $\R^2$ niech $g=e^{2u}(dx^2+dy^2)$, $u(x,y)=x^2-y^2$. Wyznacz $K_g$ i ustal, czy metryka jest płaska. Powtórz obliczenie dla $u(x,y)=x^2+y^2$.
\wskazowka{Przy zmianie $\tilde g=e^{2u}g$ na powierzchni użyj $\tilde K=e^{-2u}(K-\Delta_gu)$; sprawdź konwencję znaku laplasjanu ze skryptu.}
\item \textbf{Krzywizna rozmaitości produktowej.} Dla $M=S^2_R\times S^2_S$ z metryką produktową oblicz $\Ric$ i $\Scal$ oraz zakres możliwych krzywizn sekcyjnych w punkcie. Dla jakich $R,S$ metryka jest Einsteinowska?
\item \textbf{Wymiar trzy.} Niech na trójwymiarowej rozmaitości operator Ricciego ma wartości własne $(1,2,3)$ w bazie ortonormalnej. Oblicz krzywizny sekcyjne trzech płaszczyzn współrzędnych, korzystając z faktu, że w wymiarze trzy tensor Riemanna jest wyznaczony przez tensor Ricciego.
\wskazowka{W wymiarze trzy tensor Weyla znika. Wykorzystaj rozkład Riemanna przez $\Ric$ i $\Scal$, a nie próbuj odtwarzać wszystkich składowych niezależnie.}
\item \textbf{Płaskość nie wynika z Ricciego.} Na czterowymiarowej rozmaitości o metryce produktowej $S^2_R\times H^2_R$ oblicz tensor Ricciego i skalar krzywizny. Wyjaśnij, dlaczego znikanie skalara nie implikuje znikania tensora Riemanna.
\wskazowka{W produkcie $S^2_R\times H^2_R$ oba czynniki mają przeciwne krzywizny Gaussa. Oblicz osobno składowe $\Ric$ na obu czynnikach, a następnie ich ślad.}
\item \textbf{Transport wzdłuż małej pętli.} Dla metryki $g=dr^2+f(r)^2d\theta^2$, $f(r)=2+\cos r$, wyznacz przybliżony kąt holonomii wokół małego prostokąta współrzędnościowego o bokach $\varepsilon,\delta$ przy $r=r_0$. Uwzględnij element pola metryki.
\wskazowka{Holonomia małej pętli zaczyna się od wyrazu kwadratowego w długościach boków, proporcjonalnego do $R(X,Y)$. Dopilnuj konwencji znaku.}
\item \textbf{Równanie Jacobiego o zmiennym współczynniku.} Wzdłuż geodezyjnej na powierzchni o krzywiźnie $K(t)=2/(1+t)^2$ wyprowadź skalarne równanie pola Jacobiego poprzecznego. Znajdź jego rozwiązania w postaci potęg $(1+t)^\lambda$ i omów możliwe miejsca zerowe.
\wskazowka{Po podstawieniu $s=1+t$ szukaj rozwiązań $j=s^\lambda$; otrzymasz równanie kwadratowe na wykładnik.}
\item \textbf{Porównanie pól Jacobiego.} Dla ustalonego $\kappa>0$ porównaj rozwiązania $j''+Kj=0$, $j(0)=0$, $j'(0)=1$, w trzech przypadkach $K=-\kappa^2,0,\kappa^2$. Zbadaj zależność pierwszego zera od $\kappa$.
\wskazowka{Dla stałego $K$ równanie Jacobiego rozwiązują funkcje trygonometryczne, liniowe lub hiperboliczne; użyj tych samych warunków początkowych.}
\item \textbf{Oszacowanie średnicy.} Kompletna rozmaitość $M^4$ spełnia $\Ric\ge9g$. Jakie ograniczenie na średnicę daje twierdzenie Bonneta--Myersa? Zbuduj model sferyczny, dla którego oszacowanie jest ostre.
\wskazowka{Zapisz $\Ric\ge(n-1)kg$ i wyznacz $k$ przed zastosowaniem twierdzenia Bonneta--Myersa.}
\item \textbf{Objętość przy nieujemnym Riccim.} Niech $(M^3,g)$ będzie kompletna i $\Ric\ge0$. Z oszacowania Bishopa--Gromowa wyprowadź nierówność dla $\operatorname{Vol}(B_p(3r))/\operatorname{Vol}(B_p(r))$. Wyjaśnij, dlaczego samo twierdzenie nie daje dolnego ograniczenia tego ilorazu większego niż $1$.
\wskazowka{Funkcja $r\mapsto\mathrm{Vol}(B_p(r))/r^n$ jest nierosnąca. Podstaw dwa promienie i nie odwracaj kierunku nierówności.}
\item \textbf{Trójkąt na sferze.} Na $S_R^2$ rozważ trójkąt geodezyjny o kątach $\pi/2,\pi/2,\pi/3$. Oblicz jego pole i holonomię wokół brzegu. Wyjaśnij zgodność obu wartości z twierdzeniem Gaussa--Bonneta.
\item \textbf{Powierzchnia rodzaju $g$.} Niech $M_g$ będzie zamkniętą orientowalną powierzchnią rodzaju $g\ge2$, a $g_0$ metryką stałej krzywizny $-1$. Oblicz jej pole. Jak zmieni się całka krzywizny i pole po przeskalowaniu metryki przez $c^2$?
\item \textbf{Krzywizna przy brzegu.} Oblicz składnik brzegowy w twierdzeniu Gaussa--Bonneta dla obszaru $\{r\le R\}$ w metryce $g=dr^2+f(r)^2d\theta^2$, zakładając $f(0)=0$, $f'(0)=1$ i regularność w środku. Sprawdź wynik dla $f(r)=\sinh r$.
\wskazowka{Zastosuj $K=-f''/f$ i element pola $f(r)dr\,d\theta$; wkład brzegu zawiera pochodną $f'(R)$.}
\item \textbf{Osobliwość a zmiana współrzędnych.} Metryka $g=-d\tau^2+\tau^2d\chi^2$ dla $\tau>0$ jest płaska. Znajdź współrzędne Minkowskiego, oblicz $\Scal$ i porównaj powierzchnię $\tau=0$ z prawdziwą osobliwością krzywiznową modelu FLRW o $a(t)=t^{2/3}$.
\wskazowka{Sprawdź transformację $T=\tau\cosh\chi$, $X=\tau\sinh\chi$. Osobliwość współrzędnych odróżnisz od osobliwości krzywizny przez niezmienniki.}
\item \textbf{Tensor krzywizny przy nietypowym czynniku skali.} W płaskim przestrzennie modelu FLRW z $a(t)=t^{3/4}$, $t>0$, oblicz $\Scal$, $\Ric_{tt}$ oraz $R_{0i0j}$ w ortonormalnej ramie. Zbadaj granice przy $t\downarrow0$ i porównaj je z wynikiem dla $a(t)=t$.
\wskazowka{Najpierw oblicz $\dot a/a$ i $\ddot a/a$. We wzorach FLRW pracuj w ramie ortonormalnej, aby uniknąć pozornych potęg $a$.}
\item \textbf{Forma indeksowa.} W przestrzeni o stałej dodatniej krzywiźnie przekrojowej $K$ rozważ normalną wariację geodezyjnej jednostkowej o długości $L$ z ustalonymi końcami. Dla pola próbnego $V(t)=\sin(\pi t/L)E(t)$, gdzie $E$ jest równoległe i normalne, oblicz formę indeksową i określ warunek na $L$, przy którym nie jest dodatnia.
\wskazowka{Forma indeksowa dla pola normalnego to całka z $|D_tV|^2-K|V|^2$. Użyj całek z $\sin^2$ i $\cos^2$ na odcinku.}
\end{enumerate}


\chapter{Twierdzenie Sarda i transwersalność}
\begin{enumerate}
\item \textbf{Zbiór wartości krytycznych.} Dla $f:\R^2\to\R$, $f(x,y)=x^2y-y^3$, oblicz zbiór punktów krytycznych i wyznacz jego obraz. Sprawdź w tym przykładzie wniosek twierdzenia Sarda o mierze zbioru wartości krytycznych.
\item \textbf{Dwa rodzaje osobliwości.} Niech $F:\R^2\to\R^2$, $F(x,y)=(x^2-y^2,2xy)$. Wyznacz wartości krytyczne oraz liczbę punktów w typowym włóknie nad wartością regularną. Porównaj zachowanie przy początku i poza nim.
\item \textbf{Rodzina poziomic.} Dla $f:\R^3\to\R$, $f(x,y,z)=x^2+y^2-z^4$, wyznacz wartości krytyczne i opisz gładkość poziomic $f^{-1}(c)$ dla $c>0$, $c<0$, $c=0$.
\wskazowka{Zestaw zerowanie gradientu z równaniem poziomicy. Przy wartości krytycznej sprawdź stożkowy charakter osobliwości, zamiast automatycznie nazywać poziomicę rozmaitością.}
\item \textbf{Dwie powierzchnie.} Niech $S=\{z=x^2+y^2\}$ i $P_a=\{z=a\}$ w $\R^3$. Dla jakich $a$ przecięcie jest transwersalne? Wyznacz przestrzeń styczną i liczbę składowych przecięcia.
\item \textbf{Przecięcie okręgu i prostej.} Dla $S^1\subset\R^2$ i prostej $L_c=\{y=c\}$ rozstrzygnij transwersalność w punktach wspólnych w zależności od $c$. Wyjaśnij, co dzieje się przy przejściu przez wartość styczności.
\item \textbf{Transwersalność odwzorowania.} Dla $F:\R^2\to\R^2$, $F(x,y)=(x,y^2)$ i $Z=\{(u,v):u=v\}$ określ $F^{-1}(Z)$, zbadaj, czy $F\pitchfork Z$, i porównaj z przypadkiem $Z'=\{v=0\}$.
\wskazowka{Dla $F\pitchfork Z$ sprawdź, czy $dF(T_p\R^2)+T_{F(p)}Z=T_{F(p)}\R^2$ we wszystkich punktach przeciwobrazu.}
\item \textbf{Transwersalność i parametr.} Rozważ $F:\R\times\R\to\R$, $F(x,a)=x^3-3ax$. Wyznacz wszystkie parametry $a$, dla których funkcja $F_a$ ma niezerowy punkt krytyczny. Następnie znajdź parametry, dla których zero jest wartością regularną $F_a$.
\wskazowka{Oddziel warunek na punkt krytyczny $\partial_x F_a=0$ od warunku, że ten punkt należy do poziomicy $F_a^{-1}(0)$.}
\item \textbf{Perturbacja paraboli.} Niech $f_a(x,y)=y-x^2-a x$ i $g_b(x,y)=y-b$. Wskaż pary $(a,b)$, dla których poziomice zerowe $f_a,g_b$ nie przecinają się transwersalnie. Znajdź liczbę punktów przecięcia dla wszystkich pozostałych parametrów.
\wskazowka{Znajdź punkty styczności przez wspólne rozwiązanie równań poziomic i zależność liniową ich gradientów.}
\item \textbf{Przecięcie z brzegiem.} Rozważ dysk $D^2=\{x^2+y^2\le1\}$ oraz $f(x,y)=x$. Sprawdź, dla których $c$ poziomica $f^{-1}(c)$ jest transwersalna do $\partial D^2$. Porównaj jej topologię dla $|c|<1$, $|c|=1$ i $|c|>1$.
\item \textbf{Zera pola jako przecięcie z sekcją zerową.} Na $\R^2$ rozważ $X_a=(x^2-a)\partial_x+y\partial_y$. Wyznacz zera, sprawdź ich niezdegenerowanie jako zer przekroju $T\R^2$ i ustal, przy jakim $a$ transwersalność do sekcji zerowej zawodzi.
\wskazowka{Przekrój wiązki stycznej przecina sekcję zerową transwersalnie wtedy, gdy jego linearyzacja w zerze jest izomorfizmem.}
\item \textbf{Zera mapy do $\R^2$.} Niech $F_a(x,y)=(x^2+y^2-a,xy)$. Wyznacz wszystkie $a$, dla których $F_a$ jest transwersalna do $\{0\}$, i dla tych parametrów oblicz orientowane znaki lokalne zer.
\item \textbf{Indeksy punktów krytycznych.} Dla $f_a(x,y)=x^4-a x^2+y^2$ oblicz punkty krytyczne, ich hesjany i indeksy Morse'a. Wyznacz parametry, dla których $f_a$ jest funkcją Morse'a na $\R^2$.
\item \textbf{Transwersalność diagonalna.} Dla dwóch gładkich krzywych $\gamma,\eta:\R\to\R^2$ zdefiniuj $G(s,t)=(\gamma(s),\eta(t))$. Pokaż, że $G$ jest transwersalne do diagonali $\Delta\subset\R^2\times\R^2$ dokładnie wtedy, gdy krzywe przecinają się transwersalnie. Sprawdź to dla $\gamma(s)=(s,s^2)$, $\eta(t)=(t,a)$.
\wskazowka{Przestrzeń styczna do diagonali ma postać $\{(v,v)\}$. Zredukuj transwersalność $G$ do liniowej niezależności wektorów stycznych do krzywych.}
\item \textbf{Wymiar przestrzeni przecięć trajektorii.} Na $n$-wymiarowej rozmaitości dla funkcji Morse'a--Smale'a dane są punkty krytyczne $p,q$ o indeksach $k,\ell$. Oblicz oczekiwany wymiar $W^u(p)\cap W^s(q)$ i przestrzeni trajektorii po ilorazie przez przesunięcie czasu. Omów przypadki $k-\ell=0,1,2$.
\wskazowka{Przestrzenie $W^u(p)$ i $W^s(q)$ mają wymiary $k$ i $n-\ell$. Po poprzecznym przecięciu odejmij jeden wymiar działania translacji czasu.}
\item \textbf{Transwersalność wielu podrozmaitości.} Niech $S_1=\{x=0\}$, $S_2=\{y=0\}$ i $S_3=\{z=x+y\}$ w $\R^3$. Sprawdź transwersalność wszystkich przecięć parami i wspólnego przecięcia trzech powierzchni; sformułuj właściwy warunek na sumę przestrzeni konormalnych.
\wskazowka{Warunek dla trzech powierzchni sprawdzaj łącznie: ich konormalne muszą być liniowo niezależne. Transwersalność parami sama nie wystarcza.}
\item \textbf{Stopień odwzorowania z wartości regularnej.} Rozważ $f:S^1\to S^1$, $f(z)=z^m$, dla $m\in\Z\setminus\{0\}$. Oblicz stopień przez zliczanie znaków różniczek w przeciwobrazie wartości regularnej, a następnie określ zachowanie stopnia po złożeniu z odbiciem okręgu.
\end{enumerate}


\chapter{Zanurzenia, wiązki normalne i otoczenia tubularne}
\begin{enumerate}
\item \textbf{Wiązka normalna paraboloidy.} Dla $S=\{(x,y,x^2+2y^2)\}\subset\R^3$ znajdź lokalną jednostkową normalną, rzut ortogonalny na $T_pS$ i jawną postać odwzorowania tubularnego $(p,t)\mapsto p+t\nu(p)$ blisko $(0,0,0)$.
\wskazowka{Jednostkową normalną powierzchni wykresowej otrzymasz z iloczynu wektorowego pochodnych parametryzacji. Lokalną odwracalność mapy tubularnej sprawdź jej różniczką przy $t=0$.}
\item \textbf{Grubość otoczenia okręgu.} Dla okręgu $S^1_R\subset\R^2$ rozważ $F(\theta,t)=(R+t)(\cos\theta,\sin\theta)$. Wyznacz największy przedział symetryczny $(-\varepsilon,\varepsilon)$, dla którego $F$ jest dyfeomorfizmem na swój obraz.
\item \textbf{Wiązka normalna wykresu izometrii.} Niech $F:M\to M$ będzie izometrią rozmaitości riemannowskiej $(M,g)$, a $\Gamma_F=\{(p,F(p)):p\in M\}\subset M\times M$. Znajdź przestrzeń normalną $\nu_{(p,F(p))}\Gamma_F$ w metryce produktowej, przedstawiając jej elementy jako pary $(v,-dF_p v)$ z odpowiednim doborem bazowego wektora. Sprawdź wynik dla obrotu $F:S^1\to S^1$.
\wskazowka{Wektor styczny do wykresu to $(v,dF_pv)$. Normalność pary $(a,b)$ oznacza $g(a,v)+g(b,dF_pv)=0$ dla każdego $v$.}
\item \textbf{Krzywa w trójwymiarowej przestrzeni.} Niech $\gamma(t)=(\cos t,\sin t,\sin2t)$, $t\in\R/2\pi\Z$. Sprawdź, że jest zanurzeniem osadzonym, znajdź dwa liniowo niezależne pola normalne i skonstruuj lokalną trywializację jej wiązki normalnej.
\wskazowka{Zacznij od wektora stycznego $\gamma'(t)$, a następnie zastosuj procedurę Grama--Schmidta do dwóch dogodnych pól wzdłuż krzywej.}
\item \textbf{Brak właściwości odwzorowania.} Niech $i:(0,1)\hookrightarrow\R$ będzie inkluzją. Sprawdź, czy jest immersją, osadzeniem i odwzorowaniem właściwym. Wyjaśnij na tym przykładzie, dlaczego właściwość jest silniejszym warunkiem niż osadzenie.
\item \textbf{Właściwe zanurzenie wykresu.} Rozważ $F:\R\to\R^2$, $F(t)=(t,t^3-t)$. Sprawdź właściwość odwzorowania i wyznacz przestrzeń normalną w każdym punkcie. Skonstruuj jawny lokalny model otoczenia tubularnego w pobliżu $t=0$.
\item \textbf{Normalna do poziomicy.} Dla $M=\{x^2+2y^2+3z^2=1\}$ oblicz izomorfizm wiązki normalnej z trywialną wiązką liniową, korzystając z gradientu funkcji definiującej. Jak zmieni się wynik po usunięciu jednego punktu z $M$?
\item \textbf{Dwie framacje tego samego węzła.} Niech $K=S^1\times\{0\}\subset S^1\times\R^2$. Porównaj framacje normalne $\nu_0(\theta)=(e_1,e_2)$ i $\nu_k(\theta)$ otrzymaną przez obrót ramki o kąt $k\theta$. Wykaż, że ich klasy homotopii są różne dla różnych $k\in\Z$.
\wskazowka{Framacja określa mapę $S^1\to SO(2)$; jej stopień jest niezmiennikiem homotopijnym.}
\item \textbf{Jednostronna krzywa na butelce Kleina.} Przedstaw butelkę Kleina jako iloraz kwadratu z odpowiednimi identyfikacjami. Wskaż osadzony okrąg, którego wiązka normalna jest nietrywialna, i uzasadnij to, śledząc wektor normalny po jednym pełnym obiegu krzywej.
\wskazowka{Prześledź identyfikację brzegów fundamentalnego kwadratu i sprawdź, czy normalna wraca do siebie, czy do przeciwnego wektora.}
\item \textbf{Kołnierz na pierścieniu.} Dla $M=\{1\le x^2+y^2\le4\}$ skonstruuj jawne kołnierze obu składowych brzegu. Zwróć uwagę na różne kierunki wektora normalnego skierowanego do wnętrza.
\item \textbf{Podwojenie walca.} Zbuduj podwojenie $M=S^1\times[0,1]$ przez sklejenie dwóch kopii wzdłuż całego brzegu. Zidentyfikuj otrzymaną zamkniętą powierzchnię i wyznacz jej charakterystykę Eulera.
\item \textbf{Izotopia osadzeń.} Dla $F_t:S^1\to\R^2$, $F_t(\theta)=((2+t)\cos\theta,(1+t/2)\sin\theta)$, $t\in[0,1]$, sprawdź, że jest izotopią osadzeń. Skonstruuj izotopię otaczającą, działającą na całej $\R^2$.
\wskazowka{Rozważ radialną interpolację dodatnich funkcji skali, z funkcją odcięcia poza dużą kulą. Sprawdź globalną odwracalność na każdym etapie izotopii.}
\item \textbf{Funkcja właściwa.} Na $M=\R^2\setminus\{0\}$ rozważ $f(x)=\|x\|^2+\|x\|^{-2}$. Wykaż, że jest właściwa, znajdź wartości krytyczne i ustal, jak zmienia się typ topologiczny podpoziomic.
\item \textbf{Nieupełne i zupełne otoczenia.} Na $M=(0,1)$ rozważ metryki $g_1=dx^2$ oraz $g_2=x^{-2}(1-x)^{-2}dx^2$. Porównaj ich zupełność, a następnie wyjaśnij, jak wybór metryki wpływa na dziedziny geodezyjnego odwzorowania wykładniczego.
\wskazowka{W wymiarze jeden wystarczy całkować element długości wzdłuż przedziału. Zbieżność po obu końcach decyduje o zupełności.}
\item \textbf{Przygotowanie do chirurgii.} Dla osadzenia $S^1\times D^{n-1}\hookrightarrow M^n$ opisz topologię brzegu otoczenia tubularnego i liczbę potrzebnych niezależnych sekcji normalnych. Wyjaśnij, jakie dane zmienia obrót framacji, nawet gdy obraz osadzonej sfery pozostaje ten sam.
\wskazowka{Framacja jest trywializacją rzędu $n-1$ wiązki normalnej. Zmiana framacji to odwzorowanie sfery przyczepienia do $SO(n-1)$.}
\end{enumerate}


\chapter{Przecięcia, dualność Poincarégo i klasy charakterystyczne}
\begin{enumerate}
\item \textbf{Iloczyn kubkowy na produkcie.} Na $M=S^1\times S^2$ oznacz przez $a\in H^1(M;\Z)$ i $b\in H^2(M;\Z)$ klasy ściągnięte z czynników. Oblicz $a\smile b$, $a^2$, $b^2$ i opisz strukturę pierścienia $H^*(M;\Z)$.
\item \textbf{Pierścień kohomologii torusa trójwymiarowego.} Niech $x,y,z$ będą standardowymi generatorami $H^1(T^3;\Z)$. Wyznacz bazę $H^2$ i $H^3$ z ich iloczynów kubkowych. Oblicz $(x+2y)\smile(y+3z)\smile(z+x)$.
\wskazowka{Pierścień kohomologii $T^3$ jest algebrą zewnętrzną na trzech generatorach stopnia jeden. Przed mnożeniem uprość wyrażenie przez antyprzemienność.}
\item \textbf{Krzywe na torusie.} Dwie zorientowane krzywe proste na $T^2$ reprezentują klasy $(2,1)$ i $(1,3)$ w $H_1(T^2;\Z)$. Oblicz ich algebraiczną liczbę przecięcia i minimalną możliwą liczbę przecięć poprzecznych.
\item \textbf{Dualność na $S^1\times S^2$.} Wyznacz obrazy generatorów $H^1$ i $H^2$ przez dualność Poincarégo $H^k(M;\Z)\to H_{3-k}(M;\Z)$. Przedstaw je jako konkretne podrozmaitości.
\item \textbf{Forma przecięcia produktu.} Dla $M=S^2\times S^2$ wybierz klasy $A=[S^2\times\{p\}]$, $B=[\{q\}\times S^2]$. Oblicz macierz formy przecięcia na $H_2(M;\Z)$ i samoprzecięcia klas $A+B$ oraz $2A-3B$.
\item \textbf{Brak orientacji i homologia najwyższa.} Dla butelki Kleina oblicz $H_2(K;\Z)$ oraz $H_2(K;\Z/2\Z)$. Wyjaśnij, jakie podstawowe klasy fundamentalne istnieją i jak zmienia się sformułowanie dualności Poincarégo.
\item \textbf{Funkcja sklejająca wiązki.} Nad $S^2$ sklej dwie trywialne rzeczywiste wiązki rzędu dwa za pomocą $g_k:S^1\to SO(2)$, $g_k(e^{it})=R_{kt}$. Dla jakich $k$ otrzymujesz orientowalną wiązkę? Oblicz jej liczbę Eulera zgodnie z ustaloną orientacją.
\wskazowka{Wiązki orientowane płaszczyzn nad $S^2$ klasyfikuje stopień funkcji sklejającej $S^1\to SO(2)$. Liczba Eulera mierzy ten stopień z ustaloną konwencją znaku.}
\item \textbf{Klasa Eulera jako przeszkoda.} Dla zorientowanej wiązki rzędu dwa $E_k\to S^2$ z zadania 7 rozstrzygnij, kiedy istnieje przekrój nigdzie niezerowy. Porównaj odpowiedź z możliwością rozszczepienia $E_k\cong\underline\R\oplus L$.
\item \textbf{Samoprzecięcie przekroju zerowego.} Niech $E_k\to S^2$ będzie wiązką z zadania 7. Wyznacz liczbę samoprzecięcia sekcji zerowej we wnętrzu przestrzeni całkowitej $E_k$, przy przyjętej orientacji.
\item \textbf{Pierwsza klasa Stiefela--Whitneya.} Na $S^1\times S^1$ rozważ sumę $L_1\oplus L_2$, gdzie $L_i$ jest cofnięciem wiązki Möbiusa przez $i$-tą projekcję. Oblicz $w_1$ i $w_2$ tej wiązki oraz ustal jej orientowalność.
\wskazowka{Użyj wzoru Whitneya na $w_1,w_2$ sumy oraz faktu, że każda cofnięta wiązka Möbiusa ma niezerową klasę $w_1$.}
\item \textbf{Wiązka styczna do produktu sfer.} Korzystając ze wzorów na klasy charakterystyczne sum prostych wyznacz $w_2(T(S^2\times S^2))$. Rozstrzygnij, czy ta rozmaitość dopuszcza strukturę spinową.
\item \textbf{Klasy Cherna sumy prostych.} Niech $L,H$ będą zespolonymi wiązkami liniowymi nad $\mathbb{CP}^2$ z $c_1(L)=h$, $c_1(H)=2h$. Oblicz $c_1$ i $c_2$ wiązki $L\oplus H$ i sprawdź, czy może ona posiadać wszędzie niezerowy przekrój.
\wskazowka{Zastosuj $c(L\oplus H)=(1+c_1(L))(1+c_1(H))$. Niezerowa najwyższa klasa Cherna utrudnia istnienie globalnego nigdzie niezerowego przekroju.}
\item \textbf{Poincaré--Hopf z brzegiem.} Na domkniętym dysku w $\R^2$ rozważ pola $X_1=(x,y)$ i $X_2=(x,-y)$. Oblicz lokalne indeksy ich zer. Zbadaj zachowanie na brzegu i wyjaśnij, dla którego pola można bezpośrednio użyć wersji twierdzenia z polem skierowanym na zewnątrz.
\item \textbf{Indeks gradientu nietypowej funkcji na sferze.} Na standardowej $S^2\subset\R^3$ rozważ funkcję $f(x,y,z)=xy$ i pole $X=\operatorname{grad}_{S^2}f$. Wyznacz jego zera, oblicz indeks każdego z nich z hesjanu $f|_{S^2}$ i porównaj sumę z charakterystyką Eulera.
\wskazowka{Zera gradientu na sferze spełniają warunek, że gradient euklidesowy funkcji jest równoległy do promienia. Indeks wynika ze znaku wyznacznika hesjanu ograniczonego.}
\item \textbf{Forma przecięcia i sygnatura.} Niech $M$ będzie zamkniętą orientowalną $4$-rozmaitością, której forma przecięcia w pewnej bazie ma macierz $\begin{psmallmatrix}2&1\\1&-2\end{psmallmatrix}$. Oblicz wyznacznik, parzystość i sygnaturę formy. Rozstrzygnij, czy jest unimodularna i co mówi to o możliwości realizacji jej jako pełnej formy na wolnej części $H_2(M;\Z)$.
\wskazowka{Forma przecięcia zamkniętej orientowalnej $4$-rozmaitości jest unimodularna na części wolnej drugiej homologii. Sprawdź wyznacznik zanim interpretujesz macierz geometrycznie.}
\item \textbf{Parzystość i klasa Stiefela--Whitneya.} Niech $X$ będzie zamkniętą, zorientowaną, jednospójną $4$-rozmaitością, dla której $Q_X=\operatorname{diag}(1,-1)$. Korzystając ze wzoru Wu, rozstrzygnij, czy $X$ jest spinowa. Porównaj z przypadkiem $Q_X=H$.
\wskazowka{Zastosuj charakterystykę $w_2$ przez kongruencję $Q(x,x)\equiv\langle w_2,x\rangle\pmod2$. Porównaj parzystość obu form.}
\end{enumerate}


\chapter{Teoria Morse'a: punkty krytyczne, uchwyty i przepływ}
\begin{enumerate}
\item \textbf{Funkcja z trzema punktami krytycznymi.} Dla $f_a:\R\to\R$, $f_a(x)=x^4-2ax^2$, zbadaj punkty krytyczne, ich niezdegenerowanie i indeksy Morse'a w zależności od $a\in\R$. Naszkicuj zmiany topologii podpoziomic dla $a>0$.
\item \textbf{Dwa minima i siodło.} Dla $f:\R^2\to\R$, $f(x,y)=(x^2-1)^2+y^2$, oblicz punkty krytyczne, ich wartości i indeksy. Opisz typy homotopii podpoziomic przy przechodzeniu przez kolejne wartości krytyczne.
\item \textbf{Punkty krytyczne na produkcie.} Na $S^1\times S^2$ rozważ $f(\theta,p)=\cos\theta+2p_3$, gdzie $p_3$ to trzecia współrzędna punktu sfery. Wyznacz wszystkie punkty krytyczne, ich indeksy i liczby $m_k(f)$.
\item \textbf{Funkcja na przestrzeni rzutowej.} Niech $0<\lambda_0<\cdots<\lambda_n$ oraz $f:\mathbb{RP}^n\to\R$, $f([x])=(\sum_{j=0}^n\lambda_jx_j^2)/(\sum_{j=0}^nx_j^2)$. Znajdź punkty krytyczne i ich indeksy Morse'a. Zinterpretuj wynik w świetle homologii ze współczynnikami w $\Z/2\Z$.
\wskazowka{Punkty krytyczne ilorazu Rayleigha są prostymi własnymi diagonalnej macierzy. Indeks oblicz, porównując różnice wartości własnych w lokalnych współrzędnych rzutowych.}
\item \textbf{Gradient siodła z wyrazem mieszanym.} Dla $f(x,y)=2x^2-3y^2+2xy$ w metryce euklidesowej znajdź przepływ $-\operatorname{grad}f$ przez diagonalizację hesjanu. Wyznacz $W^s(0)$ i $W^u(0)$, ich równania i wymiary.
\wskazowka{Przepływ antygradientowy jest liniowym układem $\dot z=-Hz$. Rozkład na kierunki własne dodatnich i ujemnych wartości własnych $H$ daje rozmaitości stabilną i niestabilną.}
\item \textbf{Zmiana metryki i podprzestrzeni stabilnych.} Dla $f(x,y)=2x^2-3y^2+2xy$ i metryki $g=4dx^2+dy^2$ wyznacz operator liniowy pola $-\operatorname{grad}_gf$ i jego podprzestrzenie stabilną oraz niestabilną. Porównaj ich kierunki z wynikami zadania 5.
\wskazowka{Gradient w metryce $G$ ma macierz $G^{-1}H$, nie sam $H$. Choć operator może nie być symetryczny euklidesowo, jego wartości własne są rzeczywiste.}
\item \textbf{Podpoziomice na cylindrze.} Na $S^1\times\R$ rozważ $f(\theta,t)=t^2+\cos\theta$. Oblicz punkty krytyczne i indeksy, a następnie opisz topologię $\{f\le c\}$ dla regularnych $c$ w kolejnych przedziałach.
\wskazowka{Topologię podpoziomic można śledzić przez uchwyty przy kolejnych wartościach krytycznych. Oddziel pojawienie się nowej składowej od dołączenia uchwytu indeksu jeden.}
\item \textbf{Charakterystyka Eulera z uchwytów.} Rozmaitość zwarta powstaje przez dołączenie $m_k$ uchwytów indeksu $k$, gdzie $(m_0,m_1,m_2,m_3)=(1,4,5,2)$. Oblicz jej charakterystykę Eulera. Wyznacz ograniczenia na liczby Bettiego, gdy rozmaitość jest zamknięta i orientowalna.
\item \textbf{Nierówności Morse'a w przykładzie.} Zamknięta orientowalna powierzchnia rodzaju $3$ ma funkcję Morse'a z jednym minimum i jednym maksimum. Wyznacz dolne ograniczenie na liczbę siodeł. Czy osiągnięcie tego ograniczenia jest zgodne z mocnymi nierównościami Morse'a?
\item \textbf{Funkcja doskonała na $\mathbb{CP}^2$.} Znając grupy homologii $\mathbb{CP}^2$, podaj najmniejszą liczbę punktów krytycznych, jaką może mieć funkcja Morse'a. Skonstruuj kandydatkę za pomocą funkcji $[z_0:z_1:z_2]\mapsto (|z_1|^2+2|z_2|^2)/(\sum|z_j|^2)$ i sprawdź jej punkty krytyczne.
\wskazowka{Zwróć uwagę, że funkcja jest dobrze określona na klasach rzutowych. W mapach wokół osi współrzędnych rozwiń ją do wyrazów kwadratowych.}
\item \textbf{Dwa możliwe dołączenia $1$-uchwytu.} Rozważ sumę dwóch dysków $D^2\sqcup D^2$. Porównaj dołączenie $1$-uchwytu łączącego dyski z dołączeniem $1$-uchwytu, którego oba końce leżą na brzegu tego samego dysku. Oblicz liczbę składowych i charakterystykę Eulera w obu przypadkach.
\item \textbf{Rachunek kompleksu Morse'a.} Załóż, że kompleks Morse'a nad $\Z$ ma $C_2=\Z$, $C_1=\Z^2$, $C_0=\Z$ oraz $\partial_2(1)=(2,1)$, $\partial_1=0$. Oblicz homologie. Sprawdź, czy taki kompleks mógłby być kompleksem Morse'a zamkniętej orientowalnej powierzchni i uzasadnij.
\wskazowka{Kompleks Morse'a musi odtwarzać homologie rozmaitości. Dla zamkniętej powierzchni porównaj szczególnie $H_0$ i najwyższą grupę homologii.}
\item \textbf{Zliczanie trajektorii.} Przyjmij, że pomiędzy punktami krytycznymi $p$ indeksu $2$ i $q_1,q_2$ indeksu $1$ istnieją odpowiednio trzy i dwa izolowane przepływy gradientowe, o znakach $(+,+,-)$ oraz $(+,-)$. Oblicz współczynniki w $\partial p$ i omów wpływ odwrócenia orientacji rozmaitości niestabilnej $W^u(p)$.
\item \textbf{Trajektorie złamane.} Dla punktów krytycznych $p,q,r$ o indeksach odpowiednio $3,2,1$ wyznacz oczekiwane wymiary przestrzeni trajektorii $\mathcal M(p,q)$, $\mathcal M(q,r)$ oraz $\mathcal M(p,r)$. Wyjaśnij, dlaczego pary trajektorii przez $q$ pojawiają się na brzegu jednowymiarowej kompaktifikacji $\mathcal M(p,r)$.
\wskazowka{Przestrzenie trajektorii niezparametryzowanych mają wymiar równy różnicy indeksów minus jeden. Złamanie trajektorii pojawia się jako zdegenerowana granica przy wyjściu z przestrzeni moduli.}
\item \textbf{Wymiar i konieczność uchwytów.} Niech zamknięta, spójna, orientowalna $3$-rozmaitość ma $b_1=2$. Wyprowadź z nierówności Morse'a ograniczenia na liczby krytyczne indeksów $0,1,2,3$, a następnie wyznacz najmniejszą możliwą łączną liczbę punktów krytycznych, którą dopuszczają te ograniczenia.
\wskazowka{Połącz $b_0=b_3=1$, dualność Poincarégo $b_1=b_2$ oraz mocne nierówności Morse'a. Warunki algebraiczne nie gwarantują jeszcze realizacji funkcją Morse'a.}
\item \textbf{Niezmienność homologii przy narodzinach pary.} Do danego kompleksu Morse'a dołącz parę generatorów $a\in C_{k+1}$, $b\in C_k$ z $\partial a=b$, $\partial b=0$. Oblicz zmianę homologii i charakterystyki Eulera. Zinterpretuj ją w języku narodzin znoszącej się pary uchwytów.
\end{enumerate}


\chapter{Twierdzenie Smale'a o \texorpdfstring{$h$}{h}-kobordyzmie}
\begin{enumerate}
\item \textbf{Kobordyzm z konkretnym brzegiem.} Niech $(W;N_-,N_+)$ będzie $h$-kobordyzmem i załóż, że $N_-=S^2\times S^3$. Oblicz wszystkie grupy homologii $W$ i $N_+$ oraz charakterystyki Eulera tych trzech rozmaitości. Rozstrzygnij, czy $N_+$ mógłby być sferą $S^5$.
\item \textbf{Wykrywanie kobordyzmu, który nie jest typu $h$.} Niech $W$ będzie zwartym kobordyzmem o $\chi(W)=3$, a jego brzeg wejściowy $N_-$ ma charakterystykę Eulera równą $2$. Wyjaśnij, dlaczego inkluzja $N_-\hookrightarrow W$ nie może być równoważnością homotopijną. Podaj przykład sytuacji, w której równość charakterystyk nie daje takiej równoważności.
\item \textbf{Ograniczenie na możliwy brzeg.} Czy istnieje $h$-kobordyzm $(W;N_-,N_+)$, w którym $N_-$ jest $5$-wymiarową przestrzenią soczewkową o grupie podstawowej $\Z/5\Z$, a $N_+$ jest po prostu spójna? Uzasadnij odpowiedź. Następnie wyjaśnij, dlaczego zgodność samej grupy podstawowej nie wystarcza do sprawdzenia warunku typu $h$.
\item \textbf{Standardowy produkt.} Dla $N=S^2\times S^3$ i $W=N\times[0,1]$ zdefiniuj funkcję Morse'a bez punktów krytycznych, wyznacz względny kompleks uchwytowy i opisz dyfeomorfizmy obu brzegów z $N$.
\item \textbf{Macierz z jednostkowym wyznacznikiem.} Kompleks uchwytowy ma $C_3=C_2=\Z^2$ oraz $\partial_3=\begin{psmallmatrix}2&1\\1&1\end{psmallmatrix}$. Oblicz jego homologię, znajdź operacje elementarne sprowadzające macierz do jednostkowej i zinterpretuj je jako zmiany baz kompleksu.
\wskazowka{Sprowadzanie macierzy do jednostkowej wykonaj przez całkowitoliczbowe operacje elementarne. Interpretacja geometryczna wymaga jeszcze kontrolowania przecięć sfer.}
\item \textbf{Macierz z torsją homologiczną.} Dla $\partial_3=\begin{psmallmatrix}2&0\\0&3\end{psmallmatrix}:\Z^2\to\Z^2$ oblicz homologię względną. Wyjaśnij, dlaczego przy takim jedynym niezerowym operatorze brzegowym nie można otrzymać $h$-kobordyzmu.
\item \textbf{Pojedyncza para uchwytów.} Załóż, że względna funkcja Morse'a na kobordyzmie ma dokładnie jeden uchwyt indeksu $k$ i jeden indeksu $k+1$, a liczba przecięcia sfer przyczepienia i pasa wynosi $m\in\Z$. Wyznacz homologię względną; kiedy dane algebraiczne pozwalają oczekiwać zniesienia pary?
\wskazowka{Kompleks względny ma tylko różniczkę $\Z\xrightarrow{\cdot m}\Z$. Rozróżnij algebraiczną acykliczność od geometrycznych warunków zniesienia uchwytów.}
\item \textbf{Znak liczby przecięcia.} Dla dwóch zorientowanych sfer w poziomicy o wymiarach dopełniających się do wymiaru poziomicy oblicz, jak zmienia się algebraiczna liczba przecięcia po odwróceniu orientacji jednej ze sfer i po zamianie kolejności czynników.
\item \textbf{Przecięcia geometryczne a algebraiczne.} Dwie sfery spotykają się poprzecznie w pięciu punktach o znakach $(+,+,-,+,-)$. Oblicz algebraiczną liczbę przecięcia. Jakie dodatkowe warunki geometryczne są potrzebne, by zredukować liczbę punktów przecięcia do jednego?
\wskazowka{Znaki dają algebraiczną liczbę przecięć, lecz do triku Whitneya potrzebujesz par o znakach przeciwnych i odpowiedniej kontroli dysków oraz grupy podstawowej.}
\item \textbf{Działanie przesunięcia uchwytu.} Niech macierz liczb przecięcia ma postać $A=\begin{psmallmatrix}1&4\\0&1\end{psmallmatrix}$. Zapisz przesunięcie uchwytu odpowiadające operacji usunięcia elementu $4$ i wyznacz powstałą macierz. Czy sama redukcja algebraiczna automatycznie oznacza geometryczne zniesienie uchwytów?
\item \textbf{Wymiar triku Whitneya.} Dla dwóch podrozmaitości w $(m-1)$-wymiarowej poziomicy o wymiarach $k-1$ i $m-k$ sprawdź wymiar ich przecięcia transwersalnego. Następnie wyjaśnij, skąd biorą się przeszkody do zastosowania dysku Whitneya przy zbyt małym $m$.
\wskazowka{Wzór na wymiar przecięcia to suma wymiarów minus wymiar otoczenia. Dla dysku Whitneya istotna jest dodatkowo możliwość usunięcia jego wnętrza z innych podrozmaitości.}
\item \textbf{Liczba uchwytów a acykliczność.} Załóż, że względny rozkład uchwytowy ma tylko uchwyty indeksów $2$ i $3$, a ich liczby wynoszą odpowiednio $r$ i $s$. Wyprowadź konieczny warunek na $r,s$, by kompleks był acykliczny nad $\Q$. Czy warunek ten jest dostateczny nad $\Z$?
\wskazowka{Acykliczność nad $\Q$ oznacza, że kwadratowa macierz różniczki musi być odwracalna nad $\Q$. Nad $\Z$ wyznacznik powinien być jednostką.}
\item \textbf{Macierz i geometria.} Dla $A=\begin{psmallmatrix}1&2&0\\0&1&3\\0&0&1\end{psmallmatrix}$ wyznacz sekwencję elementarnych zmian baz sprowadzającą $A$ do $I_3$. Wymień odpowiadające im typy przesunięć uchwytów, nie powtarzając ogólnego dowodu twierdzenia Smale'a.
\item \textbf{Założenie prostej spójności.} Wyjaśnij, które etapy zastosowania triku Whitneya mogą zawodzić w obecności nietrywialnej $\pi_1(W)$. Skonstruuj dwie pętle zbudowane z dróg biegnących do przecięć, które mają różne klasy w grupie podstawowej.
\wskazowka{Etykiety przecięć zapisuj jako klasy pętli złożonych z dróg bazowych i fragmentów obu sfer; przeciwne znaki nie wystarczają, gdy klasy w $\pi_1$ się różnią.}
\item \textbf{Granice twierdzenia.} Rozważ sytuacje: $m=5$ z $\pi_1(W)=0$ oraz $m=7$ z niezerową $\pi_1(W)$. Dla każdej wskaż, którego założenia klasycznego twierdzenia Smale'a brakuje i dlaczego nie wolno zacytować jego konkluzji bez dalszych argumentów.
\end{enumerate}


\chapter{Torsja Whiteheada i twierdzenie o \texorpdfstring{$s$}{s}-kobordyzmie}
\begin{enumerate}
\item \textbf{Mnożenie w pierścieniu grupowym.} Dla $\pi=C_3=\langle t\mid t^3=1\rangle$ oblicz $(1+t-t^2)(2-t)$ w $\Z[\pi]$. Porównaj wynik z obliczeniem w $\Z[t]/(t^3-1)$.
\item \textbf{Pierścień grupowy grupy nieprzemiennej.} Niech $\pi=S_3$ oraz $a=(12)$, $b=(123)$. Oblicz $ab-ba$ w $\Z[S_3]$ i wykaż, że nie jest zerem. Wyjaśnij, dlaczego w macierzach nad $\Z[S_3]$ nie można swobodnie zamieniać kolejności czynników.
\item \textbf{Jednostki trywialne.} Dla dowolnego $g\in\pi$ sprawdź, że $\pm g$ jest jednostką w $\Z[\pi]$. Zbadaj, jakie klasy wyznaczają macierze diagonalne $\operatorname{diag}(g,g^{-1})$ i $\operatorname{diag}(-g,1)$ po przejściu do $\Wh(\pi)$.
\item \textbf{Macierze elementarne.} Niech $R$ będzie dowolnym pierścieniem z jedynką, a $a,b\in R$. Oblicz $E_{12}(a)E_{12}(b)$ i $E_{12}(a)^{-1}$. Zbadaj, co zmienia się, gdy $a,b$ nie komutują.
\item \textbf{Acykliczny kompleks nad pierścieniem grupowym.} Niech $R=\Z[t,t^{-1}]$, $C_3=C_2=R$, $\partial_3$ będzie mnożeniem przez $-t^4$. Oblicz homologię i klasę torsji w $\Wh(\Z)$, przyjmując standardowe bazy.
\item \textbf{Macierz z elementami Laurentowskimi.} Dla $R=\Z[t,t^{-1}]$ rozważ $A=\begin{psmallmatrix}t^2&1+t\\0&t^{-2}\end{psmallmatrix}$. Znajdź $A^{-1}$, wyznacz rozkład na macierze elementarne i diagonalne oraz oblicz klasę $[A]\in\Wh(\Z)$.
\wskazowka{Macierz górnotrójkątną rozdziel na część diagonalną i macierz elementarną. W $\Wh(\Z)$ klasy potęg $t$ są trywialne.}
\item \textbf{Odwracalność nie wystarcza dla torsji.} Wyjaśnij na poziomie definicji różnicę między warunkiem $A\in GL_r(\Z[\pi])$ i warunkiem $[A]=0\in\Wh(\pi)$. Wskaż, dlaczego odwracalność gwarantuje acykliczność dwuwyrazowego kompleksu, ale nie gwarantuje możliwości jego uproszczenia przez dopuszczalne ruchy.
\wskazowka{Odwracalność różniczki mówi o homologii kompleksu, zaś torsja pyta o jej klasę w ilorazie $K_1$. Rozdziel te dwa poziomy informacji.}
\item \textbf{Przykład kompleksu nieacyklicznego.} Nad $R=\Z[t,t^{-1}]$ weź $C_3=C_2=R$ i $\partial_3$ dane przez mnożenie przez $t-1$. Oblicz $H_2$ i $H_3$, uzasadnij, że $t-1$ nie jest jednostką, i wyjaśnij, dlaczego torsja acyklicznego kompleksu nie jest tu zdefiniowana w omawianym sensie.
\item \textbf{Jawna stabilizacja macierzy.} Dla $R=\Z[t,t^{-1}]$ i $A=\begin{psmallmatrix}1&t^2\\0&1\end{psmallmatrix}$ zapisz $A\oplus I_2$ i sprowadź tę macierz do $I_4$ przez operacje elementarne. Zinterpretuj dodatkowy blok jako algebraiczny obraz dwóch znoszących się par uchwytów.
\item \textbf{Zmiana bazy w dwuwyrazowym kompleksie.} Niech $C_3=C_2=R^2$ i $\partial_3=A\in GL_2(R)$. Wyznacz nową macierz operatora po zastąpieniu baz przez macierze $P$ i $Q$ w stopniach odpowiednio $3$ i $2$. Oblicz, jak zmienia się reprezentant torsji.
\wskazowka{Przy bazach zadanych macierzami $P,Q$ nową macierz różniczki otrzymasz przez przemnożenie z lewej i prawej przez odpowiednie odwrotności; pilnuj konwencji wektorów kolumnowych.}
\item \textbf{Torsja elementarnej pary.} Rozważ acykliczny kompleks $0\to R\xrightarrow{\id}R\to0$ z bazami standardowymi, skoncentrowany w stopniach $3$ i $2$. Oblicz jego torsję i porównaj z kompleksem otrzymanym po zmianie bazy w stopniu $3$ przez pomnożenie przez $g\in\pi$.
\item \textbf{Prosta homotopia CW.} Rozważ kompleks CW otrzymany z odcinka przez dołączenie pary komórek tworzących wolną ścianę: nowej $1$-komórki i nowej $2$-komórki. Zbuduj jawne zwinięcie elementarne i wyjaśnij, dlaczego takie rozszerzenie nie zmienia prostej klasy homotopii.
\item \textbf{Zastosowanie $s$-kobordyzmu do macierzy.} Załóż, że $m\ge6$, a względny kompleks uchwytowy pewnego $h$-kobordyzmu $W$ nad $R=\Z[\pi]$ ma jedyny niezerowy operator $\partial_3=\begin{psmallmatrix}1&a\\0&1\end{psmallmatrix}$, gdzie $a\in R$. Oblicz klasę torsji i wyprowadź wniosek o typie dyfeomorficznym $W$ względem jego brzegu wejściowego.
\wskazowka{Zauważ, że macierz różniczki jest elementarna. Dopiero po obliczeniu jej klasy w $\Wh(\pi)$ przywołaj twierdzenie o $s$-kobordyzmie.}
\item \textbf{Nietrywialny pierścień i zerowa torsja.} Niech $R=\Z[\pi]$ dla dowolnej grupy $\pi$ i $A=\begin{psmallmatrix}1&a&0\\0&1&b\\0&0&1\end{psmallmatrix}$ dla $a,b\in R$. Sprowadź $A$ do $I_3$ za pomocą macierzy elementarnych, pamiętając o kolejności mnożenia. Wywnioskuj wartość klasy $[A]\in\Wh(\pi)$.
\wskazowka{Usuwaj wyrazy nad przekątną w odpowiedniej kolejności; nad nieprzemiennym pierścieniem nie można bezkarnie zamieniać $ab$ na $ba$.}
\item \textbf{Działanie inwolucji.} Dla $R=\Z[\pi]$ określ $\overline{\sum n_g g}=\sum n_g g^{-1}$. Sprawdź $\overline{xy}=\bar y\,\bar x$. Oblicz inwolucję dla dwóch jawnych elementów w $\Z[S_3]$ i wyjaśnij, dlaczego nie jest ona na ogół homomorfizmem pierścieniowym, lecz antyhomomorfizmem.
\wskazowka{Inwolucja grupowa odwraca kolejność czynników: $(gh)^{-1}=h^{-1}g^{-1}$. Rozszerz tę własność liniowo na pierścień grupowy.}
\end{enumerate}


% Rozdzial 21 w dostarczonych plikach stanowi kolejna wersje materialu
% o torsji Whiteheada; ponizsze zadania rozszerzaja rozdzial 20,
% zamiast powielac jego podstawowe cwiczenia.

\chapter{Torsja Whiteheada i \texorpdfstring{$s$}{s}-kobordyzm --- zadania uzupełniające}
\begin{enumerate}
\item \textbf{Augmentacja nie wykrywa torsji.} Niech $R=\Z[C_5]$, $C_5=\langle t\mid t^5=1\rangle$, i $\varepsilon(\sum a_jt^j)=\sum a_j$. Dla jednostki $u=1-t-t^{-1}$ znanej ze skryptu oblicz $\varepsilon(u^k)$ dla $k\in\Z$. Oblicz obrazy klas $[u^k]$ przez homomorfizm $\Wh(C_5)\to\Wh(\{e\})$ indukowany augmentacją i wyjaśnij, dlaczego nie pozwala on rozróżnić tych parami różnych klas.
\wskazowka{Augmentacja przenosi jednostki pierścienia grupowego do jednostek $\Z$. Porównaj wynik z definicją $\Wh(\{e\})$.}
\item \textbf{Dwie reprezentacje pierścienia grupowego.} W $\Z[C_5]$ rozważ homomorfizmy $t\mapsto1$ oraz $t\mapsto\zeta$, gdzie $\zeta=e^{2\pi i/5}$. Dla $a=2-t-t^{-1}$ oblicz oba obrazy. Wyjaśnij, dlaczego nierówność $|a(\zeta)|\ne1$ sama nie dowodzi, że $a$ jest jednostką.
\item \textbf{Macierze nad nieprzemiennym pierścieniem.} W $R=\Z[S_3]$ przyjmij $a=(12)$ i $b=(123)$. Oblicz iloczyny $E_{12}(a)E_{23}(b)$ i $E_{23}(b)E_{12}(a)$ w $GL_3(R)$. Znajdź ich komutator i porównaj go z analogicznym rachunkiem dla $E_{23}(b)E_{12}(a)$ z zamienionymi indeksami.
\wskazowka{Zapisz $E_{ij}(a)=I+ae_{ij}$ i wykorzystaj regułę $e_{ij}e_{kl}=\delta_{jk}e_{il}$. Kolejność $ab$ jest istotna.}
\item \textbf{Macierz z trzema generatorami.} Niech $a,b,c\in R=\Z[\pi]$ i $A=\begin{psmallmatrix}1&a&c\\0&1&b\\0&0&1\end{psmallmatrix}$. Znajdź $A^{-1}$ z zachowaniem kolejności iloczynów w $R$, rozłóż $A$ na macierze elementarne i oblicz klasę torsji dwuwyrazowego kompleksu z różniczką $A$.
\wskazowka{W odwrotności macierzy trójkątnej wpis w prawym górnym rogu zawiera iloczyn $ab$. Następnie redukuj wyrazy poza przekątną za pomocą macierzy elementarnych.}
\item \textbf{Torsja sumy prostej.} Niech $C$ i $D$ będą skończonymi bazowanymi acyklicznymi kompleksami wolnych $R$-modułów. Używając definicji torsji i blokowych baz, ustal zależność między $\tau(C\oplus D)$, $\tau(C)$ i $\tau(D)$ w $\Wh(\pi)$. Zastosuj ją do sumy dwóch kompleksów z jedynymi różniczkami $(u)$ i $(u^{-1})$ dla $u\in R^\times$.
\wskazowka{Suma prosta kompleksów daje macierz blokowo diagonalną. W $K_1$ iloczyn klas przechodzi na sumę w zapisie addytywnym torsji.}
\item \textbf{Przyklejanie kobordyzmów.} Załóż, że $W_1$ jest $h$-kobordyzmem od $N_0$ do $N_1$, a $W_2$ od $N_1$ do $N_2$. Wyraź torsję sklejenia $W_1\cup_{N_1}W_2$ przez torsje $W_1,W_2$, precyzując, przez jaki izomorfizm grup Whiteheada przenosi się drugą klasę. Rozpatrz przypadek klas $\alpha$ i $-\alpha$.
\wskazowka{Skorzystaj ze wzoru addytywności torsji i naturalnego izomorfizmu grup podstawowych wywołanego inkluzją wspólnego brzegu; najpierw sprowadź obie klasy do jednej grupy.}
\item \textbf{Dwa układy baz.} W kompleksie $0\to R^2\xrightarrow{A}R^2\to0$ przyjmij $A=\operatorname{diag}(u,1)$, $u\in R^\times$. Zmień bazę w górnym stopniu za pomocą $P=\begin{psmallmatrix}1&a\\0&1\end{psmallmatrix}$, a w dolnym za pomocą $Q=\operatorname{diag}(g,1)$, $g\in\pi$. Oblicz nową macierz różniczki i sprawdź, że klasa w $\Wh(\pi)$ się nie zmienia.
\item \textbf{Kompleks o zerowej charakterystyce Eulera.} Dla $C_3=R^2$, $C_2=R^3$, $C_1=R$ i innych wyrazów równych zeru wyznacz konieczne ograniczenia na rangi różniczek acyklicznego kompleksu. Podaj przykład różniczek nad $R=\Z$, które czynią go acyklicznym, oraz przykład spełniający $\partial^2=0$, ale nieacykliczny.
\wskazowka{Acykliczność rozpatrz kolejno od stopnia pierwszego: różniczka do $C_1$ musi być surjektywna, a poprzedzająca ją musi opisywać całe jądro.}
\item \textbf{Przeszkoda w macierzy blokowej.} Dla $u\in\Z[C_5]^\times$ i $A=\begin{psmallmatrix}u&a\\0&1\end{psmallmatrix}$ oblicz $[A]\in\Wh(C_5)$. Korzystając z jednostki $u$ ze skryptu, rozstrzygnij, czy dodanie wyrazu $a$ może usunąć niezerową przeszkodę.
\wskazowka{Macierz rozkłada się na diagonalną i elementarną. Klasa Whiteheada pozostaje bez zmiany po dodaniu takiego wyrazu pozadiagonalnego.}
\item \textbf{Zmiana orientacji i podniesień komórek.} Wyjaśnij na bazowanym kompleksie nakrycia uniwersalnego, jak zmieniają się macierze brzegowe przy zmianie wybranego podniesienia jednej komórki przez działanie $g\in\pi$. Oblicz wpływ tej operacji i odwrócenia orientacji komórki na torsję Whiteheada.
\item \textbf{Wykrywanie przeszkody bez homologii.} Rozważ dwa dwuwyrazowe kompleksy nad $\Z[C_5]$ o różniczkach odpowiednio $A=I_2$ i $B=\operatorname{diag}(u,1)$, gdzie $u$ jest nietrywialną jednostką ze skryptu. Porównaj ich homologie, torsje i wyniki po zastosowaniu augmentacji $\Z[C_5]\to\Z$.
\item \textbf{Skutek geometryczny niezerowej torsji.} Niech $m\ge6$ i niech $h$-kobordyzm $(W;N_-,N_+)$ ma torsję reprezentowaną przez $\operatorname{diag}(u,u^{-1})$. Rozstrzygnij, czy $W$ musi być produktem względem $N_-$. Powtórz, gdy macierz torsji zastąpimy przez $\operatorname{diag}(u,1)$ i $[u]\ne0$.
\wskazowka{Torsję bloku diagonalnego oblicz addytywnie. Warunek z twierdzenia o $s$-kobordyzmie dotyczy klasy całej macierzy, nie każdej jej jednostki osobno.}
\end{enumerate}


\chapter{Chirurgia geometryczna i odwzorowania normalne}
\begin{enumerate}
\item \textbf{Chirurgia na rozłącznej sumie okręgów.} Dla $M=S^1\sqcup S^1$ przeprowadź chirurgię indeksu $0$ na parze punktów należących do różnych składowych oraz na parze punktów tej samej składowej. Opisz rezultat w obu przypadkach, liczbę składowych i zmianę $\chi$.
\item \textbf{Chirurgia na powierzchni rodzaju dwa.} W zamkniętej orientowalnej powierzchni $\Sigma_2$ wybierz nieoddzielającą prostą zamkniętą krzywą i wykonaj chirurgię usuwającą jej pierścieniowe otoczenie i wklejającą dwa dyski. Wyznacz wynik, grupę podstawową i charakterystykę Eulera. Co zmieni się dla krzywej oddzielającej?
\item \textbf{Sfera jako wynik chirurgii punktowej.} Dla $n\ge2$ wykonaj chirurgię indeksu $0$ na dwóch małych dyskach osadzonych w $S^n$. Wykaż, że wynik jest dyfeomorficzny z $S^{n-1}\times S^1$, i wskaż dualną chirurgię przywracającą $S^n$.
\item \textbf{Rachunek charakterystyki Eulera.} Dla chirurgii typu $(p,q)$ na $n$-rozmaitości, gdzie $p+q=n$, czyli $S^p\times D^q$ zastępujemy przez $D^{p+1}\times S^{q-1}$, wyprowadź wzór na zmianę $\chi$. Oblicz ją dla $n=4$ i $p=0,1,2$.
\wskazowka{W zmianie charakterystyki Eulera odejmujesz wkład $S^p\times D^q$ i dodajesz wkład $D^{p+1}\times S^{q-1}$. Zastosuj multiplikatywność $\chi$.}
\item \textbf{Dwa obramowania w innym wymiarze.} Wyznacz $\pi_1(SO(3))$ i opisz klasy homotopii obramowań trywialnej zorientowanej wiązki normalnej okręgu w orientowalnej $4$-rozmaitości. Wyjaśnij, które dane operacji chirurgicznej zmienia wybór obramowania.
\wskazowka{Zmiana obramowania okręgu jest pętlą w $SO(3)$; jej klasa homotopii określa różnicę dwóch framacji.}
\item \textbf{Homologia śladu chirurgii.} Niech $W$ będzie śladem pojedynczej chirurgii typu $(p,q)$ wykonanej na $M^n$. Oblicz $H_k(W,M;\Z)$ dla wszystkich $k$ i sprawdź wynik dla $p=1$, $n=5$. Opisz rolę rdzenia dołączonego uchwytu.
\wskazowka{Ślad chirurgii otrzymuje się przez dołączenie jednego uchwytu indeksu $p+1$. Względna homologia pary koncentruje się w stopniu indeksu uchwytu.}
\item \textbf{Operacja w sferze pięciowymiarowej.} Na $S^5$ wykonaj chirurgię na standardowo osadzonej $S^1$ ze standardowym obramowaniem. Wskaż możliwy model otrzymanej rozmaitości oraz policz jej grupy homologii, korzystając z rozkładu sfery na odpowiednie części.
\item \textbf{Działanie na grupę podstawową.} Rozważ chirurgię na osadzonej pętli w $M^n$ dla $n\ge5$. Opisz działanie na $\pi_1$ w zależności od tego, jaki element $[\gamma]\in\pi_1(M)$ reprezentuje pętla. Zastosuj do $M=S^1\times S^{n-1}$ i pętli stanowiącej generator $\pi_1(M)$.
\wskazowka{Dołączenie uchwytu o indeksie dwa prowadzi do relacji zabijającej klasę pętli; użyj twierdzenia van Kampena w modelu śladu.}
\item \textbf{Zanik homologii środkowej.} Załóż, że $M$ jest po prostu spójną zamkniętą $6$-rozmaitością, a klasa $x\in H_3(M;\Z)$ jest reprezentowana przez sferę z trywialną wiązką normalną. Określ, jakie dodatkowe warunki są potrzebne, by chirurgia na tej sferze mogła usunąć $x$ bez powstania niechcianej homologii. Uzasadnij odpowiedź za pomocą ciągu dokładnego pary.
\wskazowka{W środkowym wymiarze potrzebujesz kontrolować nie tylko reprezentowalność klasy przez sferę, ale też samoprzecięcia i dualne klasy przecinające.}
\item \textbf{Warunek przedłużenia mapy.} Niech $f:M\to X$ będzie odwzorowaniem ciągłym, a chirurgia przebiega wzdłuż osadzonej sfery $i:S^p\hookrightarrow M$. Wykaż, że przedłużenie $f$ na ślad uchwytu wymaga zerowania $[f\circ i]\in\pi_p(X)$ (z właściwym wyborem punktów bazowych). Czy ten warunek zapewnia przedłużenie danych normalnych?
\wskazowka{Przedłużenie mapy przez rdzeń uchwytu to problem wypełnienia mapy na sferze dyskiem. Do przedłużenia danych normalnych trzeba ponadto kontrolować obramowanie.}
\item \textbf{Względne jądro homotopii.} Dla odwzorowania stopnia $1$, $f:M\to X$, i jego wersji $f_*:H_k(M;\Z)\to H_k(X;\Z)$ oblicz jądra w przypadku $M=X\#(S^k\times S^{n-k})$ oraz mapy zgniatającej drugi składnik sumy spójnej do punktu, dla $2\le k<n-k$ i $X=S^n$.
\item \textbf{Stopień i chirurgia.} Niech $f:M^n\to X^n$ ma stopień $d$, a $F:W^{n+1}\to X$ przedłuża $f$ przez ślad chirurgii do $f':M'\to X$. Zastosuj własności klas fundamentalnych, aby porównać $\deg f$ i $\deg f'$. Podaj przykład, w którym homologiczne własności map nie wystarczają do jej przedłużenia normalnego.
\item \textbf{Granica wymiaru środkowego.} Dla chirurgii na $S^p\subset M^n$ porównaj zakresy $2p<n$ i $2p=n$. Na podstawie wymiarów oczekiwanych przecięć wskaż, dlaczego w drugim zakresie pojawiają się samoprzecięcia i forma kwadratowa.
\wskazowka{Porównaj wymiar przecięcia dwóch ogólnych zanurzonych $p$-sfer, równy $2p-n$. W wymiarze środkowym pojawiają się punkty przecięcia o wymiarze zero.}
\item \textbf{Chirurgia i orientacja.} Wyznacz, w jakich warunkach obrana orientacja na $M$ wyznacza orientację na rozmaitości po chirurgii. Przeanalizuj jawnie zgodność orientacji na wspólnym brzegu $S^p\times S^{q-1}$ dwóch sklejanych kawałków.
\item \textbf{Dwie kolejne operacje.} Na standardowym torusie $T^2$ wykonaj chirurgię na nieoddzielającej pętli, a następnie chirurgię indeksu $0$ na dwóch dyskach w otrzymanej rozmaitości. Ustal końcową liczbę składowych, genus i charakterystykę Eulera. Porównaj kolejność operacji w możliwych przypadkach.
\item \textbf{Połączenie sumy spójnej z chirurgią.} Uzasadnij, że $M^n\#(S^p\times S^{n-p})$ można dla $1\le p\le n-2$ opisać przez odpowiednio obramowaną chirurgię na $M$ tylko przy dodatkowych warunkach dotyczących wybranej sfery. Prześledź w modelu lokalnym zmianę odpowiadających części CW.
\wskazowka{Suma spójna z produktem sfer ma standardowy rozkład uchwytowy. Zidentyfikuj uchwyt tworzący generator i dualny uchwyt go znoszący.}
\item \textbf{Przykład nierozwiązywalnego problemu normalnego.} Podaj odwzorowanie $f:M\to X$ i sferę chirurgiczną $i:S^p\to M$, dla których $f\circ i$ nie jest homotopijnie trywialne. Wyjaśnij, dlaczego żaden wybór obramowania nie pozwoli przedłużyć $f$ przez dany uchwyt.
\item \textbf{Chirurgie komplementarne.} Dla $p+q=n$ porównaj indeks uchwytu śladu chirurgii i indeks uchwytu widzianego od drugiego końca. Sprawdź tę relację w wymiarach $n=3,4,5$ na nietrywialnych parach $(p,q)$.
\item \textbf{Operacja Dehna a ramowanie.} Dla niewęzła $K\subset S^3$ rozważ dołączenie $2$-uchwytu z obramowaniem całkowitym $p$. Zidentyfikuj trójwymiarowy brzeg dla $p=0,1,2$, wyznacz jego pierwszą homologię i porównaj z macierzą formy przecięcia śladu.
\wskazowka{Dla całkowitoliczbowej operacji Dehna na niewęźle użyj chirurgicznego opisu brzegu $2$-uchwytu oraz relacji $H_1\cong\operatorname{coker}(p)$.}
\end{enumerate}


\chapter{Formy kwadratowe i przeszkoda chirurgiczna}
\begin{enumerate}
\item \textbf{Inwolucja w nieprzemiennym pierścieniu.} Niech $\pi=S_3$ i $R=\Z[\pi]$ z inwolucją $\overline{\sum n_gg}=\sum n_gg^{-1}$. Dla $x=1+(12)-(123)$ i $y=(23)+(123)$ oblicz $\overline{xy}$ i $\bar y\bar x$. Zbadaj, czy $x\bar x$ jest niezmienniczy względem inwolucji.
\item \textbf{Moduł formy hermitowskiej.} Dla $R=\Z[t,t^{-1}]$ z $\bar t=t^{-1}$ rozważ $A=\begin{psmallmatrix}2&t\\t^{-1}&3\end{psmallmatrix}$. Sprawdź $A=\bar A^{\,T}$, oblicz augmentację tej macierzy i sygnaturę otrzymanej rzeczywistej formy.
\item \textbf{Zmiana bazy formy seskwiliniowej.} Jeśli $A$ reprezentuje $\varepsilon$-hermitowską formę na wolnym lewym $R$-module, wyprowadź odpowiednią regułę zmiany macierzy przy przejściu do nowej bazy. Oblicz ją dla $A=\begin{psmallmatrix}0&1\\\varepsilon&0\end{psmallmatrix}$ i diagonalnej zmiany bazy przez $u,v\in R^\times$ w przypadku przemiennego $R$.
\wskazowka{W konwencji lewych modułów dobierz konsekwentnie zapis wektorów kolumnowych i macierzy. Reguła powinna zawierać transpozycję i inwolucję pierścienia.}
\item \textbf{Hiperboliczność i lagranżjan.} Dla formy symetrycznej $H=\begin{psmallmatrix}0&I_2\\I_2&0\end{psmallmatrix}$ na $\Z^4$ znajdź dwa dopełniające się lagranżjany, oblicz sygnaturę i wskaż bazy, w których uwidacznia się para hiperboliczna.
\wskazowka{Lagrangian jest izotropową podgrupą rangi połowy całej przestrzeni, równą swojemu ortogonalnemu dopełnieniu. W macierzy hiperbolicznej widać go bez diagonalizacji.}
\item \textbf{Sygnatura a stabilizacja.} Dla $A=\operatorname{diag}(1,1,-1,-1,1,1)$ i $B=\begin{psmallmatrix}0&I_2\\I_2&0\end{psmallmatrix}$ oblicz sygnatury $A$, $B$ i $A\oplus B$. Czy dołączenie $B$ zmienia klasę w grupie Wittowskiej symetrycznych form nad $\R$?
\item \textbf{Parzysta forma unimodularna.} Rozważ $A=\begin{psmallmatrix}2&1\\1&2\end{psmallmatrix}$. Oblicz jej wyznacznik, parzystość i sygnaturę. Wyjaśnij, dlaczego nie reprezentuje unimodularnej formy przecięć zamkniętej po prostu spójnej $4$-rozmaitości.
\item \textbf{Forma antysymetryczna.} Dla $J=\begin{psmallmatrix}0&1\\-1&0\end{psmallmatrix}$ i $A=J\oplus J$ na $\Z^4$ znajdź bazy symplektyczne, oblicz wyznacznik i pokaż, że każda z dwóch widocznych par dopuszcza lagranżjan. Odnieś wynik do nieparzystego przypadku teorii chirurgii.
\item \textbf{Dwie formy kwadratowe na płaszczyźnie nad $\mathbf F_2$.} Niech $V=(\Z/2\Z)^2$ z formą dwuliniową $b((x,y),(x',y'))=xy'+x'y$. Zbadaj funkcje $q_0(x,y)=xy$ i $q_1(x,y)=xy+x+y$. Sprawdź tożsamość kwadratową i oblicz oba niezmienniki Arfa, licząc zera każdej funkcji.
\wskazowka{Niezmiennik Arfa licz przez bazę symplektyczną albo przez liczbę zer kwadratowego uszczegółowienia; sprawdź tożsamość $q(v+w)=q(v)+q(w)+b(v,w)$.}
\item \textbf{Suma ortogonalna i Arf.} Dla dwóch płaszczyzn kwadratowych $(V_0,q_0)$ i $(V_1,q_1)$ z poprzedniego zadania oblicz Arf formy $q_0\oplus q_1$. Sprawdź wynik bezpośrednio, wybierając bazę symplektyczną dla sumy.
\item \textbf{Kwadratowe uszczegółowienie.} Dla symplektycznej formy $b$ na $(\Z/2\Z)^4$ z macierzą $J_2\oplus J_2$ nad $\mathbf F_2$ podaj wszystkie możliwe wartości $q(e_i)$ dla wybranej bazy symplektycznej i skonstruuj przykład uszczegółowienia $q$ o Arf równym $1$.
\wskazowka{W bazie symplektycznej wartości $q$ na wektorach bazowych można wybierać niezależnie, po czym tożsamość kwadratowa wyznacza wartości na sumach.}
\item \textbf{Rachunek przecięcia z grupą podstawową.} Niech dwie zanurzone sfery wymiaru środkowego w nakryciu uniwersalnym mają punkty przecięcia etykietowane przez $g,h,g$ i znaki $+,-,+$. Zapisz wynik parowania w $\Z[\pi]$. Oblicz jego augmentację i wyjaśnij, kiedy dwa punkty o przeciwnych znakach mogą tworzyć parę Whitneya.
\wskazowka{Do każdej pary punktów przecięcia dołącz etykietę z $\pi_1$; augmentacja zapomina etykiety, ale nie przesądza istnienia dysku Whitneya.}
\item \textbf{Etykiety po zmianie podniesienia.} Niech $\lambda(x,y)\in\Z[\pi]$ oznacza parowanie dwóch klas, a ich podniesienia zmieniono przez $g,h\in\pi$. Wyznacz, jak zmienia się współczynnik parowania zgodnie z ustaloną konwencją lewego modułu, i sprawdź zachowanie sygnatury po augmentacji w przypadku $\pi=1$.
\wskazowka{Zapisz seskwiliniowość parowania osobno w obu zmiennych i koniecznie ustal, czy moduły są lewe, czy prawe.}
\item \textbf{Kiedy nie ma przeszkody środkowej?} Załóż, że zwarta po prostu spójna $2k$-rozmaitość ma jądro chirurgii w stopniu $k$ z formą stanowiącą sumę trzech płaszczyzn hiperbolicznych. Oblicz jej klasę w odpowiedniej grupie $L$ i opisz, jakie przeszkody geometryczne trzeba jeszcze skontrolować przed wykonaniem chirurgii.
\item \textbf{Skutek sygnatury osiem.} Rozważ nieosobliwą parzystą formę całkowitą $E_8\oplus H$, gdzie $H$ jest płaszczyzną hiperboliczną. Oblicz sygnaturę i wartość przeszkody sygnaturowej w $L_0(\Z)$ dla standardowej normalizacji. Czy dołączenie kolejnej płaszczyzny hiperbolicznej może ją wyzerować?
\wskazowka{Sygnatura płaszczyzny hiperbolicznej jest zerowa. Dla form parzystych unimodularnych porównaj wynik z normalizacją generatora $L_0(\Z)$.}
\item \textbf{Porównanie wymiarów $4k$ i $4k+2$.} Korzystając z opisu $L_n(\Z)$ przedstawionego w skrypcie, wyznacz grupy i sposób wykrywania klasy przeszkody w wymiarach $6$, $8$, $10$ i $12$. Podaj dla każdego przykładowe dane algebraiczne o zerowej i niezerowej przeszkodzie.
\item \textbf{Pierścień Laurentowski i rozszczepienie.} Dla $\pi=\Z$ wykorzystaj opis grup $L_n(\Z[t,t^{-1}])$ podany w skrypcie, aby obliczyć je w czterech kolejnych wymiarach w stosownej dekoracji. Wskaż, który składnik pochodzi z $L_{n-1}(\Z)$, i wyjaśnij rolę wyboru dekoracji.
\wskazowka{Użyj właściwej wersji rozszczepienia dla pierścienia Laurentowskiego wraz z dokładnie tą dekoracją grup $L$, która występuje w skrypcie.}
\item \textbf{Przeszkoda odwzorowania stopnia jeden.} Załóż, że odwzorowanie normalne $f:M^{4k}\to X^{4k}$ ma po prostu spójny cel i nieosobliwą formę na jądrze środkowej homologii o sygnaturze $16$. Oblicz przeszkodę chirurgiczną przy standardowej konwencji i wskaż, co wynika z jej niezerowości.
\item \textbf{Kontrola formy przez wektory izotropowe.} Dla $A=\begin{psmallmatrix}0&1&0&0\\1&0&0&0\\0&0&2&1\\0&0&1&2\end{psmallmatrix}$ znajdź prymitywne wektory izotropowe i zbadaj, czy posiada ona lagranżjan rangi dwa w sensie form unimodularnych. Jakie znaczenie ma wyznacznik $A$?
\wskazowka{Sprawdź nieosobliwość formy i warunek rangi lagranżjanu. Sama obecność jednego prymitywnego wektora izotropowego nie oznacza metaboliczności.}
\item \textbf{Relacja między $K_1$ a $L$-teorią.} Wyjaśnij na przykładzie formy hiperbolicznej nad $\Z[\pi]$, dlaczego zerowanie torsji macierzy przejścia nie jest tym samym warunkiem co zerowanie klasy chirurgicznej. Sformułuj przykład algebraiczny ujawniający różnicę między tymi dwiema przeszkodami.
\item \textbf{Metaboliczność i sygnatura.} Dla form $H$, $H\oplus H$ oraz $I_{1,1}$ wyznacz maksymalne podgrupy izotropowe. Które z tych form są metaboliczne nad $\Z$? Porównaj z nieunimodularną formą $\operatorname{diag}(1,-2)$, która ma sygnaturę zero, lecz nie zawiera niezerowego wektora izotropowego.
\wskazowka{Sygnatura zero jest warunkiem słabszym niż istnienie lagranżjanu całkowitoliczbowego. Przed szukaniem izotropii przeanalizuj równanie kwadratowe $q(v)=0$.}
\end{enumerate}


\chapter{Ciąg dokładny chirurgii i zbiór struktur}
\begin{enumerate}
\item \textbf{Struktura jako dane, nie tylko rozmaitość.} Dla dwóch równoważności homotopijnych $f_i:M_i\to X$ zapisz warunek, pod którym reprezentują ten sam element zbioru struktur $\mathcal S(X)$. Wytłumacz na przykładzie złożenia z nietrywialną samorównoważnością homotopijną $X$, dlaczego sam typ dyfeomorficzny $M_i$ nie rozstrzyga równoważności struktur.
\item \textbf{Włókno normalnego niezmiennika.} Traktując ciąg chirurgii jako ciąg zbiorów z wyróżnionymi elementami i działaniami grup, wyjaśnij, co oznacza dokładność w punkcie $[X,G/O]$. Zastosuj ją do sytuacji, w której zadany normalny niezmiennik ma niezerową przeszkodę w $L_n(\Z[\pi_1X])$.
\wskazowka{Dokładność przy normalnym niezmienniku oznacza, że jego przeszkoda chirurgiczna znika dokładnie wtedy, gdy leży w obrazie mapy ze zbioru struktur.}
\item \textbf{Punkt bazowy i działanie grupy $L$.} Zapisz na schemacie, jak element $L_{n+1}(\Z[\pi_1X])$ działa na zbiór struktur, oraz jak stabilizator wybranego elementu wpływa na rozmiar jego orbity. Uzasadnij, dlaczego nie wolno traktować całego ciągu automatycznie jako ciągu grup abelowych.
\wskazowka{Narysuj orbitę elementu zbioru struktur pod działaniem grupy $L_{n+1}$. O liczbie jej elementów decyduje iloraz przez stabilizator.}
\item \textbf{Cel po prostu spójny w wymiarze siedem.} Używając odpowiednich grup $L_n(\Z)$, uprość fragment ciągu chirurgii dla po prostu spójnej zamkniętej $7$-rozmaitości $X$. Wyznacz, w których miejscach grupa przeszkód jest zerowa, i wskaż, jakie informacje nadal mogą być potrzebne do wyznaczenia $\mathcal S(X)$.
\wskazowka{Podstaw grupy $L_{7}(\Z)$ i $L_{8}(\Z)$ do odpowiedniego fragmentu ciągu; nie utożsamiaj automatycznie zerowania przeszkody z jednoznacznością struktury.}
\item \textbf{Struktury na sferze.} Dla $X=S^7$ porównaj zbiór struktur z grupą sfer homotopijnych $\Theta_7$, uwzględniając oznaczenia i konwencje skryptu. Wyjaśnij, dlaczego fakt $\pi_1(S^7)=0$ nie wymusza jednoelementowości zbioru struktur.
\wskazowka{Porównaj definicję struktury nad $S^7$ z równoważnościami homotopijnymi sfer homotopijnych; uwzględnij wpływ automorfizmów homotopijnych celu.}
\item \textbf{Zerowa przeszkoda a brak jednoznaczności.} Skonstruuj abstrakcyjny przykład ciągu dokładnego zbiorów z działaniem, w którym dwa różne elementy zbioru struktur mają ten sam normalny niezmiennik. Wyraź algebraicznie, jakie dane o działaniu $L_{n+1}$ rozróżniają te struktury.
\item \textbf{Zadanie z ustalonym brzegiem.} Niech $(X,\partial X)$ jest zwartą rozmaitością z brzegiem, a $f:(M,\partial M)\to(X,\partial X)$ jest równoważnością homotopijną ograniczającą się do ustalonego dyfeomorfizmu brzegu. Napisz właściwą wersję zbioru struktur i wymień, które deformacje normalne muszą być relatywne do brzegu.
\item \textbf{Dysk relatywnie do brzegu.} Rozważ $D^n$ z ustalonym punkt po punkcie brzegiem, $n\ge5$. Wyjaśnij, jak sklejanie dwóch dysków wzdłuż brzegu wiąże możliwe struktury względne z homotopijnymi sferami wymiaru $n$. Które dodatkowe dane identyfikacji brzegu mają znaczenie?
\wskazowka{W konstrukcji przez podwojenie dysku decydujące znaczenie ma wybrana identyfikacja na brzegu. Porównaj ją z definicją grupy sfer homotopijnych.}
\item \textbf{Homotopijny niezmiennik normalny.} Dla dwóch homotopijnych odwzorowań do $G/O$ wyjaśnij, dlaczego reprezentują tę samą klasę normalną. Przeanalizuj rolę wybranego stabilnego izomorfizmu wiązek w definicji odwzorowania normalnego.
\item \textbf{Skutek $L_n=0$.} Załóż w odpowiednim wymiarze i kategorii, że $L_n(\Z[\pi_1X])=0$. Wyprowadź z dokładności wniosek o obrazie odwzorowania $\mathcal S(X)\to[X,G/O]$. Wyjaśnij, dlaczego nie wynika stąd automatycznie ani iniektywność, ani trywialność $\mathcal S(X)$.
\wskazowka{Z dokładności w punkcie $[X,G/O]$ wnioskuj jedynie o obrazie mapy normalnej. O jej włóknach mówi działanie poprzedniej grupy $L$.}
\item \textbf{Rola wymiaru w teorii.} Podaj po dwa przykłady wymiarów i grup podstawowych, w których klasyczny ciąg chirurgii ma zastosowanie w sformułowaniu ze skryptu albo wymaga dodatkowych założeń. Dla każdego wskaż konkretny etap argumentu, w którym pojawia się wymóg wysokiego wymiaru.
\item \textbf{Normalny niezmiennik kompozycji.} Rozważ równoważność homotopijną $f:M\to X$ i dyfeomorfizm $h:N\to M$. Porównaj normalne niezmienniki $f$ i $f\circ h$. Ustal, w jakim sensie wynik zależy od ustalonych identyfikacji wiązek normalnych.
\item \textbf{Dwie dekoracje grupy $L$.} Porównaj role $L^s$ i $L^h$ w klasyfikacji prostych i zwykłych równoważności homotopijnych. Wyjaśnij, jaki problem może powstać przy zastąpieniu jednej dekoracji drugą, jeśli $\Wh(\pi_1X)\ne0$.
\wskazowka{Dekoracja $s$ odpowiada prostym równoważnościom homotopijnym; przejście do $h$ wymaga kontroli torsji Whiteheada.}
\item \textbf{Obliczenia w czterech stopniach.} Dla $X$ po prostu spójnego i wymiarów $n=6,7,8,9$ wstaw do schematu ciągu chirurgii znane grupy $L_n(\Z)$ i $L_{n+1}(\Z)$. Porównaj miejsca, w których istnieje przeszkoda istnienia struktury, i miejsca, w których może występować nietrywialne działanie na struktury.
\item \textbf{Przeszkoda w wymiarze $4k$.} Przyjmując identyfikację $L_{4k}(\Z)\cong\Z$ przez odpowiednio znormalizowaną sygnaturę, porównaj przeszkody reprezentowane przez formy $H$, $E_8$ i $E_8\oplus(-E_8)$. Oddziel obliczenie klasy algebraicznej od wniosku o realizacji chirurgii.
\wskazowka{Klasy przeszkód są wyznaczane przez sygnaturę po właściwej normalizacji. Porównaj algebraiczną klasę z warunkiem istnienia geometrycznej realizacji.}
\end{enumerate}


\chapter{Hipoteza Poincarégo w wymiarze dwa i w wysokich wymiarach}
\begin{enumerate}
\item \textbf{Powierzchnia z daną charakterystyką.} Niech $M$ będzie zamkniętą spójną orientowalną powierzchnią o $\chi(M)=-6$. Wyznacz jej rodzaj, grupy homologii całkowitej i minimalną liczbę generatorów $\pi_1(M)$. Wyjaśnij, czy $M$ może być po prostu spójna.
\item \textbf{Dwie topologie z tą samą charakterystyką.} Znajdź orientowalną i nieorientowalną zamkniętą powierzchnię o tej samej ujemnej charakterystyce Eulera. Oblicz ich pierwsze grupy homologii z $\Z$ i z $\Z/2\Z$ i wskaż, który niezmiennik pozwala je rozróżnić.
\item \textbf{Skończona grupa podstawowa powierzchni.} Korzystając z klasyfikacji powierzchni, sklasyfikuj zamknięte spójne powierzchnie z niezerową skończoną grupą podstawową. Czy orientowalna powierzchnia może mieć niezerową skończoną $\pi_1$?
\item \textbf{Sfera homotopijna z wyciętymi dyskami.} Niech $\Sigma^n$ będzie sferą homotopijną, $n\ge6$, a $W$ powstaje przez usunięcie wnętrz dwóch rozłącznych dysków. Oblicz grupy względne $H_*(W,S^{n-1};\Z)$ i wyjaśnij, jak wykorzystać twierdzenie Smale'a do rozpoznania $W$.
\wskazowka{Po wycięciu dwóch dysków użyj długiego ciągu pary i van Kampena. Przed cytowaniem Smale'a sprawdź osobno obie inkluzje brzegowe.}
\item \textbf{Po co dodatkowe ograniczenie wymiaru?} Dla poprzedniego zadania porównaj wymiary kobordyzmów otrzymanych ze sfer homotopijnych wymiarów $4,5,6$. Wskaż dokładnie, w których przypadkach klasyczne twierdzenie o $h$-kobordyzmie ma wystarczająco duży wymiar, a w których nie rozstrzyga problemu.
\item \textbf{Sfera homotopijna a homologiczna.} Podaj zamkniętą $3$-rozmaitość mającą homologie całkowite sfery $S^3$, ale niebędącą sferą homotopijną. Wyjaśnij, którego warunku brakuje i jak zmienia się sytuacja przy założeniu prostej spójności.
\item \textbf{Homologia po prostu spójnej pięciorozmaitości.} Niech $M^5$ będzie zamknięta, zorientowana i po prostu spójna, a $H_2(M;\Z)=\Z^r\oplus T$ z $T$ skończoną. Użyj dualności Poincarégo i twierdzenia o współczynnikach uniwersalnych, aby wyrazić grupy homologii i kohomologii w pozostałych stopniach.
\wskazowka{Najpierw policz kohomologie przez dualność Poincarégo, a następnie użyj twierdzenia o współczynnikach uniwersalnych do wykrycia miejsc, gdzie pojawia się torsja.}
\item \textbf{Zakaz samotnej torsji nieparzystej.} Załóż, że $M^5$ jest zamknięta i po prostu spójna. Wyjaśnij, dlaczego $\operatorname{Tor}H_2(M;\Z)$ nie może być izomorficzna z $\Z/3\Z$, korzystając z niezdegenerowanej antysymetrycznej formy wiązania. Czy przeszkoda znika dla $(\Z/3\Z)^2$?
\wskazowka{Na części torsyjnej istnieje niesingularne parowanie antysymetryczne z wartościami w $\Q/\Z$; poszukaj partnera dla niezerowego elementu.}
\item \textbf{Macierz formy wiązania.} Na $T=(\Z/7\Z)^2$ rozważ $b((a,b),(c,d))=(ad-bc)/7\pmod\Z$. Sprawdź niezdegenerowanie i antysymetrię formy oraz wyznacz jej grupę automorfizmów wśród macierzy $2\times2$ nad $\Z/7\Z$.
\wskazowka{Dla macierzy formy nad $\Z/7$ warunek zachowania standardowej formy symplektycznej sprowadza się do równania macierzowego $A^TJA=J$.}
\item \textbf{Klasa $w_2$ na części wolnej.} Niech $H_2(M;\Z)=\Z^2$ dla zamkniętej po prostu spójnej $5$-rozmaitości. Ile różnych homomorfizmów $H_2\to\Z/2\Z$ może reprezentować $w_2$? Sklasyfikuj ich orbity względem działania $GL_2(\Z)$ i zinterpretuj różnicę między przypadkiem spinowym a niespinowym.
\wskazowka{Rozważ działanie $GL_2(\Z)$ na wektorach z $(\Z/2)^2$. Przy redukcji modulo dwa porównaj orbitę zera i niezerowych wektorów.}
\item \textbf{Dane Bardena dla dwóch modeli.} Porównaj modele $S^2\times S^3$ oraz nietrywialną zorientowaną wiązkę $S^3$ nad $S^2$. Oblicz $H_2$, porównaj $w_2$ i zastosuj twierdzenie klasyfikacyjne, aby wyjaśnić, dlaczego nie są dyfeomorficzne, mimo tej samej homologii.
\wskazowka{Dwie wiązki nad $S^2$ mają taką samą homologię przestrzeni całkowitej, ale różnią się klasą $w_2$; zestaw ją z klasyfikacją Bardena.}
\item \textbf{Składniki torsyjne w modelach $M_k$.} Dla modeli z części o klasyfikacji Bardena oblicz $H_2(M_3;\Z)$ i $H_2(M_8;\Z)$. Określ ich liczbę generatorów i porównaj formy wiązania w świetle danych podanych w skrypcie.
\item \textbf{Kiedy nie da się zastosować klasyfikacji?} Wyjaśnij, dlaczego klasyfikacja po prostu spójnych zamkniętych $5$-rozmaitości nie klasyfikuje automatycznie zwartych ściągalnych $5$-rozmaitości z brzegiem $S^4$. Zidentyfikuj różnicę w kategoriach i w założeniach.
\item \textbf{Iloczyn z $S^1$ a hipoteza Poincarégo.} Niech $\Sigma^n$ będzie sferą homotopijną. Wyznacz $\pi_1(\Sigma\times S^1)$ i jej homologie. Wyjaśnij, dlaczego nawet gdy $\Sigma$ jest homeomorficzna z $S^n$, o gładkim typie $\Sigma\times S^1$ nie wolno wnioskować jedynie z homotopii.
\end{enumerate}


\chapter{Sfery Milnora i egzotyczne struktury gładkie}
\begin{enumerate}
\item \textbf{Działanie funkcji sklejającej.} Dla kwaternionów jednostkowych $u\in S^3$ rozważ $g_{h,j}(u)(v)=u^hvu^j$. Sprawdź, że $g_{h,j}(u)$ zachowuje normę i orientację $\R^4$ oraz wyznacz jej wartość w $u=1$.
\item \textbf{Dwie klasy charakterystyczne.} Korzystając z wzorów $e(E_{h,j})=(h+j)\alpha$, $p_1(E_{h,j})=2(h-j)\alpha$ (przy spójnie ustalonej orientacji), oblicz te klasy dla $(h,j)=(3,-2), (2,-1), (4,-3)$ oraz $(-1,2)$. Które wiązki mają tę samą klasę Eulera, a które ten sam $p_1$?
\item \textbf{Gdy całkowita przestrzeń nie jest sferą homotopijną.} Wyznacz całkowite homologie wiązki sferycznej $S(E_{h,j})\to S^4$ w przypadku $h+j=3$. Wykorzystaj ciąg Gysina albo Mayera--Vietorisa i porównaj z przypadkiem $h+j=1$.
\wskazowka{Użyj ciągu Gysina dla wiązki $S^3\to M\to S^4$. Kluczowe odwzorowanie mnoży przez liczbę Eulera wiązki.}
\item \textbf{Przypadek zerowej klasy Eulera.} Załóż $h+j=0$. Oblicz grupy homologii całkowitej przestrzeni $S(E_{h,j})$ oraz porównaj je z homologiami $S^3\times S^4$. Czy sama zgodność homologii wystarcza, by uznać przestrzenie za dyfeomorficzne?
\item \textbf{Warunek na sferę homotopijną.} Dla $h+j=\pm1$ przeanalizuj ciąg Mayera--Vietorisa rozkładu na dwie części nad półsferami. Oblicz wyznacznik macierzy występującej w odwzorowaniu na $H_3$, a następnie wyznacz homologie całkowitej przestrzeni.
\wskazowka{W rozkładzie przestrzeni całkowitej na dwie trywializacje nad półsferami śledź homomorfizm na $H_3$ wyznaczany przez funkcję sklejającą.}
\item \textbf{Parametr $k=h-j$.} Przy $h+j=1$ pokaż, że $k=h-j$ jest nieparzysty, i odwrotnie, dla dowolnego nieparzystego $k$ wyznacz odpowiadającą parę $(h,j)$. Zapisz $p_1(E_{h,j})$ przez $k$.
\item \textbf{Niezmiennik Milnora w przykładach.} Dla $k=9,11,13,15,17,19$ oblicz $\lambda(M_k)\in\Z/7$ na podstawie wzoru ze skryptu. Wskaż pary różnych parametrów, których ten niezmiennik nie rozróżnia.
\wskazowka{Oblicz kwadraty parametrów modulo odpowiednio dobraną liczbę; dla wartości $\lambda$ uwzględnij tylko właściwą kongruencję ze skryptu.}
\item \textbf{Dokładniejszy niezmiennik Eellsa--Kuipera.} Dla $k=9,11,13,15,17,19$ oblicz $\mu(M_k)=(k^2-1)/224\pmod\Z$. Sprawdź, które z nich reprezentują różne klasy zorientowanego dyfeomorfizmu.
\wskazowka{Najpierw wyznacz liczby $k^2-1$, a następnie porównuj je modulo $224$, pamiętając, że niezmiennik ma wartości w $\Q/\Z$.}
\item \textbf{Kongruencja klasyfikacyjna.} Rozstrzygnij, czy $M_3$ i $M_{115}$ są zorientowanie dyfeomorficzne. Następnie znajdź najmniejsze dodatnie nieparzyste $l>3$, dla którego $M_l$ jest zorientowanie dyfeomorficzna z $M_3$.
\wskazowka{Równość niezmienników sprowadź do kongruencji $k^2\equiv l^2\pmod{224}$. Dodatkowo sprawdź zakres stosowalności kryterium klasyfikacyjnego.}
\item \textbf{Odwrócenie orientacji.} Oblicz, jak zmieniają się $\lambda$ i $\mu$ po odwróceniu orientacji sfery homotopijnej. Znajdź przykład parametrów $k,l$, dla których $M_k$ i $M_l$ są dyfeomorficzne po odwróceniu orientacji jednej z nich, ale nie są zorientowanie dyfeomorficzne.
\wskazowka{Odwrócenie orientacji zmienia znak addytywnego niezmiennika w $\Q/\Z$. Nie myl równoważności zorientowanej z niezorientowaną.}
\item \textbf{Dodawanie w $\Theta_7$.} Przyjmując $\Theta_7\cong\Z/28$ i addytywność $\mu$, wyznacz rząd klasy $[M_3]$, znając $\mu(M_3)=1/28$. Oblicz klasy $[M_3]\#[M_5]$ i $[M_3]\#(-[M_3])$ przez ich niezmienniki.
\item \textbf{Realizacja klas w rodzinie Milnora} Korzystając z wyliczonego już w skrypcie zbioru szesnastu wartości $28\mu(M_k)\pmod{28}$, rozstrzygnij, czy klasy o niezmiennikach $2/28$, $5/28$, $8/28$ i $25/28$ występują w tej rodzinie. W przypadku klasy $8/28$ znajdź jeden parametr $k$.
\item \textbf{Podwojenie wypełnienia.} Niech $B$ będzie zwartą zorientowaną spinową $8$-rozmaitością z brzegiem będącym sferą homotopijną. Rozważ jej podwojenie $B\cup_{\partial B}(-B)$. Oblicz sygnaturę podwojenia i określ, co można wnioskować o klasach Pontriagina; nie zakładaj bez uzasadnienia, że całe podwojenie ma trywialną wiązkę styczną.
\wskazowka{Przy sklejaniu $B$ i $-B$ sygnatury sumują się z uwzględnieniem odwróconej orientacji. Oddziel zachowanie liczb charakterystycznych od twierdzenia o orientowalnym spinowym podwojeniu.}
\item \textbf{Rola klasy Pontriagina.} Dla dwóch wiązek $E_{h,j}$ o $h+j=1$ i różnych $h-j$ porównaj $p_1(E)$ i znormalizowany niezmiennik rozmaitości brzegowej. Podaj przykład różnych wiązek dających zorientowanie dyfeomorficzne sfery brzegowe.
\wskazowka{Klasa $p_1(E_{h,j})$ zależy liniowo od $h-j$, a niezmiennik brzegu od odpowiedniej kongruencji kwadratowej; porównaj oba rodzaje informacji.}
\item \textbf{Grupa sfer homotopijnych.} Wyjaśnij na poziomie operacji sumy spójnej, dlaczego elementem neutralnym w $\Theta_7$ jest $[S^7]$, a odwrotnością klasy $[\Sigma]$ jest $[-\Sigma]$. Wskaż twierdzenie ze skryptu potrzebne, by przejść od $h$-kobordyzmu do dyfeomorfizmu.
\item \textbf{Homeomorfizm a dyfeomorfizm.} Dla $M_3$ ustal, jakie niezmienniki homotopijne i topologiczne zgadzają się z $S^7$, a który obliczony w skrypcie niezmiennik gładki je rozróżnia. Podaj powód, dla którego argument nie jest sprzeczny z wysokowymiarową hipotezą Poincarégo.
\item \textbf{Zmiana pary przez symetrię.} Porównaj wiązki $E_{h,j}$ i $E_{j,h}$, ich klasy Eulera i Pontriagina. W przypadku $h+j=1$ zbadaj relacje między parametrami $k$ i $-k$ i wartościami $\mu(M_k)$; opisz, czego sama równość niezmiennika nie mówi o wybranym izomorfizmie wiązek.
\item \textbf{Działanie odwrócenia orientacji na klasach sfer.} W grupie $\Theta_7\cong\Z/28$ ustal klasy o rzędzie $1$, $2$, $4$, $7$, $14$ i $28$. Dla których klas odwrócenie orientacji daje tę samą klasę w $\Theta_7$?
\end{enumerate}


\chapter{Teoria Cherna--Weila i Cherna--Gaussa--Bonneta}
\begin{enumerate}
\item \textbf{Krzywizna koneksji abelowej.} Na trywialnej zorientowanej wiązce płaszczyzn nad $\R^2$ rozważ formę koneksji $A=\begin{psmallmatrix}0&-\alpha\\\alpha&0\end{psmallmatrix}$, gdzie $\alpha=x^2\,dy-y\,dx$. Oblicz $\Omega=dA+A\wedge A$ i formę Eulera $e_\nabla(E)$.
\item \textbf{Druga koneksja, ta sama klasa.} Na torusie $T^2=(\R/\Z)^2$ w trywialnej wiązce orientowanych płaszczyzn weź $\alpha_0=0$ i $\alpha_1=\sin(2\pi x)\,dy$. Oblicz formy Eulera obu koneksji, znajdź jawnie ich różnicę jako formę dokładną i oblicz całki po torusie.
\item \textbf{Pfaffian w wymiarze cztery.} Dla macierzy form $\Omega=(\Omega_{ij})_{1\le i,j\le4}$ z $\Omega_{ji}=-\Omega_{ij}$ rozpisz $\operatorname{Pf}(\Omega)$ za pomocą trzech iloczynów klinowych. Zastosuj do bloków $\Omega_{12}=\alpha$, $\Omega_{34}=\beta$ i pozostałych wyrazów zerowych.
\wskazowka{Dla antysymetrycznej macierzy $4\times4$ Pfaffian ma trzy składniki. Wygodnie rozwinąć go według pierwszego wiersza.}
\item \textbf{Dwa kierunki krzywizny.} Niech $E=E_1\oplus E_2$ będzie ortogonalną sumą dwóch orientowanych wiązek rzędu $2$ o koneksji blokowo diagonalnej. Oblicz formy $e_\nabla(E)$ i $p_{1,\nabla}(E)$ przez $\alpha=e_\nabla(E_1)$ i $\beta=e_\nabla(E_2)$.
\wskazowka{W koneksji blokowej krzywizna nie miesza dwóch bloków. Użyj mnożeniowości klasy Eulera oraz addytywności $p_1$.}
\item \textbf{Iloczyn dwóch sfer.} Wyposaż $S^2(r)\times S^2(s)$ w metrykę produktową. Korzystając z Pfaffianu krzywizny oblicz całkę formy Eulera po tym produkcie, a następnie porównaj ją z $\chi(S^2\times S^2)$. Zbadaj zależność od $r,s$.
\wskazowka{Dla metryki produktowej krzywizny w dwóch niezależnych płaszczyznach to $1/r^2$ i $1/s^2$. W Pfaffianie porównaj je z elementem objętości.}
\item \textbf{Torus i płaska koneksja.} Na $T^2$ oblicz formę Eulera standardowej płaskiej koneksji. Następnie zinterpretuj niezerowe, lecz dokładne, formy Eulera wynikające ze zmiany koneksji. Czy forma punktowo niezerowa musi reprezentować niezerową klasę Eulera?
\item \textbf{Zmiana orientacji.} Ustal, jak zmienia się $\operatorname{Pf}(\Omega)$ po odwróceniu orientacji rzeczywistej wiązki rzędu $2m$. Sprawdź rachunek na przykładzie macierzy blokowej rzędu $4$ i porównaj wynik z zachowaniem $p_1(E)$.
\wskazowka{Pfaffian zmienia znak przy zmianie orientacji ramy o wyznaczniku ujemnym; wielomiany Pontriagina pozostają niezmienione.}
\item \textbf{Wielomian niezmienniczy.} Dla $P(A,B)=\operatorname{tr}(AB)$ na $\mathfrak{so}(n)$ sprawdź niezmienniczość przy jednoczesnym sprzężeniu przez $SO(n)$. Oblicz $P(\Omega,\Omega)$ dla koneksji blokowej z dwoma blokami $2\times2$.
\item \textbf{Transgresja dla linii zespolonej.} Na trywialnej hermitowskiej linii nad $T^2$ rozważ koneksje $d+i\alpha_0$ i $d+i\alpha_1$. Wyprowadź formę transgresji pierwszej klasy Cherna i sprawdź ją dla $\alpha_0=0$, $\alpha_1=\sin(2\pi x)\,dy$.
\wskazowka{Dla koneksji liniowej $d+i\alpha$ krzywizna jest formą $i\,d\alpha$. Różnica przedstawicieli $c_1$ jest pochodną zewnętrzną różnicy potencjałów.}
\item \textbf{Chern--Simons w rodzinie.} Dla $A_t=A_0+tB$ wyprowadź jawny wzór na formę transgresji odpowiadającą niezmienniczemu wielomianowi stopnia $2$, $P(X,Y)=\operatorname{tr}(XY)$. Zastosuj go do przypadku przemiennego, w którym $[A_0,B]=0$ i $B\wedge B=0$.
\wskazowka{Użyj rodziny krzywizn $F_t=dA_t+A_t\wedge A_t$ i oblicz $\frac{d}{dt}P(F_t,F_t)$; następnie zastosuj niezmienniczość $P$.}
\item \textbf{Gauss--Bonnet dla pierścienia.} Niech $A=\{r^2\le x^2+y^2\le R^2\}\subset\R^2$ z metryką euklidesową, $0<r<R$. Oblicz całki krzywizny geodezyjnej na obu składowych zorientowanego brzegu i sprawdź wzór Gaussa--Bonneta z brzegiem.
\item \textbf{Sektor kołowy z narożnikami.} Dla sektora dysku o kącie $\theta\in(0,2\pi)$ i promieniu $R$ oblicz wkład całek krzywizny geodezyjnej i kątów zewnętrznych narożników do wzoru Gaussa--Bonneta. Wyjaśnij, dlaczego formuła dla gładkiego brzegu wymaga odpowiedniej poprawki.
\wskazowka{Przy brzegach z narożnikami Gauss--Bonnet zawiera sumę kątów zewnętrznych. Ustal orientację obwodu przed liczeniem wkładów łuków.}
\item \textbf{Zmiana metryki na powierzchni.} Niech $g'=e^{2u}g$ na zamkniętej orientowalnej powierzchni. Korzystając z transformacji krzywizny Gaussa, sprawdź, że całki $\int K_g\,dA_g$ i $\int K_{g'}\,dA_{g'}$ są równe. Zastosuj do płaskiego torusa z niekonstantnym $u$.
\item \textbf{Przeszkoda do niezerowej sekcji.} Niech $E\to M$ będzie zorientowaną wiązką parzystej rangi nad zamkniętą rozmaitością. Wyjaśnij, jakie ograniczenie na parowanie $\langle e(E),[M]\rangle$ wynika z istnienia globalnej sekcji nigdzie niezerowej, gdy ranga $E=\dim M$. Zastosuj do $TM$ dla $M=S^2\times S^2$.
\item \textbf{Euler a Pontriagin na rozkładzie.} Załóż, że zorientowana wiązka rzeczywista rzędu $4$ rozkłada się jako $E=L_1\oplus L_2$, gdzie $L_i$ to zorientowane wiązki płaszczyzn. Oblicz $e(E)$ i $p_1(E)$ w kategoriach $e(L_1)$ i $e(L_2)$, a następnie przyjmij $e(L_1)=2x$, $e(L_2)=3y$ dla odpowiednich klas stopnia $2$.
\item \textbf{Brak krzywizny a holonomia.} Na trywialnej linii zespolonej nad $S^1$ rozważ koneksję $d+ia\,d\theta$, $a\in\R$. Oblicz krzywiznę, klasę Cherna i holonomię wokół okręgu. Wyjaśnij, dlaczego niezerowa holonomia nie wymusza niezerowej klasy charakterystycznej.
\item \textbf{Pontriagin na sumie prostych zespolonych.} Dla $E=L_1\oplus L_2$ na $\C P^2$, gdzie $c_1(L_i)=a_i h$, oblicz $c_1(E)$, $c_2(E)$ oraz $p_1(E_{\R})$. Wyznacz wszystkie pary $(a_1,a_2)$ o $c_1(E)=0$ i $p_1(E_{\R})=8h^2$; osobno rozstrzygnij, czy wartość $-8h^2$ jest możliwa.
\wskazowka{Dla rzeczywistej wiązki pod zespoloną $E$ użyj $p_1(E_\R)=c_1(E)^2-2c_2(E)$. Po podstawieniu wyrażeń w $a_i$ sprowadź warunek do równania diofantycznego.}
\end{enumerate}


\chapter{Teoria Hodge'a i metoda Bochnera}
\begin{enumerate}
\item \textbf{Gwiazda na prostokątnym torusie.} Na $T^2$ weź metrykę $g=a^2dx^2+b^2dy^2$, gdzie $a,b>0$, i orientację $dx\wedge dy$. Oblicz $*1$, $*dx$, $*dy$ i $*(dx\wedge dy)$. Sprawdź bezpośrednio $*^2$ na formach stopni $0,1,2$.
\item \textbf{Gwiazda w trzech wymiarach.} Na euklidesowej $\R^3$ wyznacz $*(dx\wedge dy)$, $*(dy\wedge dz)$ i $*(dz\wedge dx)$. Oblicz $*\alpha$ dla $\alpha=x\,dx+y\,dy+z\,dz$ oraz $d(*\alpha)$.
\item \textbf{Współróżniczka w praktyce.} Dla $\alpha=x^2\,dy-xy\,dx$ na euklidesowej $\R^2$ oblicz $d\alpha$, $\delta\alpha$ i $\Delta\alpha=(d\delta+\delta d)\alpha$. Zweryfikuj wynik we współrzędnych.
\wskazowka{Współróżniczkę oblicz przez $\delta=(-1)^{n(k+1)+1}*d*$ w przyjętej konwencji; możesz też sprawdzić wynik laplasjanem współczynnikowym.}
\item \textbf{Harmoniczne funkcje wielomianowe.} Wyznacz wszystkie jednorodne wielomiany harmoniczne stopnia $2$ w dwóch zmiennych. Następnie wyjaśnij, dlaczego istnienie niezerowej funkcji harmonicznej na $\R^2$ nie przeczy twierdzeniu o harmonicznych funkcjach na zwartej rozmaitości.
\item \textbf{Harmoniczne $1$-formy na torusie.} Dla płaskiego $T^2$ znajdź bazę przestrzeni harmonicznych $1$-form, oblicz macierz ich iloczynów $L^2$ przy metryce z zadania 1 i podaj harmonicznego reprezentanta klasy $3[dx]-2[dy]$.
\item \textbf{Kohomologia i produkt.} Niech $M=S^2\times T^2$ z metryką produktową. Oblicz liczby Bettiego i wymiary przestrzeni form harmonicznych wszystkich stopni. Skonstruuj jawne generatory form harmonicznych korzystając z form na czynnikach.
\wskazowka{Dla metryki produktowej formy harmoniczne otrzymuje się przez iloczyny zewnętrzne ściągniętych form harmonicznych czynników.}
\item \textbf{Forma dokładna na rozmaitości niezwartej.} Na euklidesowej $\R$ sprawdź, że $dx$ jest jednocześnie dokładna, domknięta i harmoniczna. Wyjaśnij, które założenie klasycznego twierdzenia Hodge'a uniemożliwia przeniesienie wniosku o zerowaniu dokładnych form harmonicznych na ten przypadek.
\item \textbf{Konforemna zmiana gwiazdy.} Dla $g'=e^{2u}g$ na zorientowanej $n$-rozmaitości oblicz zależność między $*_{g'}$ i $*_g$ na $k$-formach. Sprawdź, że w wymiarze $n=2$ gwiazda na $1$-formach nie zależy od konforemnego czynnika.
\wskazowka{Policz, jak pod skalowaniem $g\mapsto e^{2u}g$ zmieniają się norma $k$-form i element objętości. Wykorzystaj definicję gwiazdy przez parowanie.}
\item \textbf{Hodge na $4$-torusie.} W standardowej metryce na $T^4$ znajdź bazy przestrzeni samodualnych i antysamodualnych harmonicznych $2$-form. Oblicz $b_2^+$, $b_2^-$ i sygnaturę formy przecięcia.
\wskazowka{Rozdziel bazowe $2$-formy na pary wymieniane przez $*$. Z ich sum i różnic uzyskasz podprzestrzenie własne gwiazdy.}
\item \textbf{Forma przecięcia w prostym przykładzie.} Dla $M=\mathbb{CP}^2\#(-\mathbb{CP}^2)$ oblicz macierz formy przecięcia na $H^2(M;\Z)$ i rozkład $H^2(M;\R)$ na podprzestrzenie dodatnią i ujemną. Wyznacz $b_2^\pm$ i sygnaturę.
\item \textbf{Bochner a topologiczna przeszkoda do metryki.} Rozważ $M=S^1\times S^2$ i jego $3$-krotne nakrycie wzdłuż czynnika $S^1$. Oblicz $b_1(M)$ i $b_1$ nakrycia. Zastosuj twierdzenie Bochnera ze skryptu, aby rozstrzygnąć, czy któraś z tych rozmaitości może mieć metrykę o dodatnio określonym tensorze Ricciego.
\wskazowka{Zastosuj tożsamość Bochnera dla harmonicznych $1$-form: dodatni Ricci wyklucza niezerowe formy harmoniczne. Porównaj z obliczonym $b_1$.}
\item \textbf{Bochner przy półdodatniej krzywiźnie.} Na produkcie $S^2\times S^1$ z metryką produktową rozważ formę $d\theta$ z czynnika okręgu. Oblicz $\nabla(d\theta)$, $\Delta(d\theta)$ i $\operatorname{Ric}^{\#}(d\theta)$. Wyjaśnij, dlaczego półdodatnia krzywizna Ricciego nie prowadzi do zaniku $H^1$.
\item \textbf{Holonomia i pierwszy Betti.} Załóż, że zamknięta rozmaitość ma $\operatorname{Ric}\ge0$ i że reprezentacja holonomii na przestrzeni stycznej w pewnym punkcie nie ma niezerowych wektorów ustalonych. Korzystając z rezultatów Bochnera ze skryptu, wyznacz $b_1(M)$. Podaj przykład, dla którego holonomia ma dokładnie jeden niezależny wektor ustalony.
\wskazowka{Przy $\Ric\ge0$ każda harmoniczna $1$-forma jest równoległa. W punkcie jej wartość musi być nieruchoma względem holonomii.}
\item \textbf{Weitzenböck w przestrzeni płaskiej.} Dla $\alpha=(x^2-y^2)\,dx+2xy\,dy$ na $\R^2$ porównaj $\Delta_H\alpha$ z laplasjanem koneksyjnym $\nabla^*\nabla\alpha$. Wyjaśnij, dlaczego człon Ricciego znika.
\item \textbf{Zerowe mody i indeks.} Niech $D:\R^4\to\R^3$ ma macierz $\begin{psmallmatrix}1&0&0&0\\0&1&0&0\\0&0&0&0\end{psmallmatrix}$. Oblicz $\ker D$, $\ker D^*$ i indeks $\dim\ker D-\dim\ker D^*$. Porównaj z superśladem operatora $\exp[-t\operatorname{diag}(D^*D,DD^*)]$.
\item \textbf{Sygnatura produktu powierzchni.} Dla $M=\Sigma_g\times\Sigma_h$ oblicz $b_2(M)$ i sygnaturę formy przecięcia, korzystając ze wzoru Künnetha. Wyznacz $b_2^\pm(M)$ i porównaj je z wynikami dla $T^4$.
\wskazowka{Zastosuj Künnetha do $H^2(\Sigma_g\times\Sigma_h)$; parowanie na składnikach mieszanych pozwala znaleźć pary dodatnie i ujemne.}
\item \textbf{Całkowitość liczby $\widehat A$.} Dla formalnych liczb charakterystycznych $\langle p_1^2,[M^8]\rangle=16a$ i $\langle p_2,[M^8]\rangle=b$ zapisz wzór $\widehat A(M)=(7\langle p_1^2\rangle-4\langle p_2\rangle)/5760$ i wyznacz warunek podzielności wymagany dla spinowej $8$-rozmaitości. Zastosuj do $(a,b)=(0,1440)$ i $(1,0)$, nie zakładając realizowalności obu par.
\wskazowka{Indeks operatora Diraca jest liczbą całkowitą. Po podstawieniu wzoru na $\widehat A$ otrzymasz warunek kongruencyjny, nie automatyczny dowód realizowalności.}
\item \textbf{Wirtualny wymiar moduli ASD.} Dla głównej wiązki $SU(2)$ o drugim numerze Cherna $k$ nad zamkniętą po prostu spójną $4$-rozmaitością użyj wzoru $\operatorname{ind}=8k-3(1-b_1+b_2^+)$ podanego w skrypcie. Oblicz indeks dla $S^4$, $\mathbb{CP}^2$ i $S^2\times S^2$ przy $k=1,2$, z właściwą orientacją.
\item \textbf{Indeks i orientacja czterowymiarowa.} Niech $M$ będzie zamkniętą orientowalną $4$-rozmaitością. Wyjaśnij, jak zamiana orientacji przestawia $b_2^+$ i $b_2^-$ oraz zmienia znak liczby $k=\langle c_2(P),[M]\rangle$ dla ustalonej wiązki. Oblicz konsekwencję dla formalnego wzoru indeksowego ASD i wskaż, dlaczego wniosek nie oznacza utożsamienia moduli ASD obu orientacji.
\wskazowka{Odwrócenie orientacji zamienia formy samodualne z antysamodualnymi. Liczba Chernowska zależy od orientacji klasy fundamentalnej.}
\end{enumerate}


\chapter{Grupy Liego i ich algebry}
\begin{enumerate}
\item \textbf{Grupa afiniczna prostej.} Niech $G=\{\begin{psmallmatrix}a&b\\0&1\end{psmallmatrix}:a>0,\ b\in\R\}$. Wyznacz prawo mnożenia, odwrotność, przestrzeń styczną w jedności oraz nawias w jej algebrze Liego. Sprawdź, czy grupa jest przemienna i czy jest spójna.
\item \textbf{Eksponenta macierzy trójkątnej.} Dla $X=\begin{psmallmatrix}a&b\\0&0\end{psmallmatrix}$ oblicz $\exp(tX)$, uwzględniając przypadek $a=0$. Porównaj otrzymaną rodzinę z podgrupami jednoparametrowymi grupy afinicznej z poprzedniego zadania.
\item \textbf{Grupa Heisenberga.} Dla $H=\{\begin{psmallmatrix}1&x&z\\0&1&y\\0&0&1\end{psmallmatrix}:x,y,z\in\R\}$ wyznacz działanie mnożenia we współrzędnych $(x,y,z)$, centrum i grupę komutatorów. Oblicz nawias w bazie $E_{12},E_{23},E_{13}$.
\item \textbf{Eksponenta nilpotentna.} Dla $X=aE_{12}+bE_{23}+cE_{13}\in\mathfrak h$ oblicz $\exp(X)$ i $\log(I+N)$ dla nilpotentnej macierzy $N$ odpowiedniego typu. Zapisz wzór Bakera--Campbella--Hausdorffa dla $X,Y\in\mathfrak h$.
\wskazowka{Dla ściśle górnotrójkątnych macierzy $3\times3$ zachodzi $X^3=0$. Utnij szereg wykładniczy i logarytmiczny po skończonej liczbie wyrazów.}
\item \textbf{Algebra $\mathfrak{sl}_2$.} Dla $H=\begin{psmallmatrix}1&0\\0&-1\end{psmallmatrix}$, $E=\begin{psmallmatrix}0&1\\0&0\end{psmallmatrix}$, $F=\begin{psmallmatrix}0&0\\1&0\end{psmallmatrix}$ oblicz wszystkie nawiasy. Wyznacz centrum $\mathfrak{sl}_2(\R)$ i podalgebry rozpięte przez wybrane pary generatorów.
\item \textbf{Działanie adjungowane.} W $SL(2,\R)$ weź $g=\operatorname{diag}(e^t,e^{-t})$. Oblicz $\operatorname{Ad}_g(H)$, $\operatorname{Ad}_g(E)$ i $\operatorname{Ad}_g(F)$. Wyznacz różniczkę w $t=0$ i porównaj ją z $\operatorname{ad}_H$.
\item \textbf{Forma Killinga w przykładzie.} Oblicz macierz $B(X,Y)=\operatorname{tr}(\operatorname{ad}_X\operatorname{ad}_Y)$ na $\mathfrak{sl}_2(\R)$ względem bazy $(H,E,F)$. Wyznacz sygnaturę i porównaj z algebrą $\mathfrak{so}(3)$.
\wskazowka{Oblicz macierze operatorów $\operatorname{ad}_H,\operatorname{ad}_E,\operatorname{ad}_F$ w zadanej bazie; następnie bierz ślady ich iloczynów.}
\item \textbf{Centralizator.} Dla $X=\operatorname{diag}(1,1,-2)\in\mathfrak{sl}_3(\R)$ wyznacz przestrzeń $\{Y\in\mathfrak{sl}_3(\R):[X,Y]=0\}$, jej wymiar i strukturę algebraiczną.
\item \textbf{Jednoparametrowa podgrupa torusa.} Na $T^2=\R^2/\Z^2$ rozważ $\gamma(t)=[t,\alpha t]$. Ustal, kiedy obraz $\gamma$ jest zamkniętym okręgiem, a kiedy jest gęsty w $T^2$. Wyznacz jądro homomorfizmu $\gamma:\R\to T^2$ w obu przypadkach.
\wskazowka{Krzywa $t\mapsto[t,\alpha t]$ jest okresowa, gdy istnieje $t\ne0$ z $t,\alpha t\in\Z$. Dla nieokresowości użyj kryterium gęstości podgrupy torusa.}
\item \textbf{Homomorfizm a różniczka.} Dla $\Phi:GL(n,\R)\to\R_{>0}$, $\Phi(A)=|\det A|$, oblicz $d\Phi_I$ i sprawdź, że jest homomorfizmem algebr Liego do abelowej algebry $\R$. Wyznacz jądro i porównaj z $\mathfrak{sl}_n$.
\item \textbf{Nakrycie dwóch grup obrotów.} Korzystając z mapy $SU(2)\to SO(3)$ omówionej w skrypcie, wyznacz obrazy $\pm I$, jądro homomorfizmu i obrazy podgrupy diagonalnej $\operatorname{diag}(e^{it},e^{-it})$. Oblicz, po jakim okresie zamyka się odpowiednia podgrupa w $SU(2)$ i w $SO(3)$.
\wskazowka{Porównaj elementy $\pm I$ w $SU(2)$ z ich działaniem przez sprzężenie na algebrze kwaternionów urojonych. Okres po przejściu do $SO(3)$ może się zmienić.}
\item \textbf{Algebra ruchów płaszczyzny.} Dla $SE(2)=\R^2\rtimes SO(2)$ znajdź bazę algebry Liego odpowiadającą dwóm translacjom i jednemu obrotowi. Oblicz nawiasy i działanie $\operatorname{Ad}$ obrotu na translacjach.
\item \textbf{Działania przez sprzężenie.} Niech $G=GL(2,\R)$ działa na $M_2(\R)$ przez $A\cdot X=AXA^{-1}$. Dla $X=\operatorname{diag}(1,2)$ znajdź stabilizator, przestrzeń styczną do orbity i wymiar orbity.
\item \textbf{Kiedy eksponenta nie jest surjektywna?} Wskaż macierz w $SL(2,\R)$, która nie jest eksponentą żadnej rzeczywistej macierzy z $\mathfrak{sl}_2(\R)$, i uzasadnij to za pomocą własności wartości własnych rzeczywistej macierzy $2\times2$. Porównaj z sytuacją w $SO(2)$.
\wskazowka{Jeśli macierz rzeczywista $2\times2$ ma sprzężone wartości własne, to jej eksponenta podlega ograniczeniom co do ujemnych wartości własnych. Zbadaj także przypadek nietrywialnego bloku Jordana.}
\item \textbf{Lewostronnie niezmiennicze pola.} W grupie Heisenberga we współrzędnych z zadania 3 wyznacz trzy lewostronnie niezmiennicze pola odpowiadające bazie w jedności. Oblicz ich nawiasy bezpośrednio z pochodnych współrzędnych.
\item \textbf{Macierzowe rozwiązanie równania różniczkowego.} Dla stałej macierzy $X\in\mathfrak{gl}_n(\R)$ rozwiąż równanie $\dot A(t)=A(t)X$ z $A(0)=A_0\in GL(n,\R)$. Zastosuj do rotacji płaszczyzny i do macierzy nilpotentnej $E_{12}$.
\item \textbf{Podgrupy normalne i ideały.} Dla grupy Heisenberga wyznacz co najmniej dwa nietrywialne ideały jej algebry Liego. Znajdź odpowiadające spójne podgrupy normalne i opisz otrzymane ilorazy.
\item \textbf{Baker--Campbell--Hausdorff w przypadku dwustopniowo nilpotentnym.} Dla macierzy $X=E_{12}$, $Y=E_{23}$ w $\mathfrak{gl}_3(\R)$ wyznacz $\exp(X)\exp(Y)$ oraz $\log(\exp(X)\exp(Y))$. Zweryfikuj wzór $X+Y+\frac12[X,Y]$ w tym przykładzie.
\wskazowka{Ponieważ $[X,Y]$ jest centralny, wszystkie nawiasy zagnieżdżone stopnia co najmniej trzy znikają. Szereg BCH kończy się na członie $\frac12[X,Y]$.}
\end{enumerate}


\chapter*{29a. Geometria Kleina: przestrzenie jednorodne i ich symetrie}
\phantomsection\addcontentsline{toc}{chapter}{29a. Geometria Kleina: przestrzenie jednorodne i ich symetrie}
\begin{enumerate}
\item \textbf{Przestrzeń punktów zorientowanych.} Grupa $SE(2)$ działa na parach $(p,v)$, gdzie $p\in\R^2$, $v\in S^1$ jest kierunkiem. Sprawdź przechodniość działania, znajdź stabilizator wybranej pary i przedstaw przestrzeń tych par jako $G/H$.
\item \textbf{Dwie grupy, ta sama przestrzeń.} Porównaj działania grup translacji $\R^2$ i grupy ruchów $SE(2)$ na płaszczyźnie. Oblicz stabilizatory zera, wymiary orbit i odpowiednie jednorodne opisy $\R^2\cong G/H$.
\item \textbf{Niezmiennicza metryka na ilorazie.} Niech $G$ działa przechodnio na $M=G/H$. Wyprowadź konieczność niezmienniczości iloczynu skalarnego na $T_oM$ względem reprezentacji izotropowej $H\to GL(T_oM)$. Zastosuj do standardowego działania $GL(n,\R)\ltimes\R^n$ na $\R^n$ i rozstrzygnij, czy może istnieć niezerowa niezmiennicza metryka Riemannowska.
\wskazowka{Metryka niezmiennicza na $G/H$ odpowiada dodatnio określonemu iloczynowi skalarnemu niezmienniczemu względem izotropii. Dla pełnej grupy liniowej sprawdź już działanie dylatacji.}
\item \textbf{Przestrzeń dodatnich macierzy.} Niech $GL(n,\R)$ działa na dodatnio określonych macierzach symetrycznych przez $g\cdot A=gAg^T$. Znajdź stabilizator $I$, udowodnij przechodniość i oblicz wymiar przestrzeni ilorazowej $GL(n,\R)/O(n)$.
\wskazowka{Zastosuj rozkład Cholesky'ego lub dodatni pierwiastek macierzy $A$. Stabilizator $I$ wyznacza równanie $gg^T=I$.}
\item \textbf{Grassmannian jako iloraz.} Dla działania $SO(4)$ na zorientowanych $2$-płaszczyznach w $\R^4$ opisz stabilizator płaszczyzny $\operatorname{span}(e_1,e_2)$ i wyznacz wymiar $SO(4)/(SO(2)\times SO(2))$.
\item \textbf{Działanie niezupełnie swobodne.} Grupa $C_m$ działa na $\C$ przez obroty $z\mapsto e^{2\pi i/m}z$. Opisz stabilizatory punktów, topologię ilorazu i lokalną strukturę orbifoldową w obrazie zera. Wyjaśnij, dlaczego sam homeomorfizm ilorazu z $\R^2$ nie usuwa informacji o grupie lokalnej.
\wskazowka{Stabilizatory są trywialne poza początkiem, lecz w zerze działa całe $C_m$. Rozróżnij topologię przestrzeni ilorazowej od danych orbifoldowych.}
\item \textbf{Iloraz torusa przez symetrię.} Na $T^2$ rozważ inwolucję $[x,y]\mapsto[-x,y]$. Wyznacz zbiór punktów stałych i opisz iloraz topologiczny jako powierzchnię z brzegiem. Porównaj jego strukturę z ilorazem przez inwolucję $[x,y]\mapsto[-x,-y]$.
\item \textbf{Działanie niewłaściwe.} Niech $\R$ działa na $T^2$ przez $t\cdot[x,y]=[x+t,y+\alpha t]$, gdzie $\alpha\notin\Q$. Wyznacz orbity i stabilizatory. Sprawdź, że działanie nie jest właściwe, analizując odwzorowanie $\R\times T^2\to T^2\times T^2$.
\wskazowka{Właściwe działanie powinno dawać przeciwobrazy zwartych podzbiorów w $T^2\times T^2$ będące zwartymi w $\R\times T^2$. Wykorzystaj gęstość orbity do zbudowania ciągu przeczącego temu warunkowi.}
\item \textbf{Dlaczego stabilizator musi być domknięty?} Rozważ grupę topologiczną $G=\R$ i jej niedomkniętą podgrupę $H=\Q$. Wykaż, że topologia ilorazowa na $G/H$ nie jest Hausdorffa. Porównaj z gęstą orbitą działania jednoparametrowego na $T^2$ z zadania 8, której stabilizator jest domknięty.
\wskazowka{W ilorazie $\R/\Q$ każdy niepusty otwarty zbiór obrazu ma gęste pełne przeciwobrazy; porównaj warunek Hausdorffa z separowaniem dwóch klas.}
\item \textbf{Hiperboliczna droga pionowa.} W półpłaszczyźnie Poincarégo z metryką $g=(dx^2+dy^2)/y^2$ oblicz długość pionowego odcinka od $(0,a)$ do $(0,b)$, gdzie $a,b>0$. Wyznacz działanie dylatacji $(x,y)\mapsto(cx,cy)$ na tę długość.
\item \textbf{Podgrupa paraboliczna.} Dla działania $PSL(2,\R)$ na półpłaszczyźnie oblicz działanie macierzy $\begin{psmallmatrix}1&t\\0&1\end{psmallmatrix}$. Znajdź jej punkty stałe na brzegu idealnym i wyznacz horocykle niezmiennicze przez to działanie.
\item \textbf{Geometria cusp w wymiarze dwa.} Rozważ $C=\{(x,y):y>1\}/(x\sim x+1)$ z metryką hiperboliczną. Oblicz pole obszaru $C$, długości przekrojów $y=\mathrm{const}$ i ich zachowanie przy $y\to\infty$.
\item \textbf{Torus: które automorfizmy są izometriami?} Dla płaskiego torusa $T^2=\R^2/\Z^2$ sprawdź, kiedy liniowe przekształcenie $A\in GL_2(\Z)$ indukuje izometrię standardowej metryki. Porównaj z warunkiem indukowania dyfeomorfizmu torusa.
\item \textbf{Butelka Kleina jako płaski iloraz.} Niech grupa $\Gamma$ jest generowana przez izometrie $a(x,y)=(x+1,y)$ i $b(x,y)=(-x,y+1)$ płaszczyzny. Sprawdź, że działa swobodnie, właściwie nieciągle i kokompaktowo. Wyznacz relację $bab^{-1}=a^{-1}$ i opisz fundamentalną dziedzinę ilorazu.
\wskazowka{Użyj grupy krystalograficznej z relacją $bab^{-1}=a^{-1}$. Dla properności i kokompaktowości wystarczy wskazać fundamentalny pasek i kontrolować skończoność przecięć przesuniętych zbiorów zwartych.}
\item \textbf{Stabilizator i krzywizna.} Załóż, że grupa izometrii działa przechodnio na spójnej powierzchni Riemannowskiej. Uzasadnij, że krzywizna Gaussa jest stała. Podaj przykład przechodniego działania grupy afinicznej, dla którego nie można zastosować tego argumentu.
\item \textbf{Niezmiennicza metryka na sferze.} Wykaż, że każdy $SO(3)$-niezmienniczy iloczyn skalarny na $T_pS^2$ jest proporcjonalny do standardowego. Wywnioskuj, że każda metryka Riemannowska na $S^2$ niezmiennicza względem standardowego działania $SO(3)$ jest okrągła z dokładnością do skali.
\wskazowka{Izotropia punktu na sferze zawiera obroty $SO(2)$ płaszczyzny stycznej. Iloczyn skalarny niezmienniczy pod wszystkimi obrotami musi być skalarną wielokrotnością euklidesowego.}
\item \textbf{Właściwość działania na przestrzeni zwartej.} Niech niezwarta grupa Liego $G$ działa na zwartej przestrzeni Hausdorffa $M$. Rozstrzygnij, czy takie działanie może być właściwe. Zastosuj do działania $\R$ na okręgu przez obroty.
\end{enumerate}


\chapter{Potok Ricciego przed osobliwościami}
\begin{enumerate}
\item \textbf{Potok na rozmaitości hiperbolicznej.} Niech $(M^n,g_0)$ będzie zwartą rozmaitością o stałej krzywiźnie sekcyjnej $-1$. Znajdź rozwiązanie potoku Ricciego postaci $g(t)=c(t)g_0$ i oblicz krzywiznę sekcyjną, skalarną oraz objętość w funkcji $t$.
\item \textbf{Metryka Einsteina o dodatniej stałej.} Jeżeli $\operatorname{Ric}(g_0)=\lambda g_0$, $\lambda>0$, oblicz maksymalny dodatni czas istnienia rozwiązania homotetycznego, tempo zaniku objętości oraz skalara $R(t)$. Wyraź odpowiedź przez $n,\lambda$ i $\operatorname{Vol}(M,g_0)$.
\item \textbf{Dwa czynniki o różnych znakach krzywizny.} Weź $M=S^2\times\Sigma_g$, $g\ge2$, z metrykami czynników o krzywiznach odpowiednio $+1$ i $-1$. Wyznacz zwykły potok Ricciego, oba czynniki skali i maksymalny przedział czasu. Oblicz zachowanie objętości w pobliżu końca istnienia.
\wskazowka{Dla produktu dwóch powierzchni potok zachowuje postać metryki produktowej, a każdy czynnik skali spełnia osobne równanie liniowe.}
\item \textbf{Dwie skale krzywizny.} Niech $M=M_1\times M_2$, gdzie $\operatorname{Ric}(g_i)=\lambda_i g_i$. Rozwiąż potok dla początkowej metryki $a_0g_1\oplus b_0g_2$ i ustal pierwszy czas osobliwy przy $\lambda_1>0$, $\lambda_2\le0$.
\item \textbf{Ewolucja całkowitej krzywizny skalarnej.} Dla zwartej rozmaitości bez brzegu wyprowadź wzór $\frac{d}{dt}\int_M R\,d\mu=\int_M(2|\operatorname{Ric}|^2-R^2)\,d\mu$ przy zwykłym potoku. Sprawdź, że w wymiarze dwa pochodna znika, i powiąż to ze wzorem Gaussa--Bonneta.
\wskazowka{Użyj $\partial_t d\mu=-R\,d\mu$ oraz równania ewolucji $\partial_tR=\Delta R+2|\Ric|^2$. Na zwartej rozmaitości scałkuj laplasjan przez części.}
\item \textbf{Pole powierzchni rodzaju trzy.} Zamknięta orientowalna powierzchnia rodzaju $3$ ma początkowe pole $A_0=20\pi$. Oblicz jego zmianę w zwykłym potoku Ricciego i ustal, czy globalny argument oparty tylko na polu wymusza skończony czas osobliwy.
\item \textbf{Tempo zmiany objętości.} Niech zamknięta $n$-rozmaitość ma w chwili $t=0$ skalar $R\ge c>0$. Zapisz $\left.\frac{d}{dt}\operatorname{Vol}(M,g(t))\right|_{t=0}$ i wyprowadź nierówność zależną od $c$ i objętości początkowej.
\item \textbf{Reakcja i dyfuzja skalara.} Na zwartej powierzchni rozważ $\partial_tR=\Delta R+R^2$. Zastosuj zasadę maksimum do oszacowania z dołu $\min_M R(t)$ przez rozwiązanie odpowiedniego ODE. Przeanalizuj oddzielnie ujemne i dodatnie minimum początkowe.
\wskazowka{W równaniu ewolucji skalara oddziel człon dyfuzyjny $\Delta R$ i reakcyjny $2|\Ric|^2$. Do nierówności punktowych wykorzystaj Cauchy'ego--Schwarza dla wartości własnych $\Ric$.}
\item \textbf{Skalowanie konkretnego rozwiązania.} Niech $g_0$ będzie metryką hiperboliczną w wymiarze $3$ z $K=-1$ i $g(t)=(1+4t)g_0$. Dla $t_0=1$, $Q=9$ oblicz $\widetilde g(s)=Q g(t_0+s/Q)$, jej krzywiznę skalarną i objętość względem objętości $g_0$. Wyznacz przedział nowego czasu $s$ wynikający z $t\ge0$.
\item \textbf{Osobliwość modelowa i blow-up.} Dla $g(t)=(1-4t)g_{S^3}$, $0\le t<1/4$, wybierz chwile $t_j\to1/4$ i skale $Q_j=(1-4t_j)^{-1}$. Oblicz metryki $Q_jg(t_j+s/Q_j)$ i ich przedziały istnienia w nowym czasie.
\wskazowka{Przy zmianie skali $g\mapsto Qg$ normy krzywizny skalują się przez $Q^{-1}$. Dopasuj zmianę zmiennej czasu tak, aby potok pozostał rozwiązaniem.}
\item \textbf{Zmiana zegara w potoku znormalizowanym.} Dla metryki Einsteina $\operatorname{Ric}(g_0)=\lambda g_0$ oblicz potok znormalizowany do stałej objętości. Wyznacz jawny związek czasu zwykłego i znormalizowanego, zgodnie z konwencją przyjętą w skrypcie.
\item \textbf{Metryka Bergera: dwa warunki początkowe.} Dla równań $A'=-8+4C/A$, $C'=-4C^2/A^2$ rozważ $(A(0),C(0))=(2,1)$ i $(1,3)$. Oblicz pochodne $A',C'$ i $\rho'=\frac d{dt}(C/A)$ w chwili zero. Wyjaśnij kierunek zmiany anizotropii.
\wskazowka{Metrykę Bergera zapisz w nieruchomej bazie lewostronnie niezmienniczych form; równania potoku dla współczynników otrzymasz z wartości własnych $\Ric$.}
\item \textbf{Wspólny kierunek w metrykach Bergera.} Użyj równań z poprzedniego zadania, aby pokazać, że dla każdej dodatniej pary $(A,C)$ iloraz $C/A$ nie może przekroczyć wartości $1$ od dołu ani przejść przez nią od góry w chwili regularnej ewolucji. Jaki wniosek daje jednoznaczność równań ODE?
\item \textbf{Skalowanie nieokrągłej metryki Bergera.} Dla metryki Bergera z parametrami $A=2$, $C=1$ oblicz składowe tensora Ricciego w bazie pól lewostronnie niezmienniczych. Porównaj otrzymane składowe i krzywiznę skalarną ze składowymi i skalarem dla metryki o parametrach $A=6$, $C=3$.
\item \textbf{Lokalna transformacja przez dyfeomorfizm.} Niech $\phi_t$ będzie przepływem pola $X=x\partial_x-y\partial_y$ na $\R^2$ i $g_0=dx^2+dy^2$. Oblicz $\left.\frac d{dt}\right|_{0}\phi_t^*g_0=\mathcal L_Xg_0$. Wyjaśnij rolę takiego członu w modyfikacji DeTurcka.
\item \textbf{Pochodna tensora Ricciego w rozwiązaniu Einsteina.} Dla $g(t)=c(t)g_0$ o $\operatorname{Ric}(g_0)=\lambda g_0$ oblicz $\partial_t\operatorname{Ric}(g(t))$ we współrzędnych stałych. Sprawdź zgodność z ewolucją tensora Ricciego opisaną w skrypcie.
\item \textbf{Laplacjan i stała zmiana skali.} Dla funkcji $f$ i $g'=cg$, $c>0$, oblicz zależność między $\Delta_{g'}f$ i $\Delta_gf$. Wykorzystaj ją do sprawdzenia zgodności wymiarów w równaniu ewolucji krzywizny skalarnej.
\item \textbf{Krzywizna dwóch iloczynów.} Niech $g=g_{S^2}\oplus g_{T^2}$, gdzie pierwsza metryka ma krzywiznę $+1$, a druga jest płaska. Oblicz wartości własne operatora Ricciego, $R$ i $|\operatorname{Ric}|^2$. Sprawdź wzór na $\partial_tR$ w chwili zero.
\item \textbf{Nieciągłość przedłużenia w chwili osobliwej.} Dla modelu $S^2\times\Sigma_g$ z zadania 3 oblicz normę krzywizny w czasie i uzasadnij, dlaczego nawet jeśli współczynniki metryki nie wybuchają, nie można przedłużyć potoku przez pierwszy czas osobliwy jako gładkiej metryki Riemannowskiej.
\wskazowka{Aby wykluczyć gładkie przedłużenie, znajdź skalarne niezmienniki krzywizny rozbieżne przy zbliżaniu do chwili osobliwej.}
\item \textbf{Oszacowanie przedłużenia.} Załóż, że rozwiązanie potoku na zwartej $M$ jest określone dla $0\le t<T<\infty$ i spełnia $\sup_{M\times[0,T)}|\operatorname{Rm}|\le K$. Sformułuj wniosek z kryterium Hamiltona, a następnie zastosuj kontrapozycję do każdego rozwiązania o rzeczywistym skończonym czasie maksymalnym.
\item \textbf{Soliton gradientowy.} Niech $g$ ma $\operatorname{Ric}(g)=\lambda g$. Sprawdź, że spełnia równanie solitonowe $\operatorname{Ric}+\nabla^2f=\lambda g$ dla stałej funkcji $f$. Podaj przykład zwężającego się, stacjonarnego i rozszerzającego się solitonu tego rodzaju.
\wskazowka{Soliton gradientowy spełnia równanie tensorowe zawierające $\Ric$ i $\nabla^2f$. Najpierw zbadaj jego ślad i skalowanie metryki.}
\item \textbf{Własności znormalizowanego potoku na powierzchni.} Na zamkniętej powierzchni rodzaju $g\ge2$ wykaż, że metryka o stałej krzywiźnie $K=-1$ jest punktem stałym znormalizowanego potoku. Zbadaj, co dzieje się z początkowym mnożnikiem skali $c>0$ metryki hiperbolicznej.
\end{enumerate}


\chapter*{30a. Hipoteza Poincarégo w wymiarze trzy: potok Ricciego z chirurgią}
\phantomsection\addcontentsline{toc}{chapter}{30a. Hipoteza Poincarégo w wymiarze trzy: potok Ricciego z chirurgią}
\begin{enumerate}
\item \textbf{Krzywizny w wymiarze trzy.} Niech tensor Ricciego w ortonormalnej bazie ma wartości własne $r_1,r_2,r_3$. Wyznacz krzywizny sekcyjne trzech płaszczyzn współrzędnych oraz skalar. Rozstrzygnij, czy dodatniość skalara oznacza dodatniość wszystkich krzywizn sekcyjnych.
\wskazowka{W wymiarze trzy operator krzywizny jest wyznaczony przez Ricciego. Skorzystaj z tożsamości $K_{ij}=r_i+r_j-R/2$ w bazie diagonalizującej $\Ric$.}
\item \textbf{Zmiana skali na walcu.} Dla metryki $g=r^2g_{S^2}+dz^2$ na $S^2\times\R$ oblicz krzywiznę skalarną, krzywizny sekcyjne i normę tensora Ricciego. Znajdź skalę, w której skalar w wybranym punkcie jest równy $1$.
\item \textbf{Blow-up osobliwości.} Załóż, że $Q_i=|\mathrm{Rm}|(x_i,t_i)=i^2\to\infty$, a kula o promieniu $i^{-1}$ przy $x_i$ ma objętość $v_i i^{-3}$. Dla $g_i(s)=Q_i g(t_i+s/Q_i)$ oblicz $|\mathrm{Rm}_{g_i}|(x_i,0)$ oraz promień i objętość tej kuli po przeskalowaniu. Co oznacza warunek $\inf_i v_i>0$?
\wskazowka{Pamiętaj, że przy $g'=Qg$ odległości mnożą się przez $\sqrt Q$, a objętości w wymiarze trzy przez $Q^{3/2}$. Normy krzywizny mnożą się przez $Q^{-1}$.}
\item \textbf{Płaski model zapadania.} Na $\R^2\times S^1_\varepsilon$ z okręgiem długości $2\pi\varepsilon$ oblicz asymptotykę objętości kuli promienia $1$ przy $\varepsilon\to0$. Porównaj ją z normą krzywizny i wyjaśnij, dlaczego do uzyskania trójwymiarowej granicy potrzeba kontroli objętości.
\item \textbf{Normalizacja entropii.} W definicji entropii $\mathcal W(g,f,\tau)$ ze skryptu przyjmij $f$ stałą na zamkniętej $3$-rozmaitości o objętości $V$. Wyznacz stałą $f$ wymuszoną przez normalizację całki do $1$, a następnie zapisz $\mathcal W$ za pomocą $V,\tau$ i średniej krzywizny skalarnej.
\wskazowka{Znormalizuj gęstość $(4\pi\tau)^{-3/2}e^{-f}d\mu$ do masy jeden. Dopiero potem podstaw stałą $f$ do funkcjonału.}
\item \textbf{Warunek równości w monotoniczności.} Załóż, że na gładkim przedziale potoku entropia Perelmana pozostaje stała. Wyprowadź z podanego w skrypcie wzoru warunek tensorowy $\operatorname{Ric}+\nabla^2 f=g/(2\tau)$. Sprawdź go dla odpowiednio przeskalowanej metryki na okrągłej sferze.
\wskazowka{W pochodnej entropii Perelmana występuje całka nieujemnego kwadratu tensora. Równość wymaga zaniku tego kwadratu w rozważanym rozwiązaniu.}
\item \textbf{Czapa asymptotycznie walcowa.} Rozważ metrykę obrotową $g=ds^2+\tanh^2(s)g_{S^2}$, $s\ge0$. Sprawdź warunek gładkości w biegunie $s=0$, oblicz krzywizny sekcyjne płaszczyzn radialnych i stycznych do sfer oraz opisz zachowanie metryki przy $s\to\infty$.
\wskazowka{Dla metryki $ds^2+f(s)^2g_{S^2}$ użyj wzorów $K_{\rm rad}=-f''/f$ i $K_{\rm tan}=(1-(f')^2)/f^2$; gładkość w czapie wymaga właściwego zachowania $f$ przy zerze.}
\item \textbf{Geometria przekroju szyjki.} Niech dokładna szyjka ma metrykę $g=h^2(2g_{S^2}+dz^2)$. Wyznacz pole centralnego przekroju, skalar i objętość odcinka $S^2\times[-L,L]$. Oblicz kombinacje bezwymiarowe uzyskane przez podzielenie pola przez $h^2$, a objętości przez $h^3$.
\item \textbf{Niezapadanie a naturalna skala.} Na cylindrze $S^2(h)\times\R$ porównaj przybliżone zachowanie $\operatorname{Vol}(B(x,r))/r^3$ w zakresach $r\ll h$ i $r\gg h$. Wyjaśnij, dlaczego sam spadek tego ilorazu dla dużych $r$ nie zaprzecza warunkowi $\kappa$-niezapadania przy ograniczeniu $|\mathrm{Rm}|\le r^{-2}$.
\wskazowka{Warunek $\kappa$-niezapadania stosuj tylko w skalach, dla których spełniona jest hipoteza ograniczenia krzywizny. Duże promienie na cienkim walcu nie spełniają jej automatycznie.}
\item \textbf{Topologia cięcia jednej szyjki.} Przyjmij, że centralny przekrój szyjki jest dwustronną sferą $S^2$ w zamkniętej orientowalnej $3$-rozmaitości $M$. Opisz dwie możliwości: sfera oddzielająca i nieoddzielająca. Porównaj wynik cięcia i wypełnienia przekrojów kulami z rozkładem $M$ na sumy spójne.
\item \textbf{Grupa podstawowa sumy spójnej.} Dla $M=L(5,1)\#(S^2\times S^1)\#L(7,1)$ oblicz $\pi_1(M)$ i jej abelianizację. Rozstrzygnij, czy założenia wersji twierdzenia o wygasaniu podanej w skrypcie są spełnione.
\item \textbf{Kontrola oszacowania na sferze okrągłej.} Na okrągłej $S^3$ o początkowej krzywiźnie sekcyjnej $1$ rozważ całkowicie geodezyjną sferę równikową $S^2$. Oblicz jej pole podczas zwykłego potoku Ricciego do chwili osobliwej, a następnie podstaw wartości $\operatorname{Scal}$ i $|A|^2$ do wzoru ze skryptu na pochodną pola minimalnej sfery. Sprawdź, czy zachodzi równość.
\wskazowka{Pod homotetycznym potokiem okrągłej $S^3$ promień zależy od czasu, a pole równika skaluje się jak kwadrat promienia. Porównaj pochodną z równaniem wariacyjnym.}
\item \textbf{Pętla grupowa a wygasanie.} Porównaj grupy podstawowe $S^3$, $S^2\times S^1$, $T^3$ i $\Sigma_g\times S^1$ dla $g\ge2$. Dla których modeli warunek grupowy twierdzenia Perelmana o wygasaniu ze skryptu jest spełniony? Nie zakładaj odwrotności tego twierdzenia.
\item \textbf{Model cięcia w $S^2\times S^1$.} Przetnij $S^2\times S^1$ wzdłuż przekroju $S^2\times\{p\}$. Opisz otrzymaną rozmaitość z brzegiem, a następnie dołącz kule $D^3$ do obu jej składowych brzegowych. Rozpoznaj otrzymaną zamkniętą $3$-rozmaitość i odtwórz $S^2\times S^1$ przez operację odwrotną.
\item \textbf{Rozpoznanie końcowej topologii.} Załóż, że zamknięta zorientowana $3$-rozmaitość jest połączoną sumą sferycznych form przestrzennych i kopii $S^2\times S^1$. Pokaż, że dodatkowy warunek $\pi_1(M)=1$ wymusza zniknięcie wszystkich nietrywialnych czynników, korzystając z van Kampena i twierdzenia o rozkładzie na czynniki pierwsze.
\end{enumerate}


\chapter*{30b. Geometryzacja Thurstona i geometrie trójwymiarowe}
\phantomsection\addcontentsline{toc}{chapter}{30b. Geometryzacja Thurstona i geometrie trójwymiarowe}
\begin{enumerate}
\item \textbf{Przykłady ośmiu geometrii.} Dla $S^3$, $T^3$, zamkniętej rozmaitości hiperbolicznej, $S^2\times S^1$, $\Sigma_2\times S^1$, niebanalnej wiązki okręgowej nad $T^2$, wiązki jednostkowych wektorów stycznych do $\Sigma_2$ oraz wiązki torusowej z hiperboliczną monodromią przyporządkuj jedną z ośmiu geometrii Thurstona. Dla każdego modelu wskaż wyróżniający niezmiennik.
\item \textbf{Przestrzeń soczewkowa.} Dla $L(7,2)$ wyznacz grupę podstawową, pierwszą homologię i rząd grupy działającej swobodnie na $S^3$. Sprawdź bezpośrednio, że działanie z definicji przestrzeni soczewkowej nie ma punktów stałych poza identycznością.
\item \textbf{Skończoność grupy podstawowej.} Niech $M=S^3/\Gamma$ będzie ilorazem przez skończoną grupę izometrii działającą swobodnie. Oblicz $\pi_1(M)$ i rozstrzygnij, czy $M$ może mieć taką samą grupę podstawową jak $T^3$ lub $S^2\times S^1$.
\item \textbf{Wiązki produktowe o dwóch rodzajach krzywizny.} Dla $M=\Sigma_g\times S^1$ z metryką produktową oblicz krzywizny sekcyjne płaszczyzn stycznych do powierzchni i płaszczyzn mieszanych. Porównaj $g=0$, $g=1$ oraz $g\ge2$.
\item \textbf{Krata w grupie Heisenberga.} Dla macierzy $\begin{psmallmatrix}1&x&z\\0&1&y\\0&0&1\end{psmallmatrix}$ z całkowitymi $x,y,z$ oblicz prawo mnożenia, komutator dwóch elementów i centrum grupy. Wykaż, że grupa jest nilpotentna klasy dwa, ale nie jest abelowa.
\item \textbf{Homologia nilrozmaitości.} Niech $M$ będzie nietrywialną orientowaną wiązką okręgową nad $T^2$ o liczbie Eulera $e=2$. Korzystając z ciągu Gysina, oblicz $H^1(M;\Z)$ i $H^2(M;\Z)$ oraz porównaj je z odpowiednimi grupami $T^3$.
\wskazowka{W ciągu Gysina mnożenie przez klasę Eulera $e=2$ jest kluczową strzałką między kohomologiami $T^2$. Oddziel część wolną i torsyjną.}
\item \textbf{Grupa i monodromia geometrii Sol.} Dla $A=\begin{psmallmatrix}3&2\\1&1\end{psmallmatrix}\in GL(2,\Z)$ oblicz wyznacznik, wartości własne i ich moduły. Rozstrzygnij, czy torusowa wiązka z monodromią $A$ spełnia warunek hiperbolicznej monodromii prowadzący do geometrii Sol.
\item \textbf{Abelianizacja grupy wiązki torusowej.} Dla $A=\begin{psmallmatrix}2&1\\1&1\end{psmallmatrix}\in SL(2,\Z)$ rozważ $G=\Z^2\rtimes_A\Z$. Oblicz $G_{\mathrm{ab}}$ i porównaj z $H_1$ odpowiadającej wiązki nad $S^1$. Powtórz dla $A=\begin{psmallmatrix}2&1\\1&2\end{psmallmatrix}$, wyjaśniając wcześniej, dlaczego druga macierz nie definiuje automorfizmu torusa.
\wskazowka{Abelianizacja $\Z^2\rtimes_A\Z$ ma postać $\Z\oplus\operatorname{coker}(A-I)$. Najpierw sprawdź jednak, czy $A\in GL(2,\Z)$.}
\item \textbf{Pole podstawowe orbifoldu.} Oblicz $\chi^{\mathrm{orb}}(S^2(2,2,5))$, $\chi^{\mathrm{orb}}(S^2(2,4,4))$, $\chi^{\mathrm{orb}}(S^2(2,4,5))$ i $\chi^{\mathrm{orb}}(T^2(2,2))$. W każdym przypadku wskaż znak i odpowiadający typ geometrii powierzchniowej, gdy orbifold jest dobry.
\wskazowka{Dla orbifoldu z punktami stożkowymi użyj $\chi^{\rm orb}=\chi(|B|)-\sum_i(1-1/m_i)$; znak nie rozstrzyga jakości orbifoldu bez dodatkowego kryterium.}
\item \textbf{Orbifold złej jakości.} Porównaj $S^2(2)$ i $S^2(2,2)$: oblicz charakterystyki orbifoldowe, sprawdź kryterium dobrego orbifoldu podane w skrypcie i wyjaśnij, dlaczego dodatnia charakterystyka sama nie rozstrzyga o istnieniu pokrycia przez powierzchnię.
\item \textbf{Lokalny model włókna wyjątkowego.} Dla działania $C_5$ na $D^2\times S^1$ danego przez $(z,u)\mapsto(e^{2\pi i/5}z,e^{4\pi i/5}u)$ sprawdź swobodność działania na przestrzeni całkowitej oraz stabilizator środka dysku w bazie. Oblicz rząd włókna wyjątkowego.
\item \textbf{Geometria Seiferta zależna od skręcenia.} Dla orientowanych wiązek okręgowych nad $S^2$, $T^2$ i $\Sigma_2$ przyjmij kolejno liczby Eulera $0$ i $1$. Korzystając z tabeli geometrii Seiferta, przyporządkuj sześć typów geometrii i wyjaśnij, dlaczego żadna nie jest typu $\mathrm{Sol}$.
\wskazowka{Dla wiązek Seiferta sprawdź dwie dane: znak orbifoldowej charakterystyki Eulera bazy oraz znikanie bądź nieznikanie skręcenia Eulera.}
\item \textbf{Wiązka styczna do powierzchni hiperbolicznej.} Oblicz liczbę Eulera wiązki $U\Sigma_g\to\Sigma_g$ dla $g=2,3$. Ustal geometrię przestrzeni całkowitej i porównaj z geometrią $\Sigma_g\times S^1$.
\item \textbf{Wykrywanie granicy między Nil i Euklidesem.} Dla zorientowanej wiązki okręgowej $M_e\to T^2$ opisz, jak zmienia się abelianizacja $\pi_1(M_e)$ przy $e=0$ i $e\ne0$. Pokaż na przykładzie $e=1,3$, że samo pierwsze Betti nie wyznacza liczby Eulera wiązki.
\item \textbf{Rozkład wzdłuż torusa.} Rozważ sklejanie dwóch zwartych rozmaitości trójwymiarowych wzdłuż składowych brzegowych będących torusami. Opisz, jakie informacje o wkładach do $\pi_1$ daje twierdzenie van Kampena i dlaczego grupa fundamentalna rezultatu zależy od wybranego dyfeomorfizmu sklejającego.
\wskazowka{Dla sklejania wzdłuż nieściśliwego torusa skorzystaj z iloczynu amalgamowanego grup podstawowych. Homomorfizmy włożenia zależą od mapy sklejenia.}
\item \textbf{Pierwszość i nierozkładalność.} Zbadaj $S^2\times S^1$, $\mathbb{RP}^3\#\mathbb{RP}^3$ i $T^3$ pod kątem pierwszości i nierozkładalności. Wyjaśnij, dlaczego pojęć tych nie należy utożsamiać w dowodzie geometryzacji.
\item \textbf{Torus ściskalny.} Niech $T=\partial(S^1\times D^2)$ będzie torusem brzegowym pełnego torusa. Oblicz jądro homomorfizmu $\pi_1(T)\to\pi_1(S^1\times D^2)$ i podaj krzywą prostą reprezentującą jego generator. Porównaj z torusem nieściśliwym w rozkładzie JSJ.
\item \textbf{Przeszkoda do geometrii hiperbolicznej.} Wyjaśnij, dlaczego zamknięta rozmaitość $T^3$ nie może być zamkniętą rozmaitością hiperboliczną, porównując grupy podstawowe i ich podgrupy abelowe. Przeanalizuj podobnie $\Sigma_g\times S^1$ dla $g\ge2$.
\item \textbf{Od potoku do rozkładu geometrycznego.} Przyjmij jako dane twierdzenie geometryzacyjne i opis długoczasowego potoku z chirurgią ze skryptu. Dla rozmaitości będącej sumą spójną części sferycznej, torusowej wiązki typu Sol i części hiperbolicznej wskaż, które podziały odpowiadają sferom, a które torusom JSJ, oraz jakich geometrii można oczekiwać po odpowiednich rozcięciach.
\wskazowka{Najpierw rozdziel rozkład na sumę spójną (cięcia wzdłuż sfer) od rozkładu JSJ (cięcia wzdłuż torusów nieściśliwych). Dopiero potem przypisz geometrię elementom.}
\end{enumerate}

\chapter{Formy przecięcia i uchwyty w wymiarze cztery}
\begin{enumerate}
\item \textbf{Jednoparametrowa rodzina form.} Dla $m\in\Z$ rozważ $A_m=\begin{psmallmatrix}m&1\\1&0\end{psmallmatrix}$. Wyznacz wyznacznik, sygnaturę i parzystość odpowiadającej formy. Dla parzystego $m$ podaj całkowitą zmianę bazy sprowadzającą ją do formy hiperbolicznej $H$. Co zmienia się dla nieparzystego $m$?
\item \textbf{Forma trójdiagonalna.} Zbadaj formę o macierzy $A=\begin{psmallmatrix}2&-1&0\\-1&2&-1\\0&-1&2\end{psmallmatrix}$: oblicz jej wyznacznik, sygnaturę oraz parzystość. Czy może być formą przecięcia zamkniętej, jednospójnej, zorientowanej $4$-rozmaitości? Uzasadnij odpowiedź.
\wskazowka{Wyznacznik i parzystość macierzy rozpatruj osobno od sygnatury; do postaci diagonalnej nad $\R$ można użyć kryterium Sylvestera, ale do klasyfikacji całkowitej wymagane są macierze unimodularne.}
\item \textbf{Samoprzecięcia i zmiana bazy.} Niech $Q$ ma macierz $\begin{psmallmatrix}1&2\\2&0\end{psmallmatrix}$ w bazie $(u,v)$. Oblicz kwadraty wektorów $u+kv$ dla $k=-2,-1,0,1,2$ oraz macierz formy w bazie $(u+v,v)$. Sprawdź niezmienniczość wyznacznika i sygnatury.
\item \textbf{Forma przecięcia sum spójnych.} Dla $X=3\C P^2\#2\overline{\C P^2}$ oblicz $b_2(X)$, $\sigma(X)$, $\chi(X)$, macierz $Q_X$ w naturalnej bazie oraz parzystość. Zbadaj, czy $X$ jest spinowa.
\item \textbf{Dwie powierzchnie produktowe.} W $\Sigma_g\times\Sigma_h$ rozważ klasy $[\Sigma_g\times\{q\}]$ i $[\{p\}\times\Sigma_h]$. Oblicz ich samoprzecięcia i wzajemne przecięcie. Jaką formę otrzymujemy na podgrupie, którą generują?
\item \textbf{Pierwiastki i odbicia.} W kracie $E_8$ wybierz wektor bazowy $v$ o $Q(v,v)=2$. Wykaż, że odwzorowanie $s_v(x)=x-Q(x,v)v$ zachowuje kratę i formę. Oblicz $s_v(v)$ i działanie na dwóch wektorach bazowych sąsiadujących z $v$ w grafie $E_8$.
\wskazowka{Odbicie względem wektora o odpowiednim kwadracie ma postać $s_v(x)=x-2Q(x,v)Q(v,v)^{-1}v$. Sprawdź, kiedy zachowuje kratę całkowitą.}
\item \textbf{Dodanie płaszczyzny hiperbolicznej.} Wyznacz rząd, wyznacznik, sygnaturę i parzystość form $E_8\oplus H$, $-E_8\oplus2H$ oraz $E_8\oplus(-E_8)$. Które z nich są określone dodatnio, a które nieokreślone?
\wskazowka{Dodanie płaszczyzny hiperbolicznej zmienia rangę i parę dodatnich/ujemnych kierunków, ale nie sygnaturę. Oddziel niezmienniki stabilne od bezwzględnych.}
\item \textbf{Charakterystyczne wektory modulo dwa.} Dla unimodularnej macierzy symetrycznej $A=\begin{psmallmatrix}1&1\\1&0\end{psmallmatrix}$ znajdź wszystkie klasy $k\in(\Z/2)^2$ spełniające $k^TAx\equiv x^TAx\pmod2$ dla każdego $x\in\Z^2$. Zinterpretuj wynik jako algebraiczny odpowiednik klasy Wu.
\wskazowka{Wektor charakterystyczny $k$ spełnia $Q(x,x)\equiv Q(x,k)\pmod2$ dla każdego $x$ w kracie. Wystarczy sprawdzić warunek na bazie.}
\item \textbf{Obramowany splot dwuskładowy.} Przyjmij, że diagram uchwytów $2$ ma macierz sprzężeń $A=\begin{psmallmatrix}2&1\\1&3\end{psmallmatrix}$. Wyznacz formę przecięcia śladu, jej wyznacznik i sygnaturę, a następnie oblicz $H_1$ trójwymiarowego brzegu jako $\operatorname{coker}A$.
\item \textbf{Plumbowanie liniowe.} Niech trzy sfery w plumbowaniu mają kwadraty $-2,-3,-2$, a sąsiednie przecinają się raz. Zapisz macierz formy, sprawdź jej określoność i oblicz grupę $H_1$ brzegu, posługując się postacią Smitha.
\wskazowka{Macierz plumbowania odczytuj z samoprzecięć sfer i ich geometrycznych przecięć. Homologię brzegu wiąże z nią dokładny ciąg pary.}
\item \textbf{Przesunięcie uchwytu.} Dla $A=\begin{psmallmatrix}0&1\\1&0\end{psmallmatrix}$ wykonaj zmianę bazy $u'=u+2v$, $v'=v$. Oblicz $P^TAP$ i wyjaśnij, dlaczego odpowiadające przesunięcie uchwytu nie zmienia izometrii formy mimo zmiany macierzy.
\wskazowka{Przesunięcie uchwytu odpowiada całkowitoliczbowej zmianie bazy $Q\mapsto P^TQP$. Pilnuj, by $\det P=\pm1$.}
\item \textbf{Brzeg przy osobliwej macierzy.} Rozważ $A=\begin{psmallmatrix}0&0\\0&4\end{psmallmatrix}$. Oblicz $\ker A$ i $\operatorname{coker}A$ nad $\Z$. Określ część wolną i torsyjną pierwszej homologii brzegu odpowiedniego śladu $2$-uchwytów.
\item \textbf{Parzystość a zmiana bazy.} Sprawdź, czy forma $H\oplus[1]$ może być izometryczna nad $\Z$ z $H\oplus[-1]$. Porównaj sygnatury, a następnie wyjaśnij, dlaczego samo porównanie parzystości nie rozstrzyga tego przykładu.
\item \textbf{Liczba geometrycznych przecięć.} Dwie zorientowane sfery przecinają się w sześciu punktach o znakach $+,+,+,-,-,-$. Co można wywnioskować o ich liczbie przecięcia? Jakie dodatkowe informacje o pętlach Whitneya i obramowaniach są potrzebne, by próbować usunąć punkty przecięcia w wymiarze cztery?
\wskazowka{Algebraiczna liczba przecięć jest sumą znaków; geometryczna jest liczbą punktów. Para punktów o przeciwnych znakach nie zawsze daje się usunąć bez dodatkowych warunków.}
\item \textbf{Dwie formy o tej samej sygnaturze.} Porównaj formy $2H$ oraz $I_{2,2}=\operatorname{diag}(1,1,-1,-1)$. Wykaż, że nie są izometryczne nad $\Z$, chociaż mają tę samą rangę, wyznacznik i sygnaturę.
\item \textbf{Charakterystyka Eulera a forma.} Niech $X$ będzie zamkniętą, spójną, jednospójną, zorientowaną topologiczną $4$-rozmaitością i $b_2(X)=r$. Korzystając z dualności Poincarégo, wyznacz wszystkie jej liczby Bettiego i $\chi(X)$. Zastosuj wynik do form rang $0$, $2$ i $8$.
\item \textbf{Obramowanie a sygnatura.} Dla rodziny $A_t=\begin{psmallmatrix}t&1\\1&-2\end{psmallmatrix}$, $t\in\Z$, wyznacz wartości $t$, dla których forma jest dodatnio określona, ujemnie określona albo nieokreślona. Dla których $t$ jest unimodularna?
\item \textbf{Konsekwencja dla gładkości.} Korzystając z twierdzenia Donaldsona (rozdział 34), oceń, które z form $E_8$, $E_8\oplus E_8$, $H\oplus E_8$ podlegają bezpośrednio przeszkodzie diagonalizacji formy określonej. Wyraźnie odróżnij wniosek negatywny od przypadku, którego to twierdzenie nie rozstrzyga.
\wskazowka{Rozpatrz łącznie określoność formy, jej parzystość i unimodularność. Twierdzenia Freedmana i Donaldsona mają różne hipotezy oraz kategorie wniosku.}
\item \textbf{Homologia brzegu a wyznacznik.} Przyjmij, że ślad $2$-uchwytów ma nieosobliwą macierz sprzężeń $A$ rozmiaru $r$. Wyznacz rząd skończonej grupy $H_1(\partial W;\Z)$ w dwóch nowych przykładach: $A=[7]$ i $A=\begin{psmallmatrix}0&3\\3&0\end{psmallmatrix}$. Porównaj jej strukturę w obu przypadkach.
\wskazowka{Dla uchwytowej rozmaitości czterowymiarowej pierwszą homologię brzegu w wielu przykładach opisuje cokernel macierzy łączenia. Sprawdź pełny dokładny ciąg pary.}
\end{enumerate}

\chapter{Uchwyty Cassona i zanurzanie dysków}
\begin{enumerate}
\item \textbf{Drzewo o nieregularnym rozgałęzieniu.} Na pierwszym piętrze narysuj uchwyt z trzema załamaniami o znakach $+,+,-$. W drugim piętrze przypisz odpowiednim trzem uchwytom kolejno $1$, $2$ i $3$ załamania, a każdemu uchwytowi trzeciego piętra po dwa. Podaj liczby załamań $r_1,r_2,r_3$ i rangi grup $\pi_1(T_1),\pi_1(T_2),\pi_1(T_3)$.
\item \textbf{Regularna wieża.} Załóż, że każdy uchwyt na każdym piętrze ma dokładnie trzy punkty podwójne. Oblicz liczbę uchwytów w piętrze $n$, liczbę wszystkich uchwytów w $T_n$ oraz rangę $\pi_1(T_n)$. Porównaj tempo wzrostu liczby uchwytów z zachowaniem grupy podstawowej w przejściu do granicy.
\wskazowka{Uchwyt Cassona koduje się nieskończonym drzewem zakorzenionym, a liczba wierzchołków na kolejnych poziomach zależy od rozgałęzienia. Liczby skończonych pięter nie są jeszcze typem granicy.}
\item \textbf{Różne znaki, te same rangi.} Skonstruuj dwa różne drzewa znakowane o jednakowej liczbie gałęzi w każdym wierzchołku, ale różnych znakach. Wyjaśnij, jaką informację o wieżach wykrywają rangi $\pi_1(T_n)$, a jakiej informacji o plumbowaniach nie wykrywają.
\item \textbf{Etykiety przecięć.} Niech $\pi=\langle t\rangle\cong\Z$. Dla czterech punktów przecięcia o danych $(\varepsilon,g)=(+,1),(+,t),(-,1),(-,t)$ oblicz $\lambda(A,B)$. Powtórz rachunek po zastąpieniu ostatniej etykiety przez $t^2$. W którym przypadku same dane algebraiczne dopuszczają sparowanie wszystkich punktów?
\wskazowka{Etykieta przecięcia składa się ze znaku i elementu grupy podstawowej. Rozdziel tę informację od zwykłej sumy algebraicznych znaków.}
\item \textbf{Zerowa zwykła liczba przecięcia to za mało.} Niech $\pi=\Z$ i $\lambda(A,B)=1-t\in\Z[t,t^{-1}]$. Oblicz augmentację tego elementu. Dlaczego z jej znikania nie wolno wnioskować, że da się znaleźć parę punktów z nullhomotopijnym okręgiem Whitneya?
\wskazowka{Rozważ dwa punkty o przeciwnych znakach, ale o różnych etykietach grupowych. W liczbie całkowitej zniosą się, lecz nie muszą tworzyć dopuszczalnej pary Whitneya.}
\item \textbf{Liczba samoprzecięcia i jej redukcja.} W $\Z[t,t^{-1}]/\langle t-t^{-1}\rangle$ porównaj klasy $\mu_1=2+t+t^{-1}$, $\mu_2=t+t^{-1}$ oraz $\mu_3=2+t-t^{-1}$. Wskaż, które stają się równe po dodatkowym ilorazowaniu przez $\Z\cdot1$.
\item \textbf{Zamiana dwóch płatów.} Dla jednego punktu podwójnego z etykietą $g$ opisz zmianę etykiety po zamianie kolejności płatów. Oblicz odpowiednie klasy w grupie samoprzecięć dla $\pi=\Z/3$ i porównaj je z klasami w pierścieniu grupowym bez ilorazowania.
\item \textbf{Macierz dualności.} Dwie immersyjne sfery $B_1,B_2$ i dwa dyski $A_1,A_2$ mają macierz $\bigl(\lambda(A_i,B_j)\bigr)=\begin{psmallmatrix}1&1-t\\0&1\end{psmallmatrix}$ nad $\Z[t,t^{-1}]$. Sprawdź, czy jest odwracalna. Znajdź nowy układ formalnych generatorów, dla którego macierz parowania staje się jednostkowa, i wyjaśnij, dlaczego nie jest to jeszcze geometryczny dowód zanurzania dysków.
\wskazowka{Wyznacz macierz parowań między sferami i dyskami dualnymi. Do zastosowania twierdzenia o dyskach wymagany jest odpowiedni układ dualności, nie sama zgodność rang.}
\item \textbf{Sprawdzenie hipotezy grupowej.} Zbuduj ciągi podnormalne o ilorazach skończonych lub nieskończenie cyklicznych dla grup $\Z^2$, $\Z/6$ i $\Z\rtimes(\Z/2)$, gdzie element nietrywialny działa przez $n\mapsto-n$. Które z nich obejmuje cytowana w skrypcie wersja twierdzenia Freedmana--Quinna?
\wskazowka{Sprawdź dokładnie, w której grupie podstawowej należy testować hipotezę o „dobrej” grupie. Nie należy przenosić własności automatycznie między otoczeniem dysków i całą rozmaitością.}
\item \textbf{Wniosek z twierdzenia o dyskach.} W rozmaitości $M^4$ z $\pi_1(M)=\Z^2$ dane są trzy właściwie immersyjne obramowane dyski i trzy sfery spełniające $\lambda(A_i,B_j)=\delta_{ij}$, $\lambda(B_i,B_j)=0$ i $\mu(B_j)=0$. Sformułuj dokładnie, jaki rodzaj dysków zapewnia twierdzenie o zanurzaniu, oraz czego nie orzeka o gładkości i regularnej homotopii wyjściowych immersji.
\item \textbf{Nieuprawniona implikacja.} Załóż, że dla nieskończonej wieży $CH$ stwierdzono tylko $\pi_1(CH)=0$. Czy wynika stąd homeomorfizm pary $(CH,\partial_-CH)$ ze standardowym otwartym $2$-uchwytem? Wskaż, jakiej dodatkowej informacji topologicznej dostarcza twierdzenie Freedmana.
\item \textbf{Obramowanie dysku Whitneya.} Wyjaśnij na przykładzie pary punktów przecięcia, dlaczego zerowa liczba przecięcia w $\Z[\pi]$ nie wystarcza do wykonania gładkiego ruchu Whitneya w wymiarze cztery. Osobno uwzględnij możliwość przecinania wnętrza dysku i problem jego obramowania.
\wskazowka{Przy dysku Whitneya zbadaj trywialność jego wiązki normalnej względnie do obramowania brzegu; sam fakt istnienia topologicznego dysku nie rozstrzyga framacji.}
\item \textbf{Meridian i sfera dualna.} Rozważ obramowaną, lokalnie płasko zanurzoną sferę $S\subset M^4$ i sferę dualną $B$ przecinającą $S$ geometrycznie w jednym punkcie. Zastosuj opis chirurgii ze skryptu do dwóch niezależnych takich par, aby określić wpływ kolejnych operacji na generatory grupy podstawowej dopełnienia.
\item \textbf{Dyski z relacjami grupowymi.} Niech $\pi=\langle a,b\mid ab=ba\rangle$. Dla punktów z etykietami $a,b,ab,a^{-1}b^{-1}$ i znakami odpowiednio $+,+,-,-$ oblicz element $\lambda\in\Z[\pi]$. Czy zniknięcie augmentacji gwarantuje, że $\lambda=0$?
\item \textbf{Zależność od wybranych ścieżek.} Wskazanie punktu bazowego i ścieżek do dwóch powierzchni może zmieniać etykiety punktów przecięcia przez mnożenie z lewej i prawej strony elementami grupy. Zbadaj algebraicznie, jak przechodzi $\lambda=1+t$ w $\Z[t,t^{-1}]$ przy równoczesnym pomnożeniu wszystkich etykiet przez $t^2$.
\wskazowka{Zmiana ścieżek bazowych sprzęga lub przemnaża etykiety przecięć elementami $\pi_1$. Oddziel element zależny od wyboru od informacji niezmienniczej.}
\item \textbf{Topologiczne a gładkie przedłużenie.} Załóż, że dane twierdzenie daje lokalnie płaski dysk topologiczny z zadanym obramowanym brzegiem. Wyjaśnij, dlaczego nie wolno bez dodatkowych hipotez uważać tego dysku za gładki i zastosować do niego zwykłego twierdzenia o otoczeniu tubularnym z kategorii DIFF.
\item \textbf{Skończone piętra i granica.} Dla wieży z $r_1=2$, $r_2=5$, $r_3=7$ oblicz grupy podstawowe trzech skończonych etapów oraz obrazy homomorfizmów między nimi. Czy którykolwiek z tych skończonych etapów musi być jednospójny?
\wskazowka{Każde skończone piętro jest tylko etapem zanurzonych dysków. Topologiczne osadzenie uzyskuje się dopiero po zastosowaniu twierdzenia do odpowiedniej nieskończonej wieży.}
\item \textbf{Granice zastosowalności.} Porównaj grupy $\pi_1(M)=1$, $\pi_1(M)=\Z^3$ i $\pi_1(M)=F_2$ pod względem możliwości bezpośredniego zastosowania wersji twierdzenia o dyskach podanej w rozdziale. W przypadku, w którym założenia twierdzenia nie są spełnione, wyjaśnij, dlaczego samo to nie stanowi kontrprzykładu do istnienia zanurzonych dysków.
\end{enumerate}

\chapter{Klasyfikacja topologicznych rozmaitości czterowymiarowych}
\begin{enumerate}
\item \textbf{Dopuszczalne pary niezmienników.} Dla form $H$, $E_8$, $2E_8$, $E_8\oplus H$ i $[1]\oplus[-1]$ oblicz sygnaturę, parzystość i wszystkie wartości $k\in\Z/2$, dla których twierdzenie Freedmana zapewnia zamkniętą jednospójną topologiczną $4$-rozmaitość o danych $(Q,k)$.
\wskazowka{Dla form unimodularnych sprawdź oddzielnie parzystość, określoność i dopuszczalny niezmiennik Kirby'ego--Siebenmanna. Nie każda formalna para danych jest realizowalna.}
\item \textbf{Dwie rozmaitości o tej samej formie.} Dla nieparzystej unimodularnej formy $Q=\operatorname{diag}(1,1,-1)$ skonstruuj na poziomie twierdzenia klasyfikacyjnego dwa różne zorientowane typy topologiczne. Podaj niezmiennik, który je rozróżnia.
\item \textbf{Orientacja i klasyfikacja.} Wyznacz, co dzieje się z $Q_X$, sygnaturą oraz $\operatorname{ks}(X)$ po odwróceniu orientacji $X$. Zastosuj do przypadku $Q_X=E_8\oplus[-1]$.
\wskazowka{Odwrócenie orientacji zastępuje $Q$ przez $-Q$ i odpowiednio zmienia sygnaturę. Niezmiennik $\operatorname{ks}$ ma wartości w $\Z/2$.}
\item \textbf{Dlaczego wymagana jest unimodularność?} Rozważ $Q=[3]$ oraz $Q=\begin{psmallmatrix}0&2\\2&0\end{psmallmatrix}$. Oblicz ich wyznaczniki i wyjaśnij, dlaczego nie są dopuszczalnymi formami przecięcia zamkniętych jednospójnych $4$-rozmaitości.
\item \textbf{Ranga i charakterystyka Eulera.} Dla rozmaitości objętych twierdzeniem Freedmana oblicz $\chi(X)$ na podstawie rangi $Q_X$. Porównaj przypadki $Q=0$, $H\oplus H$ oraz $E_8\oplus H$.
\item \textbf{Niezmiennik Kirby'ego--Siebenmanna sum spójnych.} Niech $X$ i $Y$ będą jednospójnymi topologicznymi $4$-rozmaitościami o niezmiennikach $(E_8,1)$ oraz $([1],1)$. Oblicz $(Q,\operatorname{ks})$ dla $X\#Y$, a następnie dla $X\#Y\#Y$. Zbadaj warunek parzystości nowej formy.
\wskazowka{Przy sumie spójnej formy przecięcia sumują się ortogonalnie, a niezmiennik Kirby'ego--Siebenmanna jest addytywny modulo dwa.}
\item \textbf{Parzysta forma z zerową przeszkodą stabilną.} Dla $Q=E_8\oplus E_8\oplus H$ oblicz wymuszoną wartość $\operatorname{ks}$ i wskaż, co można stwierdzić o istnieniu modelu topologicznego. Czy samo zerowanie $\operatorname{ks}$ daje w wymiarze cztery strukturę gładką?
\wskazowka{Znikanie stabilnej przeszkody Kirby'ego--Siebenmanna nie gwarantuje istnienia gładkiej struktury w wymiarze cztery. Porównaj także określone formy z Donaldsonem.}
\item \textbf{Porównanie dwóch sposobów odróżniania.} Dwie zorientowane jednospójne topologiczne $4$-rozmaitości mają izometryczne nieparzyste formy przecięcia, lecz różne $\operatorname{ks}$. Czy są homeomorficzne z zachowaniem orientacji? Co zmieniłoby się, gdyby znać jedynie równość sygnatur i liczb Bettiego?
\item \textbf{Macierz w niestandardowej bazie.} Niech $Q$ ma macierz $A=\begin{psmallmatrix}1&1\\1&0\end{psmallmatrix}$. Znajdź całkowitą bazę, w której ma postać $\operatorname{diag}(1,-1)$. Zastosuj twierdzenie Freedmana do określenia liczby dopuszczalnych typów topologicznych dla tej formy.
\wskazowka{Dla dwóch macierzy tej samej formy szukaj zmiany bazy z $GL_n(\Z)$, a nie tylko izometrii rzeczywistej. Wymagany warunek to $P^TQP=Q'$.}
\item \textbf{Forma z torsją?} Wyjaśnij, dlaczego dla zamkniętej jednospójnej $4$-rozmaitości $H_2(X;\Z)$ nie ma torsji, mimo że w ogólnej definicji formy przecięcia używa się ilorazu przez $\operatorname{Tor}$. Wskaż dokładnie miejsce zastosowania założenia o grupie podstawowej.
\item \textbf{Obliczenie klasy $w_2$ z formy.} Dla $Q=H\oplus[1]$ wyznacz funkcjonał $x\mapsto Q(x,x)\pmod2$. Czy odpowiadająca mu jednospójna $4$-rozmaitość może mieć strukturę spinową? Porównaj z $2H$.
\wskazowka{Zastosuj formułę Wu: klasa $w_2$ jest elementem charakterystycznym redukcji formy przecięcia modulo dwa.}
\item \textbf{Ograniczenie do kategorii TOP.} Wyjaśnij, dlaczego twierdzenie Freedmana pozwala porównywać topologicznie rozmaitości o zadanych formach i $\operatorname{ks}$, ale nie przesądza o liczbie struktur gładkich. Podaj przykład z tego i następnego rozdziału, ilustrujący istotność rozróżnienia.
\item \textbf{Brak informacji o grupie podstawowej.} Porównaj $S^4$ oraz $S^1\times S^3$ pod względem rangi $H_2$, charakterystyki Eulera i grupy podstawowej. Wskaż, dlaczego nie wolno zastosować twierdzenia klasyfikacyjnego do stwierdzenia ich homeomorficzności.
\item \textbf{Niespójność danych.} Student podaje parę $Q=E_8$, $k=0$, a następnie parę $Q=H$, $k=1$, jako niezmienniki zamkniętych jednospójnych $4$-rozmaitości. Zidentyfikuj osobno wszystkie naruszone warunki klasyfikacji.
\wskazowka{Dane formy i $\operatorname{ks}$ podlegają relacjom zgodności; sprawdź ograniczenia sygnaturowe w przypadku form parzystych.}
\item \textbf{Podstawowe modele.} Wyznacz $Q$, sygnaturę, parzystość i $\operatorname{ks}$ dla $\C P^2\#\overline{\C P^2}$, $S^2\times S^2$ i $\C P^2\#\C P^2$. Które pary można odróżnić bez obliczania $\operatorname{ks}$?
\item \textbf{Dwie różne sumy.} Czy forma $H\oplus[1]$ jest izometryczna z formą $[1]\oplus[1]\oplus[-1]$? Skonstruuj odpowiednią zmianę bazy lub znajdź przeszkodę. Co wynika dla modeli o równym $\operatorname{ks}$?
\item \textbf{Hipoteza Poincarégo a wzmacnianie wniosku.} Załóż, że $X$ jest zamkniętą, jednospójną topologiczną $4$-rozmaitością i $H_2(X;\Z)=0$. Ustal jej dopuszczalne dane $(Q,\operatorname{ks})$ i wskaż dokładny wniosek klasyfikacyjny. Wyjaśnij, dlaczego nie jest to automatycznie twierdzenie o dyfeomorfizmie.
\item \textbf{Przeszkoda stabilna i gładka.} Dla form $E_8$, $2E_8$ oraz $H\oplus E_8$ porównaj wnioski z niezmiennika Kirby'ego--Siebenmanna i z twierdzenia Donaldsona o formach określonych. Rozdziel przypadki wykluczone od przypadków, których te dwa narzędzia nie rozstrzygają.
\wskazowka{Oddziel stabilną wygładzalność od istnienia gładkiej struktury na samej rozmaitości. Użyj twierdzenia Donaldsona tylko przy spełnieniu jego hipotez.}
\item \textbf{Skończona lista kandydatów.} Wypisz wszystkie typy topologiczne gwarantowane przez twierdzenie Freedmana, których forma jest izometryczna z $\operatorname{diag}(1,-1,-1)$, a następnie porównaj ich liczby Bettiego i klasy spinowe.
\end{enumerate}

\chapter{Twierdzenie Donaldsona i teoria połączeń ASD}
\begin{enumerate}
\item \textbf{Rozkład form na składowe własne gwiazdy.} W standardowo zorientowanej $\R^4$ rozłóż $\omega=2\,\dd x^1\wedge\dd x^2+3\,\dd x^3\wedge\dd x^4+\dd x^1\wedge\dd x^3$ na $\omega^+$ i $\omega^-$. Oblicz ich normy punktowe i zweryfikuj ortogonalność.
\item \textbf{Nowe lokalne połączenie ASD.} Niech $T\in\mathfrak{su}(2)$ będzie stałą niezerową macierzą i $A=T(x^1\dd x^3+x^2\dd x^4)$ na $\R^4$. Oblicz $F_A$, sprawdź antysamodualność oraz wyjaśnij, dlaczego energia tego połączenia na całej przestrzeni jest nieskończona.
\wskazowka{Dla macierzowej formy koneksji oblicz $F_A=dA+A\wedge A$. Następnie rozłóż $2$-formy na wartości własne gwiazdy Hodge'a.}
\item \textbf{Cechowanie w podgrupie abelowej.} Dla $A=T(x^1\dd x^2-x^3\dd x^4)$ oraz $g(x)=\exp(T\varphi(x))$, gdzie $\varphi$ jest dowolną funkcją gładką, wyznacz $A^g$ i $F_{A^g}$. Zinterpretuj różnicę między zmianą lokalnego opisu a zmianą krzywizny.
\item \textbf{Energia a rozkład krzywizny.} Na zamkniętej zorientowanej $4$-rozmaitości niech $k=c_2(E)[X]$. Zakładając $\|F_A^+\|_{L^2}^2=3\pi^2$, wyznacz $\|F_A^-\|_{L^2}^2$ i energię całkowitą dla $k=0,1,2$. Które wyniki są zgodne z nieujemnością norm?
\wskazowka{Użyj ortogonalnego rozkładu $F=F^++F^-$ oraz relacji między $\int\operatorname{tr}(F\wedge F)$ i liczbą Cherna; pilnuj konwencji znaku.}
\item \textbf{Przypadek samodualny.} Dla połączenia spełniającego $F_A^-=0$ wykorzystaj tożsamość energii do ustalenia znaku $k$ i wartości energii. Porównaj ze znakiem liczby Cherna dla połączeń ASD przy konwencjach skryptu.
\item \textbf{Wpływ odwrócenia orientacji.} Jak zmieniają się operacja $*$ na $2$-formach, rozkład $F_A=F_A^++F_A^-$ oraz liczba $c_2(E)[X]$ po odwróceniu orientacji $X$? Zapisz odpowiadające równanie dla połączenia, które było ASD przed zmianą orientacji.
\item \textbf{Skalowanie instantonu BPST.} Dla gęstości energii $|F_{A_\lambda}|^2=48\lambda^4/(r^2+\lambda^2)^4$ oblicz, jak zmieniają się maksimum, szerokość charakterystyczna i całka energii przy $\lambda\mapsto c\lambda$. Zinterpretuj koncentrację krzywizny w granicy $\lambda\to0^+$.
\wskazowka{W podstawieniu $x=\lambda y$ w czterowymiarowej całce energii czynnik objętości i potęga gęstości krzywizny mogą się zrównoważyć.}
\item \textbf{Energia w kuli.} Dla tej samej gęstości oblicz wyrażenie $\int_{|x|\le R}|F_{A_\lambda}|^2\,\dd^4x$ jako funkcję $R/\lambda$. Sprawdź granice dla $R/\lambda\to0$ i $R/\lambda\to\infty$.
\wskazowka{Przejdź na współrzędne radialne w $\R^4$, gdzie $d^4x=2\pi^2r^3dr$ po całkowaniu po sferze. Zmienne bezwymiarowe upraszczają całkę.}
\item \textbf{Krzywizna i globalna trywialność.} Wyjaśnij, dlaczego lokalny wzór $A_\lambda$ na $S^4\setminus\{\infty\}$ nie świadczy o istnieniu globalnej trywializacji wiązki, skoro $c_2(E)[S^4]=1$. W jaki sposób funkcja przejścia może kodować niezerową klasę Cherna?
\item \textbf{Cechowanie i geometria moduli.} Dla płaskiego połączenia $A=0$ na trywialnej wiązce nad torusem $T^4$ wypisz połączenia otrzymane przez przekształcenia cechowania $g:T^4\to\mathrm{SU}(2)$. Czy wszystkie są różnymi punktami przestrzeni moduli? Uzasadnij odpowiedź.
\item \textbf{Wymiar oczekiwany.} Korzystając z wzoru $d=8k-3(1+b_2^+(X))$ dla jednospójnej $4$-rozmaitości i wiązki $\mathrm{SU}(2)$, oblicz $d$ dla $k=1,2,3$ na $S^4$, $\C P^2$ i $S^2\times S^2$ z podanymi orientacjami.
\wskazowka{Dla ustalonej orientacji oblicz najpierw $b_2^+$ z formy przecięcia. Podstaw go do wzoru indeksowego dopiero po sprawdzeniu założenia $b_1=0$.}
\item \textbf{Wybór orientacji a indeks.} Dla $X=\C P^2\#2\overline{\C P^2}$ wyznacz $b_2^+(X)$ w standardowej orientacji i po jej odwróceniu. Porównaj wartości oczekiwanego wymiaru moduli dla $k=2$.
\item \textbf{Zakres twierdzenia o diagonalizacji.} Przeanalizuj formy $2E_8$, $H\oplus E_8$ i $H\oplus(-E_8)$ w świetle twierdzenia Donaldsona. Dla których można stwierdzić niemożność bycia formą przecięcia gładkiej zamkniętej jednospójnej $4$-rozmaitości, a dla których twierdzenie nie ma zastosowania?
\wskazowka{Twierdzenie o diagonalizacji wymaga zamkniętej, gładkiej, jednospójnej rozmaitości oraz formy określonej. Dla form nieokreślonych sam ten wynik niczego nie przesądza.}
\item \textbf{Forma określona w innej bazie.} Niech $Q=-\begin{psmallmatrix}2&1\\1&1\end{psmallmatrix}$. Wykaż, że jest ujemnie określona i unimodularna, po czym znajdź bazę całkowitą, w której $Q=-I_2$. Porównaj liczbę par wektorów $\{\pm v\}$ o kwadracie $-1$ w obu bazach.
\wskazowka{Szukaj wektora o kwadracie $-1$, potem jego ortogonalnego dopełnienia w kracie. W dwuwymiarowym przykładzie pomoże całkowitoliczbowa zmiana bazy.}
\item \textbf{Liczenie wektorów o kwadracie $-1$.} Dla $Q=-I_4$ oraz $Q=-2I_4$ oblicz liczbę par $\{\pm v\}$ spełniających $Q(v,v)=-1$. Zbadaj osobno unimodularność i oceń, które przypadki można porównywać z twierdzeniem o diagonalizacji.
\item \textbf{Dwie przeszkody do gładkości.} Porównaj pary danych $(Q,\operatorname{ks})=(E_8,1)$ i $(2E_8,0)$. Wyjaśnij, którą przeszkodę do wygładzenia wykrywa już $\operatorname{ks}$, a kiedy konieczna jest informacja o formie określonej.
\item \textbf{Rokhlin a Donaldson.} Przyjmując twierdzenie Rokhlina, zbadaj formy $E_8$, $2E_8$ i $3E_8$ pod kątem możliwej sygnatury gładkiej zamkniętej rozmaitości spinowej. Porównaj tę próbę z ograniczeniem Donaldsona i wskaż, dla których form argument Rokhlina nie wystarcza.
\wskazowka{Twierdzenie Rokhlina ogranicza sygnaturę gładkich rozmaitości spinowych, a Donaldson formy określone. Najpierw sprawdź spinowość przez parzystość formy.}
\item \textbf{Granica argumentu indeksowego.} Załóż, że formuła indeksu dla moduli ASD daje dodatni wymiar oczekiwany. Wyjaśnij, dlaczego sama dodatniość indeksu nie gwarantuje istnienia niepłaskiego połączenia ASD, gładkości moduli ani zwartości przestrzeni rozwiązań.
\item \textbf{Osobliwości i energia.} Rozważ rodzinę instantonów BPST o skali $\lambda_n\to0$. Oszacuj energię poza dowolną ustaloną kulą $B_\varepsilon(0)$ i opisz jej granicę. Wyjaśnij, dlaczego tego zjawiska nie widać z samego ograniczenia całkowitej energii.
\wskazowka{Po przeskalowaniu promienia przez $\lambda$ większość energii przenosi się do coraz mniejszej kuli. Użyj wyrażenia z zadania 8.}
\item \textbf{Efekt dodania hiperbolicznej formy.} Niech $Q=E_8\oplus H$ będzie formą topologicznej $4$-rozmaitości. Sprawdź, że twierdzenie Donaldsona o diagonalizacji form \emph{określonych} nie rozstrzyga jej wygładzalności. Dlaczego dopisanie jednego składnika $H$ zmienia zakres stosowalności argumentu?
\wskazowka{Dodanie $H$ czyni formę nieokreśloną. Zbadaj, które hipotezy twierdzenia Donaldsona przestają być spełnione.}
\item \textbf{Krzywizna niebędąca antysamodualną.} Dla $F=T(\dd x^1\wedge\dd x^2+2\dd x^3\wedge\dd x^4)$ oblicz $F^+$ i $F^-$ i stosunek ich norm. Zinterpretuj wynik jako punktowy wkład w tożsamość energii.
\end{enumerate}

\chapter{Egzotyczne struktury gładkie na \texorpdfstring{$\R^4$}{R4}}
\begin{enumerate}
\item \textbf{Trzy rodzaje odwzorowań.} Sformułuj różnicę między homeomorfizmem, dyfeomorfizmem a zanurzeniem gładkim dla dwóch rozmaitości homeomorficznych z $\R^4$. Podaj, które z tych relacji zachowują w szczególności możliwość gładkiego zanurzenia zadanego zwartego świadka $K$.
\item \textbf{Nieegzotyczna zmiana współrzędnych.} Rozważ strukturę gładką na topologicznej $\R^4$ przeniesioną przez homeomorfizm $h(x_1,x_2,x_3,x_4)=(x_1^3,x_2,x_3,x_4)$ ze standardowej struktury. Wyjaśnij, dlaczego powstała rozmaitość jest dyfeomorficzna ze standardową $\R^4$, chociaż odwzorowanie identycznościowe między odpowiednimi atlasami nie musi być dyfeomorfizmem.
\item \textbf{Ślad splotu o dwóch składowych.} Niech $L$ będzie splotem dwóch niezaczepionych niewęzłów z obramowaniami zerowymi. Oblicz $H_2(M_L;\Z)$, formę przecięcia i $\partial M_L$. Następnie porównaj ze śladem splotu Hopfa z dwoma obramowaniami zerowymi; ustal formę przecięcia i wyjaśnij, dlaczego ten splot nie może być topologicznie plastrowalny.
\wskazowka{Dla śladu dołączenia $2$-uchwytów forma przecięcia jest macierzą obramowań i liczb wiązania. Gładka plastrowalność wymaga dodatkowej kontroli dysków w $B^4$.}
\item \textbf{Gdy dyski są tylko topologiczne.} Dla splotu topologicznie plastrowalnego, o którym nie wiadomo, czy jest gładko plastrowalny, zapisz dokładnie, jaki typ zanurzenia $M_L\hookrightarrow\R^4$ wynika z twierdzenia o śladzie. Wskaż, dlaczego nie można automatycznie zastąpić słowa „topologiczne” słowem „gładkie”.
\wskazowka{Twierdzenie o śladzie w kategorii TOP nie daje automatycznie gładkiego zanurzenia. Zachowaj kategorię przy przechodzeniu od dysków do otoczeń tubularnych.}
\item \textbf{Zwarty świadek niezanurzalności.} Niech $K\subset R$ będzie zwartą gładką podrozmaitością z brzegiem, której nie da się gładko zanurzyć w $R'$. Wyjaśnij, co wynika o istnieniu zanurzenia $R\hookrightarrow R'$ i dyfeomorfizmu $R\to R'$. Zastosuj obserwację do dwóch różnych świadków $K_1,K_2$.
\wskazowka{Jeśli $R$ zanurza się w $R'$, ograniczenie tego zanurzenia do dowolnej zwartej podrozmaitości $K\subset R$ również jest gładkim zanurzeniem.}
\item \textbf{Od przeszkody Gompfa do standardowej $\R^4$.} Załóż, że zwarta rozmaitość $X$ z twierdzenia Gompfa nie zanurza się z zachowaniem orientacji w żadnej zamkniętej jednospójnej gładkiej $4$-rozmaitości o ujemnie określonej formie przecięcia. Wyjaśnij, jak otrzymać niezanurzalność $X$ w standardowej $\R^4$ przez osadzenie tej ostatniej w odpowiedniej zamkniętej rozmaitości. Zaznacz rolę orientacji.
\wskazowka{Załóż istnienie zanurzenia w standardowej $\R^4$ i skomponuj je z odpowiednim standardowym otwartym osadzeniem w zamkniętej rozmaitości. Sprawdź orientacje obu kroków.}
\item \textbf{Zachowanie orientacji.} Dlaczego przeszkoda do zanurzenia $R\hookrightarrow R'$ sformułowana tylko dla zanurzeń zachowujących orientację nie musi automatycznie wykluczać zanurzeń odwracających orientację? Przepisz wniosek, który jest uprawniony bez dodatkowych założeń.
\item \textbf{Suma przy końcu dwóch standardowych przestrzeni.} W konkretnym modelu brzegowej sumy kul przeanalizuj sumę przy końcu dwóch kopii standardowej $\R^4$. Wyznacz jej typ topologiczny i gładki. Wyjaśnij, dlaczego w rozdziale wybór konkretnego modelu sumy przy końcu jest wystarczający.
\wskazowka{W wybranym modelu sumy przy końcu rozważ usuwanie odpowiednich półprzestrzeni i sklejanie cylindrycznych końców; odróżnij homeomorfizm od konkretnego modelu gładkiego.}
\item \textbf{Suma przy końcu z zwartymi świadkami.} Dwie ogolone rozmaitości $R,R'$ zawierają zwarte gładkie rozmaitości $K,K'$. Korzystając z własności sumy przy końcu, wskaż model, w którym oba świadki występują rozłącznie. Jaką przeszkodę do zanurzenia w rozmaitości docelowej otrzymamy, jeśli choć jeden ze świadków nie może się w niej zanurzyć?
\wskazowka{Umieść obie zwarte podrozmaitości poza obszarem sklejania końców. Każde późniejsze zanurzenie sumy przy końcu ogranicza się do zanurzeń obu świadków.}
\item \textbf{Wnioskowanie indukcyjne.} W konstrukcji Gompfa załóż, że na kroku $n$ wybiera się nowego zwartego świadka $X_n$ niezanurzalnego we wcześniejszych modelach, po czym buduje $R_{n+1}$ zawierającą $X_n$. Zapisz precyzyjnie implikację wykluczającą zanurzenie $R_{n+1}\hookrightarrow R_i$ dla $i\le n$.
\wskazowka{W argumentacji indukcyjnej wykorzystywany jest konkretny świadek nowego etapu. Zapisz sprzeczność między jego obecnością w $R_{n+1}$ a niedopuszczalnym zanurzeniem w $R_i$.}
\item \textbf{Niezanurzalność nie jest symetryczna.} Zbadaj logicznie relację „$R$ zanurza się gładko w $R'$”. Czy sama niezanurzalność $R\hookrightarrow R'$ wyklucza zanurzenie w przeciwnym kierunku? Dlaczego już jedna kierunkowa przeszkoda wystarcza, by wykluczyć dyfeomorfizm?
\item \textbf{Małe i duże egzotyczne $\R^4$.} Porównaj wnioski z twierdzenia De Michelisa--Freedmana i z konstrukcji Gompfa pod względem zanurzania w standardowej $\R^4$. Wyjaśnij, dlaczego „mała” egzotyczna $\R^4$ nie musi być dyfeomorficzna ze standardową mimo istnienia odpowiedniego zanurzenia.
\wskazowka{Małość oznacza istnienie gładkiego zanurzenia w standardowej $\R^4$, a nie istnienie dyfeomorfizmu na całą przestrzeń. Zestaw to z definicją dużej egzotycznej $\R^4$.}
\item \textbf{Rola otwartości.} Porównaj zbiory niezmienników, których użyto do klasyfikacji zamkniętych jednospójnych $4$-rozmaitości, z argumentem przez zwarte przeszkody dla otwartej $\R^4$. Wyjaśnij, dlaczego sama forma przecięcia nie może rozróżniać wszystkich egzotycznych struktur na $\R^4$.
\item \textbf{Homologia przestrzeni euklidesowej.} Dla dowolnej gładkiej rozmaitości $R$ homeomorficznej z $\R^4$ oblicz jej homologie, grupę podstawową, charakterystykę Eulera w sensie homotopijnym i liczbę końców. Czy te niezmienniki mogą odróżnić standardową i egzotyczną strukturę gładką?
\item \textbf{Gładkie dyski i splot.} Niech $L\subset S^3$ jest splotem, którego składowe ograniczają parami rozłączne gładkie dyski w $B^4$. Jakie własności śladu $M_L$ wynikają z faktu, że te dyski posiadają otoczenia normalne? Jakie twierdzenie trzeba dodatkowo wykorzystać, by wyprowadzić zanurzenie śladu w $\R^4$?
\item \textbf{Nierozstrzygnięta alternatywa.} Załóż, że topologicznie plastrowalny splot $L$ ma ślad, który nie zanurza się gładko w standardowej $\R^4$. Co można stwierdzić o gładkiej plastrowalności $L$? Czy wynika z tego, że $L$ nie jest topologicznie plastrowalny? Uzasadnij obie odpowiedzi.
\item \textbf{Suma przy końcu a znikanie przeszkód.} Omów, dlaczego dołączenie kolejnej ogolonej $\R^4$ w konstrukcji sumy przy końcu nie eliminuje automatycznie przeszkody pochodzącej od istniejącej zwartej podrozmaitości. W odpowiedzi wykorzystaj zachowanie jej gładkiego zanurzenia w nowej przestrzeni.
\wskazowka{Przeszkoda pochodząca od zwartego podzbioru przetrwa, o ile jego włączenie do sumy przy końcu pozostaje gładkim zanurzeniem. Nie wymaga to klasyfikacji całej otwartej rozmaitości.}
\item \textbf{Liczność rodziny struktur.} Z twierdzenia De Michelisa--Freedmana wynika istnienie kontinuum różnych typów gładkich egzotycznych $\R^4$. Porównaj to stwierdzenie z wnioskiem o nieskończenie wielu typach z konstrukcji Gompfa. Wyjaśnij, których własności dodatkowych te twierdzenia wymagają i dlaczego ich wniosków nie należy utożsamiać.
\item \textbf{Zestawienie pięciu rozdziałów.} Zbuduj diagram logiczny łączący: formy przecięcia, twierdzenie Freedmana, uchwyty Cassona, twierdzenie Donaldsona oraz konstrukcję egzotycznych $\R^4$. Przy każdej strzałce wskaż, czy chodzi o bezpośredni rachunek, zastosowanie twierdzenia ze skryptu, czy o zewnętrzny rezultat o istnieniu.
\wskazowka{Ułóż diagram oddzielający strzałki TOP i DIFF: Freedman klasyfikuje topologicznie, Donaldson ogranicza gładkie formy określone, a uchwyty Cassona i konstrukcje zanurzeniowe łączą oba światy.}
\end{enumerate}

\end{document}
